% La promotion de la culture de la paix — Allemand 1ère A — Cours signé OnBuch+
% Programme officiel MINESEC. Compiler avec : tectonic ALA26-la-promotion-de-la-culture-de-la-paix.tex
\newcommand{\DOCMATIERE}{Allemand}
\newcommand{\DOCNIVEAU}{1ère A}
\newcommand{\DOCMODULE}{Chapitre 5 : Citoyenneté}
\newcommand{\DOCLECON}{ALA26}
\newcommand{\DOCTITRE}{La promotion de la culture de la paix}
% OnBuch+ — préambule « cours signés » ENRICHI (compilé avec tectonic / XeLaTeX)
% Système visuel « Pop » d'OnBuch+ : fond crème (jamais blanc), bordures encre
% épaisses, ombres « dures » décalées, titres Archivo Black en capitales.
% Version MATHÉMATIQUES : graphes (pgfplots/tikz), boîtes pédagogiques
% (définition, propriété, méthode, attention, le savais-tu, exercice résolu,
% exercice, TP, bilan), page de garde et signature OnBuch+ ancrée en bas.
%
% Macros attendues dans main.tex AVANT \input{preamble} :
%   \DOCMATIERE  \DOCNIVEAU  \DOCMODULE  \DOCLECON (numéro)  \DOCTITRE
\documentclass[11pt]{article}

\usepackage{fontspec}
\usepackage[ngerman,french]{babel} % français = langue principale ; l'allemand via \all{..} / textebox
\usepackage[a4paper,margin=2cm,top=3.4cm,bottom=2.8cm,headheight=1.4cm,footskip=1.3cm]{geometry}
\usepackage{amsmath,amssymb,amsfonts}
\usepackage{mathtools} % \xmapsto, \coloneqq... souvent utilisés par les modèles
\usepackage{cancel} % simplifications barrées (bilans de forces)
% \abs{z} (module, valeur absolue) et \norm{u} (norme), utilisés par les modèles
\providecommand{\abs}{}\let\abs\relax\DeclarePairedDelimiter{\abs}{\lvert}{\rvert}
\providecommand{\norm}{}\let\norm\relax\DeclarePairedDelimiter{\norm}{\lVert}{\rVert}
\usepackage[table]{xcolor} % [table] : fournit aussi \rowcolors (alternance de lignes)
\usepackage{pagecolor}
\usepackage{tikz}
\usetikzlibrary{arrows.meta,positioning,calc,shapes.geometric,shapes.misc,shapes.arrows,decorations.pathmorphing,decorations.pathreplacing,decorations.markings,patterns,shadows,mindmap,trees,fit,backgrounds,matrix,intersections,angles,quotes}
\usepackage{pgfplots}
\usepackage[version=4]{mhchem} % équations nucléaires \ce{^{235}_{92}U}
\usepackage[european,siunitx]{circuitikz} % schémas électriques
\pgfplotsset{compat=1.18}
\usepgfplotslibrary{fillbetween,groupplots} % aires entre courbes (calcul intégral)
\usepackage{enumitem}
\usepackage{fancyhdr}
% [most] charge toutes les bibliothèques tcolorbox (listings, minted, raster,
% documentation...) et gonfle inutilement la mémoire TeX sur un document long
% (~300 boîtes) : on ne charge que ce qui est réellement utilisé.
\usepackage[skins,breakable]{tcolorbox}
\usepackage{graphicx}
\usepackage{colortbl}
\usepackage{booktabs}
\usepackage{tabularx}
\usepackage{diagbox} % case d'en-tête diagonale des tableaux à double entrée
% Colonnes extensibles centrée / gauche / droite (C, L, R), souvent utilisées par les modèles
\newcolumntype{C}{>{\centering\arraybackslash}X}
\newcolumntype{L}{>{\raggedright\arraybackslash}X}
\newcolumntype{R}{>{\raggedleft\arraybackslash}X}
\usepackage{array}
\usepackage{multicol}
\usepackage{multirow}
\usepackage{siunitx}
% parse-numbers=false : accepte aussi \SI{2,5 \times 10^{-3}}{...}
\sisetup{locale=FR,output-decimal-marker={,},group-minimum-digits=4,parse-numbers=false}
\DeclareSIUnit\ohm{\ensuremath{\symup{\Omega}}} % U+2126 absent de la police
\DeclareSIUnit\division{div} % réglages d'oscilloscope (ms/div, V/div)
\DeclareSIUnit\tr{tr} % tours (vitesse de rotation, ex. \SI{1500}{\tr\per\minute})
\DeclareSIUnit\molar{M} % concentration molaire (ex. \SI{3}{\molar}, \SI{1}{\micro\molar})
\DeclareSIUnit\an{an} % année (le modèle confond parfois avec \year, un compteur TeX) — ex. \SI{5}{\kilo\meter\per\an}
\DeclareSIUnit\FCFA{FCFA} % franc CFA (ex. \SI{500}{\FCFA\per\kilo\gram})
\DeclareSIUnit\degreeBrix{\textdegree Brix} % teneur en sucre d'un sirop/jus (réfractométrie)
\DeclareSIUnit\month{mois} % le modèle utilise parfois \month, un compteur TeX, comme unité de durée
\usepackage{qrcode}
% \degree : commande fréquemment utilisée par les modèles pour un angle ou une
% température en dehors de \SI{}{} (non fournie par les paquets chargés).
\providecommand{\degree}{^{\circ}}
% \qed (fin de démonstration), utilisé par les modèles sans amsthm
\providecommand{\qed}{\hfill$\square$}
% \barre : double barre des valeurs interdites dans un tableau de variations (tkz-tab)
\providecommand{\barre}{\ensuremath{\|}}
% environnement proof (amsthm non chargé) utilisé par les modèles
\ifdefined\proof\else\newenvironment{proof}[1][Démonstration]{\par\noindent\textit{#1.}\ }{\qed\par}\fi
\usepackage{lastpage}
\usepackage{hyperref}
\hypersetup{hidelinks}

% ============================================================
% Palette « Pop » OnBuch+ — mêmes hex que la classe P (plus_ui.dart)
% ============================================================
\definecolor{poporange}{HTML}{F59321}
\definecolor{popdark}{HTML}{E07A0C}
\definecolor{popcream}{HTML}{FFF6E8}
\definecolor{popink}{HTML}{1C1714}
\definecolor{poppurple}{HTML}{7A5AE0}
\definecolor{popgreen}{HTML}{2E9E6A}
\definecolor{poppink}{HTML}{E0457B}
\definecolor{popblue}{HTML}{2F7DE1}
\definecolor{popgold}{HTML}{C98A12}
\definecolor{popmuted}{HTML}{8A7F76}
\definecolor{popline}{HTML}{EADFD2}
\definecolor{popgray}{HTML}{8A7F76}
\colorlet{popgrey}{popgray}
% Alias tolérants (le modèle nomme parfois les couleurs différemment)
\colorlet{popwhite}{white}
\colorlet{popblack}{popink}
\colorlet{popred}{poppink}
\colorlet{popyellow}{popgold}
% Teintes claires (fonds de tableaux/encadrés) — le modèle nomme parfois
% les couleurs avec un suffixe L (ex. popblueL) sans les avoir définies.
\colorlet{poporangeL}{poporange!20}
\colorlet{popdarkL}{popdark!20}
\colorlet{poppurpleL}{poppurple!20}
\colorlet{popgreenL}{popgreen!20}
\colorlet{poppinkL}{poppink!20}
\colorlet{popblueL}{popblue!20}
\colorlet{popgoldL}{popgold!20}
\colorlet{popredL}{popred!20}
\colorlet{popgrayL}{popgray!20}
% Teintes claires pour fonds de tableaux / zones
\colorlet{popblueL}{popblue!10!white}
\colorlet{poporangeL}{poporange!14!white}
\colorlet{popgreenL}{popgreen!12!white}
\colorlet{poppurpleL}{poppurple!12!white}
\colorlet{poppinkL}{poppink!12!white}
\colorlet{popgoldL}{popgold!14!white}

\pagecolor{popcream}

% ============================================================
% Typographie
% ============================================================
\defaultfontfeatures{Ligatures=TeX}
\setmainfont{PlusJakartaSans-Medium.ttf}[Path=../../../fonts/,BoldFont=PlusJakartaSans-Bold.ttf]
% Maths en sans-serif (Fira Math) avec chiffres et lettres droites de Plus
% Jakarta, pour que les formules se fondent dans le texte.
\usepackage[math-style=ISO,bold-style=ISO]{unicode-math}
\setmathfont{FiraMath-Regular.otf}[Path=../../../fonts/]
\setmathfont{PlusJakartaSans-Medium.ttf}[Path=../../../fonts/,range={up/num,up/{Latin,latin},\mathup}]
\usepackage{icomma} % virgule décimale sans espace en mode math
% Caractères mathématiques Unicode absents des polices de texte (Plus Jakarta,
% Archivo Black) : rendus par leur équivalent mathématique, en texte comme en
% formule — sinon ils s'affichent en carré vide (ex. « Arithmétique dans ℤ »).
\usepackage{newunicodechar}
\newunicodechar{ℝ}{\ensuremath{\mathbb{R}}}
\newunicodechar{ℤ}{\ensuremath{\mathbb{Z}}}
\newunicodechar{ℂ}{\ensuremath{\mathbb{C}}}
\newunicodechar{ℕ}{\ensuremath{\mathbb{N}}}
\newunicodechar{ℚ}{\ensuremath{\mathbb{Q}}}
\newunicodechar{𝔻}{\ensuremath{\mathbb{D}}}
\newunicodechar{α}{\ensuremath{\alpha}}
\newunicodechar{β}{\ensuremath{\beta}}
\newunicodechar{γ}{\ensuremath{\gamma}}
\newunicodechar{δ}{\ensuremath{\delta}}
\newunicodechar{ε}{\ensuremath{\varepsilon}}
\newunicodechar{θ}{\ensuremath{\theta}}
\newunicodechar{λ}{\ensuremath{\lambda}}
\newunicodechar{μ}{\ensuremath{\mu}}
\newunicodechar{σ}{\ensuremath{\sigma}}
\newunicodechar{τ}{\ensuremath{\tau}}
\newunicodechar{φ}{\ensuremath{\varphi}}
\newunicodechar{ψ}{\ensuremath{\psi}}
\newunicodechar{ω}{\ensuremath{\omega}}
\newunicodechar{Σ}{\ensuremath{\Sigma}}
% majuscules produites par \MakeUppercase dans les titres (α-aminés -> Α)
\newunicodechar{Α}{\ensuremath{\alpha}}
\newunicodechar{Β}{\ensuremath{\beta}}
\newunicodechar{Γ}{\ensuremath{\Gamma}}
\newunicodechar{Δ}{\ensuremath{\Delta}}
\newunicodechar{Ε}{\ensuremath{\varepsilon}}
\newunicodechar{Θ}{\ensuremath{\Theta}}
\newunicodechar{Λ}{\ensuremath{\Lambda}}
\newunicodechar{Μ}{\ensuremath{\mu}}
\newunicodechar{Τ}{\ensuremath{\tau}}
\newunicodechar{Φ}{\ensuremath{\Phi}}
\newunicodechar{Ψ}{\ensuremath{\Psi}}
\newunicodechar{Ω}{\ensuremath{\Omega}}
\newunicodechar{↦}{\ensuremath{\mapsto}}
\newunicodechar{∈}{\ensuremath{\in}}
\newunicodechar{∉}{\ensuremath{\notin}}
\newunicodechar{∧}{\ensuremath{\wedge}}
\newunicodechar{⇔}{\ensuremath{\Leftrightarrow}}
\newunicodechar{⟺}{\ensuremath{\Longleftrightarrow}}
\newunicodechar{⇒}{\ensuremath{\Rightarrow}}
\newunicodechar{⟹}{\ensuremath{\Longrightarrow}}
\newunicodechar{∘}{\ensuremath{\circ}}
\newunicodechar{∪}{\ensuremath{\cup}}
\newunicodechar{∩}{\ensuremath{\cap}}
\newunicodechar{≡}{\ensuremath{\equiv}}
\newunicodechar{⊂}{\ensuremath{\subset}}
\newunicodechar{⊥}{\ensuremath{\perp}}
\newunicodechar{∃}{\ensuremath{\exists}}
\newunicodechar{∀}{\ensuremath{\forall}}
\newunicodechar{∛}{\ensuremath{\sqrt[3]{\,}}}
\newunicodechar{∜}{\ensuremath{\sqrt[4]{\,}}}
\newunicodechar{⃗}{\ensuremath{{}^{\to}}}
\newunicodechar{ⁿ}{\ifmmode^{n}\else\textsuperscript{\ensuremath{n}}\fi}
\newunicodechar{ˣ}{\ifmmode^{x}\else\textsuperscript{\ensuremath{x}}\fi}
\newunicodechar{ᵇ}{\ifmmode^{b}\else\textsuperscript{\ensuremath{b}}\fi}
\newunicodechar{ʳ}{\ifmmode^{r}\else\textsuperscript{\ensuremath{r}}\fi}
\newunicodechar{ᵘ}{\ifmmode^{u}\else\textsuperscript{\ensuremath{u}}\fi}
\newunicodechar{ʸ}{\ifmmode^{y}\else\textsuperscript{\ensuremath{y}}\fi}
\newunicodechar{ᵃ}{\ifmmode^{a}\else\textsuperscript{\ensuremath{a}}\fi}
\newunicodechar{ᵉ}{\ifmmode^{e}\else\textsuperscript{\ensuremath{e}}\fi}
\newunicodechar{ᵏ}{\ifmmode^{k}\else\textsuperscript{\ensuremath{k}}\fi}
\newunicodechar{ᵖ}{\ifmmode^{p}\else\textsuperscript{\ensuremath{p}}\fi}
\newunicodechar{ᴺ}{\ifmmode^{N}\else\textsuperscript{\ensuremath{N}}\fi}
\newunicodechar{ᶜ}{\ifmmode^{c}\else\textsuperscript{\ensuremath{c}}\fi}
\newunicodechar{ʰ}{\ifmmode^{h}\else\textsuperscript{\ensuremath{h}}\fi}
\newunicodechar{ᵠ}{\ifmmode^{q}\else\textsuperscript{\ensuremath{q}}\fi}
\newunicodechar{ₙ}{\ifmmode_{n}\else\textsubscript{\ensuremath{n}}\fi}
\newunicodechar{ₐ}{\ifmmode_{a}\else\textsubscript{\ensuremath{a}}\fi}
\newunicodechar{ᵢ}{\ifmmode_{i}\else\textsubscript{\ensuremath{i}}\fi}
\newunicodechar{ᵣ}{\ifmmode_{r}\else\textsubscript{\ensuremath{r}}\fi}
\newunicodechar{ₖ}{\ifmmode_{k}\else\textsubscript{\ensuremath{k}}\fi}
\newunicodechar{ᵦ}{\ifmmode_{\beta}\else\textsubscript{\ensuremath{\beta}}\fi}
\newunicodechar{ₓ}{\ifmmode_{x}\else\textsubscript{\ensuremath{x}}\fi}
\newfontfamily\popdisplay{ArchivoBlack-Regular.ttf}[Path=../../../fonts/]
\newfontfamily\poplogo{SpaceGrotesk-Bold.ttf}[Path=../../../fonts/]
\newfontfamily\popsemi{PlusJakartaSans-SemiBold.ttf}[Path=../../../fonts/]
\newfontfamily\popxbold{PlusJakartaSans-ExtraBold.ttf}[Path=../../../fonts/]
\newfontfamily\popxboldls{PlusJakartaSans-ExtraBold.ttf}[Path=../../../fonts/,LetterSpace=3.0]

\setlist[enumerate]{leftmargin=1.7em,itemsep=5pt,topsep=4pt}
\setlist[itemize]{leftmargin=1.7em,itemsep=4pt,topsep=4pt,label={\color{poporange}\textbullet}}
\setlength{\parindent}{0pt}
\setlength{\parskip}{5pt}
\renewcommand{\arraystretch}{1.25}

% pgfplots : style Pop par défaut
\pgfplotsset{
  every axis/.append style={axis line style={popink,line width=1pt},tick style={popink},
    grid style={popline,line width=0.5pt},label style={font=\small},tick label style={font=\footnotesize}},
  popcurve/.style={poporange,line width=1.8pt,no markers,smooth},
}
% Même style disponible dans un \draw TikZ simple (hors pgfplots)
\tikzset{popcurve/.style={poporange,line width=1.8pt}}
\tikzset{popgrid/.style={popline,very thin}}
\colorlet{popcurve}{poporange} % le nom du style est parfois utilisé comme couleur

% ============================================================
% Pill / squiggle
% ============================================================
\newcommand{\poppill}[3]{%
  \tikz[baseline=-0.55ex]\node[draw=popink,line width=1.2pt,rounded corners=2.6pt,%
    fill=#2,inner xsep=6.5pt,inner ysep=3pt]{%
    \popxboldls\fontsize{7}{7}\selectfont\color{#3}\MakeUppercase{#1}};%
}
\newcommand{\popsquiggle}[1]{%
  \tikz[baseline=-0.4ex]{
    \draw[#1,line width=2.6pt,line cap=round]
      plot [smooth] coordinates {(0,0.09) (0.34,0.01) (0.68,0.21) (1.02,0.01) (1.36,0.10)};
  }%
}
% Mot-clé surligné (marqueur orange sous le texte)
\newcommand{\cle}[1]{{\popxbold\color{popdark}#1}}

% ============================================================
% En-tête / pied
% ============================================================
\pagestyle{fancy}
\fancyhf{}
\renewcommand{\headrulewidth}{0pt}
\renewcommand{\footrulewidth}{0pt}
\lhead{\raisebox{-0.15ex}{\poplogo\fontsize{13}{13}\selectfont\color{popink}OnBuch\color{poporange}+}}
\rhead{\poppill{\DOCMATIERE\ \textbullet\ \DOCNIVEAU\ \textbullet\ Leçon \DOCLECON}{white}{popink}}
\lfoot{\popxbold\fontsize{7.6}{7.6}\selectfont\color{popmuted}\MakeUppercase{Cours signé OnBuch+}\ \textbullet\ {\popsemi\DOCTITRE}}
\rfoot{\tikz[baseline=-0.6ex]\node[rounded corners=8.5pt,fill=popink,minimum height=17pt,minimum width=17pt,inner xsep=5pt,inner sep=0]{\popxbold\fontsize{8}{8}\selectfont\color{popcream}\thepage\,/\,\pageref*{LastPage}};}

% ============================================================
% Titres
% ============================================================
\newcommand{\chaptitre}[2]{%
  \begin{tcolorbox}[enhanced,colback=poporange,colframe=popink,boxrule=2.6pt,arc=11pt,
      drop shadow=popink,left=15pt,right=15pt,top=12pt,bottom=11pt,before skip=3pt,after skip=18pt]
    \poppill{#2}{popink}{popcream}\par\vspace{9pt}
    {\raggedright\hyphenpenalty=10000\exhyphenpenalty=10000\popdisplay\fontsize{22}{24}\selectfont\color{popcream}\MakeUppercase{#1}\par}
  \end{tcolorbox}
}
% Section numérotée : pastille orange avec le numéro + titre Archivo
\newcounter{popsec}
\newcommand{\coursec}[1]{%
  \stepcounter{popsec}\setcounter{popsub}{0}%
  \par\vspace{16pt}\noindent
  \begin{minipage}{\linewidth}
  \tikz[baseline=-0.7ex]\node[draw=popink,line width=1.6pt,fill=poporange,rounded corners=4pt,minimum size=22pt,inner sep=0]{\popdisplay\fontsize{12}{12}\selectfont\color{popcream}\thepopsec};%
  \hspace{8pt}{\popdisplay\fontsize{15}{17}\selectfont\color{popink}\MakeUppercase{#1}}\par
  \vspace{4pt}\noindent{\color{poporange}\rule{36pt}{3.6pt}}
  \end{minipage}\par\nopagebreak\vspace{8pt}
}
\newcounter{popsub}
\newcommand{\courssub}[1]{%
  \stepcounter{popsub}%
  \par\vspace{9pt}\noindent{\popxbold\fontsize{11.5}{13}\selectfont\color{popink}{\color{poporange}\thepopsec.\thepopsub}\hspace{6pt}#1}\par\nopagebreak\vspace{4pt}
}

% ============================================================
% Boîtes pédagogiques — même recette : fond blanc, bordure épaisse colorée,
% ombre dure encre, badge plein. Titre optionnel [#1].
% ============================================================
\tcbset{popbase/.style={breakable,enhanced,colback=white,boxrule=2.3pt,arc=9pt,
  shadow={3pt}{-3pt}{0pt}{popink},left=14pt,right=14pt,top=11pt,bottom=11pt,before skip=11pt,after skip=15pt,
  fontupper=\popsemi\fontsize{10.6}{14.5}\selectfont\color{popink},
  fontlower=\popsemi\fontsize{10.6}{14.5}\selectfont\color{popink}}}
% Coche dessinée (le glyphe ✓ n'existe pas dans les polices du document)
\newcommand{\popcheck}[1]{\tikz[baseline=-0.5ex]\draw[#1,line width=1.6pt,line cap=round,line join=round](-0.13,0.0)--(-0.04,-0.09)--(0.13,0.1);}
\providecommand{\coche}{\popcheck{popgreen}}
\providecommand{\resultat}[1]{\par\smallskip\noindent\fcolorbox{popgreen}{popgreenL}{\parbox{\dimexpr\linewidth-2\fboxsep-2\fboxrule\relax}{\textbf{Résultat.} #1}}\par\smallskip}
\newcommand{\popboxhead}[3]{\poppill{#1}{#2}{white}\ifx\relax#3\relax\else\hspace{6pt}{\popxbold\fontsize{10.6}{12}\selectfont\color{popink}#3}\fi\par\vspace{7pt}}

\newtcolorbox{aretenir}[1][]{popbase,colframe=popgreen,before upper={\popboxhead{À retenir}{popgreen}{#1}}}
% Texte en allemand (document de lecture, dialogue, extrait) : cadre
% d'étude. Un titre optionnel (source, auteur) s'affiche dans l'en-tête.
\newtcolorbox{textebox}[1][]{popbase,colframe=popblue,before upper={\popboxhead{Text}{popblue}{#1}\begin{otherlanguage}{ngerman}},after upper={\end{otherlanguage}}}
% Marques ✓ / ✗ (forme correcte / fautive) souvent utilisées par les modèles
\providecommand{\cross}{\textcolor{poppink}{\ensuremath{\times}}}
\providecommand{\xmark}{\cross}\providecommand{\crossmark}{\cross}\providecommand{\cmark}{\ensuremath{\checkmark}}
% Ligne de réponse / justification à compléter (inventée par les modèles)
\providecommand{\justification}[1]{\par\noindent\textit{Justification~:} \dotfill\par}
% \textipa (package tipa non chargé) : la police ne couvre pas l'API -> simple passage
\providecommand{\textipa}[1]{#1}
% Soulignement (ulem non chargé) : \uline, \ul
\providecommand{\uline}[1]{\underline{#1}}\providecommand{\ul}[1]{\underline{#1}}
% \glossaire{\item ...} inventée par les modèles : liste de glossaire
\providecommand{\glossaire}[1]{\begin{itemize}#1\end{itemize}}
% \alert (beamer) inventée par les modèles : texte mis en évidence
\providecommand{\alert}[1]{\textbf{\textcolor{poppink}{#1}}}
% variantes ASCII d'ordinaux écrites en macro
\providecommand{\iere}{\textsuperscript{re}}\providecommand{\ieme}{\textsuperscript{e}}\providecommand{\ere}{\textsuperscript{re}}\providecommand{\eme}{\textsuperscript{e}}\providecommand{\er}{\textsuperscript{er}}\providecommand{\ie}{\textsuperscript{e}}
% Ordinaux français écrits en macro par les modèles : 1\ière, 3\ième
\newcommand{\ière}{\textsuperscript{re}}\newcommand{\ième}{\textsuperscript{e}}
% \exemple (inventée par les modèles) : étiquette « Exemple : »
\providecommand{\exemple}{\textbf{Exemple~:}\ }
% \square utilisable aussi hors mode math (cases à cocher)
\AtBeginDocument{\let\squareMath\square\renewcommand{\square}{\ensuremath{\squareMath}}}
% Mot ou phrase en allemand à l'intérieur d'un paragraphe français
\newcommand{\all}[1]{\ifmmode\text{\textit{#1}}\else\begingroup\selectlanguage{ngerman}\itshape #1\endgroup\fi}
\let\esp\all % alias toléré
\newtcolorbox{exemplebox}[1][]{popbase,colframe=poporange,before upper={\popboxhead{Exemple}{poporange}{#1}}}
\newtcolorbox{definition}[1][]{popbase,colframe=popblue,before upper={\popboxhead{Définition}{popblue}{#1}}}
\newtcolorbox{propriete}[1][]{popbase,colframe=poppurple,before upper={\popboxhead{Propriété}{poppurple}{#1}}}
\newtcolorbox{methode}[1][]{popbase,colframe=popgold,before upper={\popboxhead{Méthode}{popgold}{#1}}}
\newtcolorbox{attention}[1][]{popbase,colframe=poppink,before upper={\popboxhead{Attention}{poppink}{#1}}}
\newtcolorbox{experience}[1][]{popbase,colframe=popblue,colback=popblueL,before upper={\popboxhead{Expérience}{popblue}{#1}}}
% « Le savais-tu ? » : carte violette pleine (contexte, culture, Cameroun)
\newtcolorbox{savaistu}[1][]{popbase,colback=poppurple,colframe=popink,
  fontupper=\popsemi\fontsize{10.4}{14.2}\selectfont\color{white},
  before upper={\poppill{Le savais-tu\,?}{white}{poppurple}\ifx\relax#1\relax\else\hspace{6pt}{\popxbold\color{white}#1}\fi\par\vspace{7pt}}}
% Exercice résolu : énoncé en haut, \tcblower puis la solution sur fond crème
\newcounter{popexo}\newcounter{popexr}
\newtcolorbox{exoresolu}[1][]{popbase,colframe=poporange,colbacklower=popcream,
  segmentation style={popink,line width=1.2pt,dashed},
  before upper={\stepcounter{popexr}\popboxhead{Exercice résolu \thepopexr}{poporange}{#1}},
  before lower={\poppill{Solution}{popink}{popcream}\par\vspace{6pt}}}
% Exercice (non résolu en place) — la correction est dans la partie Corrigés
\newtcolorbox{exercice}[1][]{popbase,colframe=popink,
  before upper={\stepcounter{popexo}\popboxhead{Exercice \thepopexo}{popink}{#1}}}
% Corrigé
\newtcolorbox{corrige}[1][]{popbase,colframe=popgreen,colback=popgreenL,
  before upper={\popboxhead{Corrigé}{popgreen}{#1}}}
% Objectifs / prérequis / bilan
\newtcolorbox{objectifs}{popbase,colframe=popink,colback=poporangeL,
  before upper={\popboxhead{Objectifs de la leçon}{popink}{}}}
\newtcolorbox{prerequis}{popbase,colframe=popink,colback=popblueL,
  before upper={\popboxhead{Prérequis}{popblue}{}}}
\newtcolorbox{bilan}{popbase,colframe=popink,colback=white,boxrule=2.8pt,
  before upper={\popboxhead{Fiche bilan}{poporange}{L'essentiel à connaître pour l'examen}}}
% Figure encadrée avec légende
\newtcolorbox{popfigure}[1][]{enhanced,breakable=false,colback=white,colframe=popline,boxrule=1.4pt,arc=8pt,
  left=10pt,right=10pt,top=10pt,bottom=8pt,before skip=10pt,after skip=12pt,halign=center,
  lower separated=false,halign lower=center,
  fontlower=\popsemi\fontsize{9}{11.5}\selectfont\color{popmuted},
  #1}
\newcommand{\legende}[1]{\par\vspace{6pt}{\popsemi\fontsize{9}{11.5}\selectfont\color{popmuted}#1}}

% ============================================================
% Signature OnBuch+
% ============================================================
\newcommand{\poplogoinline}[1]{{\poplogo\fontsize{#1}{#1}\selectfont\color{popink}OnBuch\color{poporange}+}}
\newcommand{\signatureonbuch}{%
  \begin{tcolorbox}[enhanced,colback=popink,colframe=popink,boxrule=2.2pt,arc=10pt,
      drop shadow=poporange,left=14pt,right=14pt,top=10pt,bottom=10pt,before skip=0pt,after skip=0pt]
    \tikz[baseline=-0.7ex]\node[circle,fill=popgreen,draw=popcream,line width=1.3pt,minimum size=22pt,inner sep=0]{\popcheck{white}};%
    \hspace{9pt}\begin{minipage}[c]{0.62\linewidth}
      {\popxbold\fontsize{12}{13}\selectfont\color{popcream}Signé\ \textbullet\ {\poplogo OnBuch\color{poporange}+}}\par\vspace{2pt}
      {\raggedright\popsemi\fontsize{8}{10.5}\selectfont\color{popline}\MakeUppercase{Équipe pédagogique OnBuch+}\\\MakeUppercase{Conforme au programme officiel MINESEC}\par}
    \end{minipage}\hfill
    {\poplogo\fontsize{22}{22}\selectfont\color{popcream}OnBuch\color{poporange}+}
  \end{tcolorbox}%
}

% ============================================================
% Page de garde
% #1 titre · #2 matière · #3 niveau · #4 module · #5 n° leçon · #6 durée ·
% #7 description / objectifs (texte ou itemize)
% ============================================================
\newcommand{\pagegarde}[7]{%
  \thispagestyle{empty}
  \newgeometry{a4paper,margin=1.7cm,top=1.6cm,bottom=1.4cm}%
  \begin{tikzpicture}[remember picture,overlay]
    \fill[poporange!18] ([xshift=-3.2cm,yshift=-3.6cm]current page.north east) circle (4.6cm);
    \fill[poppurple!14] ([xshift=2.4cm,yshift=7.6cm]current page.south west) circle (3.8cm);
  \end{tikzpicture}%
  {\poplogo\fontsize{30}{30}\selectfont\color{popink}OnBuch\color{poporange}+}\hfill
  \raisebox{0.6ex}{\poppill{Cours signé \textbullet\ Premium}{popink}{popcream}}\par
  \vspace{1pt}\popsquiggle{poppurple}\par
  \vspace{8pt}
  \begin{tcolorbox}[enhanced,colback=poporange,colframe=popink,boxrule=2.8pt,arc=13pt,
      drop shadow=popink,left=18pt,right=18pt,top=12pt,bottom=13pt,after skip=10pt]
    \poppill{#2 \textbullet\ #3}{popink}{popcream}\hspace{5pt}\poppill{#4}{white}{popink}\par\vspace{8pt}
    {\popdisplay\fontsize{14}{14}\selectfont\color{popink}LEÇON #5}\par\vspace{4pt}
    {\raggedright\hyphenpenalty=10000\exhyphenpenalty=10000\popdisplay\fontsize{25}{27}\selectfont\color{popcream}\MakeUppercase{#1}\par}
    \vspace{8pt}
    {\popxbold\fontsize{9.5}{9.5}\selectfont\color{popink}\MakeUppercase{Durée conseillée : #6}}
  \end{tcolorbox}
  \begin{tcolorbox}[enhanced,colback=white,colframe=popink,boxrule=2.2pt,arc=10pt,
      drop shadow=popink,left=16pt,right=16pt,top=10pt,bottom=10pt,after skip=8pt]
    \poppill{Dans cette leçon}{poppurple}{white}\par\vspace{5pt}
    \popsemi\fontsize{9.3}{12.6}\selectfont\color{popink}#7
  \end{tcolorbox}
  \noindent\begin{minipage}[t]{0.487\linewidth}
    \begin{tcolorbox}[enhanced,colback=popgreen,colframe=popink,boxrule=2.2pt,arc=10pt,
        drop shadow=popink,left=11pt,right=11pt,top=10pt,bottom=10pt,height=3.1cm,valign=center]
      \tcbox[colback=white,colframe=popink,boxrule=1.3pt,arc=5pt,boxsep=0pt,nobeforeafter,
        left=3pt,right=3pt,top=3pt,bottom=3pt,valign=center]{\qrcode[height=1.9cm]{https://play.google.com/store/apps/details?id=com.luvvix.onbuch&hl=fr}}%
      \hspace{8pt}\begin{minipage}[c]{0.5\linewidth}
        {\popdisplay\fontsize{11}{12.5}\selectfont\color{white}\MakeUppercase{Télécharge OnBuch}}\par\vspace{4pt}
        {\raggedright\popsemi\fontsize{8.4}{11}\selectfont\color{white}Gratuit \textbullet\ Google Play\\Tuteur IA, annales, concours, résultats\par}
      \end{minipage}
    \end{tcolorbox}
  \end{minipage}\hfill
  \begin{minipage}[t]{0.487\linewidth}
    \begin{tcolorbox}[enhanced,colback=poppurple,colframe=popink,boxrule=2.2pt,arc=10pt,
        drop shadow=popink,left=11pt,right=11pt,top=10pt,bottom=10pt,height=3.1cm,valign=center]
      \tikz[baseline=-0.7ex]\node[circle,fill=white,draw=popink,line width=1.3pt,minimum size=1.6cm,inner sep=0]{\popdisplay\fontsize{24}{24}\selectfont\color{poppurple}+};%
      \hspace{8pt}\begin{minipage}[c]{0.6\linewidth}
        {\popdisplay\fontsize{11}{12.5}\selectfont\color{white}\MakeUppercase{Passe à OnBuch+}}\par\vspace{4pt}
        {\raggedright\popsemi\fontsize{8.4}{11}\selectfont\color{white}Tous les cours signés, Tuteur IA illimité, fiches PDF — onglet OnBuch+ de l'app\par}
      \end{minipage}
    \end{tcolorbox}
  \end{minipage}
  \vspace{8pt}
  \popcontenu
  \vfill
  \signatureonbuch
  \restoregeometry
  \newpage
}

% Rangée « Ce cours contient » (page de garde)
\newcommand{\poptile}[3]{% #1 couleur #2 gros texte #3 légende
  \begin{tcolorbox}[enhanced,colback=#1,colframe=popink,boxrule=2pt,arc=9pt,shadow={2.5pt}{-2.5pt}{0pt}{popink},
      left=6pt,right=6pt,top=9pt,bottom=9pt,halign=center,valign=center,height=1.75cm,width=\linewidth,nobeforeafter]
    {\popdisplay\fontsize{12.5}{13}\selectfont\color{white}\MakeUppercase{#2}}\par\vspace{3pt}
    {\centering\popsemi\fontsize{7.8}{9.5}\selectfont\color{white}#3\par}
  \end{tcolorbox}}
\newcommand{\popcontenu}{%
  \par\noindent{\popxboldls\fontsize{7.5}{7.5}\selectfont\color{popmuted}CE COURS SIGNÉ CONTIENT}\par\vspace{7pt}
  \noindent\begin{minipage}[t]{0.235\linewidth}\poptile{popblue}{Cours}{complet, illustré, pas à pas}\end{minipage}\hfill
  \begin{minipage}[t]{0.235\linewidth}\poptile{poporange}{Méthodes}{et exercices résolus}\end{minipage}\hfill
  \begin{minipage}[t]{0.235\linewidth}\poptile{poppink}{Exercices}{type examen + corrigés}\end{minipage}\hfill
  \begin{minipage}[t]{0.235\linewidth}\poptile{popgreen}{Fiche bilan}{l'essentiel à retenir}\end{minipage}\par}

% Sommaire
\newtcolorbox{sommaire}{enhanced,colback=white,colframe=popink,boxrule=2.2pt,arc=10pt,
  drop shadow=popink,left=18pt,right=18pt,top=13pt,bottom=13pt,before skip=6pt,after skip=20pt,
  before upper={\poppill{Sommaire}{poppurple}{white}\par\vspace{9pt}\popsemi\fontsize{10.6}{14.5}\selectfont\color{popink}}}

% Signature de fin de cours — poussée tout en bas de la dernière page
\newcommand{\findecours}{%
  \par\vfill\nopagebreak
  \begin{center}\popsquiggle{poporange}\end{center}\vspace{4pt}
  \signatureonbuch
}
\AtBeginDocument{\def\llbracket{\mathopen{[\mkern-3.5mu[}}\def\rrbracket{\mathclose{]\mkern-3.5mu]}}} % intervalles d'entiers (glyphe ⟦ absent de la police)
% Glyphes absents de Fira Math : redéfinis à partir de glyphes disponibles
\AtBeginDocument{%
  \def\setminus{\mathbin{\backslash}}\let\smallsetminus\setminus
  \def\vdots{\mathord{\vbox{\baselineskip3.2pt\lineskiplimit0pt\kern4pt\hbox{.}\hbox{.}\hbox{.}}}}%
  \def\ddots{\mathinner{\mkern1mu\raise6.5pt\hbox{.}\mkern2mu\raise3.6pt\hbox{.}\mkern2mu\raise0.7pt\hbox{.}\mkern1mu}}%
  \def\triangle{\mathord{\scalebox{0.9}{$\symup{\Delta}$}}}%
  \def\nabla{\mathord{\raisebox{\depth}{\scalebox{1}[-1]{$\symup{\Delta}$}}}}%
  \def\checkmark{\popcheck{popdark}}%
  \def\star{\mathbin{\ast}}\def\bigstar{\mathord{\ast}}%
  \def\diamond{\mathbin{\scalebox{0.6}{$\Diamond$}}}%
  \def\bigcup{\mathop{\mathchoice{\raisebox{-0.3em}{\scalebox{1.7}{$\cup$}}}{\scalebox{1.25}{$\cup$}}{\cup}{\cup}}\displaylimits}%
  \def\bigcap{\mathop{\mathchoice{\raisebox{-0.3em}{\scalebox{1.7}{$\cap$}}}{\scalebox{1.25}{$\cap$}}{\cap}{\cap}}\displaylimits}%
  \def\lesseqqgtr{\mathrel{\lessgtr}}\def\bigoplus{\mathop{\oplus}\displaylimits}%
}
\providecommand{\ssi}{si et seulement si} % abréviation fréquente
\AtBeginDocument{\providecommand{\wideparen}{\overparen}} % arc de cercle

\begin{document}

\pagegarde{La promotion de la culture de la paix}{Allemand}{1ère A}{Chapitre 5 : Citoyenneté}{ALA26}{2 h}{Cette leçon propose la lecture et le commentaire d'un texte allemand sur la promotion de la culture de la paix. Elle permet d'acquérir le vocabulaire thématique (Frieden, Krieg, Versöhnung, Dialog, Toleranz, Verständigung) ainsi que deux outils grammaticaux essentiels : l'impératif pour conseiller et le futur pour envisager l'avenir.
\vspace{4pt}
\begin{multicols}{2}\small
\begin{itemize}
\item À la fin de cette leçon, je sais comprendre globalement puis en détail un texte allemand sur la promotion de la culture de la paix.
\item Je sais employer correctement le vocabulaire du thème (der Frieden, der Krieg, die Versöhnung, der Dialog, die Toleranz, die Verständigung) avec articles et pluriels.
\item Je sais former et utiliser l'impératif allemand (du, ihr, Sie) pour donner un conseil ou une consigne.
\item Je sais conjuguer et employer le futur I (werden + infinitif) pour parler de l'avenir et formuler des prévisions.
\item Je sais distinguer futur I et présent + complément de temps pour exprimer le futur.
\item Je sais répondre à des questions de compréhension en phrases complètes.
\end{itemize}\end{multicols}}

\begin{sommaire}
\begin{multicols}{2}
\begin{enumerate}
\item Texte support : Frieden beginnt bei uns
\item Compréhension du texte
\item Lexique thématique : Frieden und Konflikt
\item Grammaire 1 : l'impératif et le conseil
\item Grammaire 2 : le futur
\item Méthode du commentaire et production guidée
\item Activité d'intégration
\item Méthodes et exercices résolus
\item Exercices
\item Corrigés des exercices
\item Fiche bilan
\end{enumerate}
\end{multicols}
\end{sommaire}

\begin{objectifs}
\begin{itemize}
\item À la fin de cette leçon, je sais comprendre globalement puis en détail un texte allemand sur la promotion de la culture de la paix.
\item Je sais employer correctement le vocabulaire du thème (der Frieden, der Krieg, die Versöhnung, der Dialog, die Toleranz, die Verständigung) avec articles et pluriels.
\item Je sais former et utiliser l'impératif allemand (du, ihr, Sie) pour donner un conseil ou une consigne.
\item Je sais conjuguer et employer le futur I (werden + infinitif) pour parler de l'avenir et formuler des prévisions.
\item Je sais distinguer futur I et présent + complément de temps pour exprimer le futur.
\item Je sais répondre à des questions de compréhension en phrases complètes.
\item Je sais résumer un texte en quelques phrases avec mes propres mots.
\item Je sais rédiger un court paragraphe guidé pour promouvoir la paix dans mon école ou ma communauté.
\end{itemize}
\end{objectifs}

\begin{prerequis}
\begin{itemize}
\item La conjugaison des verbes réguliers et irréguliers au présent (Seconde)
\item Les articles définis et indéfinis et le genre des noms (Seconde)
\item Les verbes modaux, notamment sollen et können, utiles pour le conseil (Seconde)
\item La place du verbe dans la phrase déclarative et l'inversion sujet-verbe (Seconde)
\item Le vocabulaire de base de la vie scolaire et de la société (début de 1ère)
\end{itemize}
\end{prerequis}

\newpage

\coursec{Texte support : Frieden beginnt bei uns}

\courssub{Introduction : la paix, une affaire de tous les jours}

Après les crises post-électorales ou les conflits dans certaines régions du Cameroun, on entend souvent parler de « culture de la paix » à la radio et dans les journaux. À Yaoundé comme à Douala, les clubs de paix des lycées organisent des débats, des marches et des journées de réflexion. En Allemagne, pays marqué par deux guerres mondiales, la paix (\all{der Frieden}) est au cœur de l'éducation civique : chaque année, le \all{Friedenspreis des Deutschen Buchhandels} (le prix de la paix des libraires allemands) récompense une personnalité engagée pour la réconciliation entre les peuples.

Mais la paix ne se construit pas seulement dans les traités internationaux : elle commence dans la cour de récréation, dans la famille, au marché, entre voisins. Comment parler de paix, de dialogue et de tolérance en allemand ? Comment donner des conseils pour résoudre un conflit et imaginer l'avenir ? C'est ce que nous allons découvrir à travers un texte sur la promotion de la culture de la paix.

\textbf{Problématique :} comment un simple projet scolaire peut-il transformer les conflits du quotidien en occasions de dialogue et de réconciliation ?

\begin{definition}[La culture de la paix]
\all{Die Kultur des Friedens} (la culture de la paix) désigne l'ensemble des valeurs — \all{die Toleranz} (la tolérance), \all{der Dialog} (le dialogue), \all{der Respekt} (le respect), \all{die Versöhnung} (la réconciliation) — et des comportements qui permettent de \textbf{prévenir} et de \textbf{résoudre} les conflits \textbf{sans violence} (\all{gewaltfrei}). Elle s'apprend et se pratique : à l'école, en famille, dans la communauté et en politique.
\end{definition}

\courssub{Lecture globale : découvrir le texte}

\begin{methode}[Comment lire un texte allemand ?]
\begin{enumerate}
\item \textbf{Lire le titre} et deviner le thème : \all{Frieden beginnt bei uns} — « La paix commence avec nous ». De quoi va parler le texte ? Probablement de la paix et de notre responsabilité personnelle.
\item \textbf{Lire une première fois sans dictionnaire}, rapidement, pour saisir le sens général : qui ? où ? quoi ?
\item \textbf{Repérer les mots transparents} : \all{Konflikt, Toleranz, Dialog, organisieren, diskutieren, Politik, Problem} se devinent facilement pour un francophone.
\item \textbf{Relire paragraphe par paragraphe} avec le glossaire, en soulignant les mots inconnus.
\item \textbf{Résumer} chaque paragraphe en une phrase, en français ou en allemand simple.
\end{enumerate}
\end{methode}

\begin{attention}[Ne traduis pas mot à mot !]
L'erreur classique du francophone est de traduire chaque mot dans l'ordre avant de comprendre la phrase. En allemand, le verbe conjugué est à la \textbf{deuxième place} et l'infinitif ou le participe va souvent \textbf{à la fin}. Cherche d'abord le \textbf{sens global} : qui fait quoi ? Consulte le glossaire seulement après avoir formulé une hypothèse de sens. Exemple : \all{Sie helfen, Probleme gewaltfrei zu lösen.} — ne t'arrête pas sur chaque mot : le sens global est « ils aident à résoudre les problèmes sans violence ».
\end{attention}

\courssub{Le texte : Frieden beginnt bei uns}

Voici un article paru dans le journal d'un lycée de Stuttgart, en Allemagne. Lis-le d'abord en silence, puis à voix haute.

\begin{textebox}[Frieden beginnt bei uns — Artikel aus einer Schülerzeitung in Stuttgart]
In unserer Schule in Stuttgart gibt es viele Konflikte im Alltag: Streit in der Pause, Beleidigungen in der Klasse, manchmal auch Gewalt auf dem Schulhof. Viele Schüler sind wütend, aber sie wissen nicht, wie sie ihre Probleme lösen können.

Aber seit einem Jahr arbeiten zwölf Schüler aus verschiedenen Klassen als Streitschlichter. Sie tragen ein blaues T-Shirt mit dem Wort „Frieden“. Ihre Arbeit ist einfach, aber wichtig: Sie hören zu, sie vermitteln, sie helfen, Probleme gewaltfrei zu lösen. Wenn zwei Schüler einen Konflikt haben, kommen sie zu den Streitschlichtern. Jeder erzählt seine Version, und gemeinsam suchen sie eine Lösung. „Redet miteinander statt übereinander!“, sagt ihre Lehrerin Frau Berger. „Und vergesst nie: Zuhören ist der erste Schritt zum Frieden.“

Die Ergebnisse sind positiv: Die Schule ist ruhiger, die Schüler respektieren einander besser, und es gibt weniger Gewalt. Auch die Lehrer sind stolz auf das Projekt.

Nächstes Jahr werden wir ein Friedensfest organisieren. Wir werden alle Klassen einladen, wir werden Musik machen, und wir werden über Toleranz und Dialog diskutieren. Vielleicht werden auch andere Schulen zu uns kommen.

Denn Frieden beginnt nicht in der Politik, sondern bei jedem von uns — in der Schule, in der Familie und in der Gemeinde.
\end{textebox}

\textbf{Glossaire :}

\begin{center}
\begin{tabularx}{\linewidth}{|X|X|}\hline
\rowcolor{popblueL} \textbf{Deutsch} & \textbf{Français} \\ \hline
die Streitschlichter (Pl.) & les médiateurs (élèves qui règlent les conflits) \\ \hline
der Alltag, -e & le quotidien, la vie de tous les jours \\ \hline
lösen (hat gelöst) & résoudre \\ \hline
gewaltfrei & non violent, sans violence \\ \hline
der Streit, -e & la dispute, la querelle \\ \hline
die Beleidigung, -en & l'insulte \\ \hline
die Gewalt (Sg.) & la violence \\ \hline
vermitteln (hat vermittelt) & servir de médiateur, arbitrer \\ \hline
die Lösung, -en & la solution \\ \hline
das Friedensfest, -e & la fête de la paix \\ \hline
\end{tabularx}
\end{center}

\begin{attention}[Articles et pluriels : apprends toujours le nom complet]
En allemand, un nom s'apprend toujours avec son \textbf{article} et son \textbf{pluriel} : \all{der Konflikt, -e} ; \all{die Lösung, -en} ; \all{das Fest, -e}. Sans l'article, impossible de décliner correctement ! Et n'oublie pas la \textbf{majuscule} à tous les noms : on écrit \all{die Schule}, jamais \all{die schule}.
\end{attention}

\courssub{Repérage des informations globales : qui ? où ? quoi ?}

Après une première lecture, réponds mentalement à trois questions simples :

\begin{itemize}
\item \textbf{Qui ?} (\all{Wer?}) — Douze élèves de différentes classes, formés comme médiateurs (\all{Streitschlichter}), et leur professeure, Frau Berger.
\item \textbf{Où ?} (\all{Wo?}) — Dans un lycée de Stuttgart, en Allemagne.
\item \textbf{Quoi ?} (\all{Was?}) — Un projet scolaire de médiation (\all{ein Schülerprojekt}) pour résoudre les conflits sans violence, et un projet futur : une fête de la paix (\all{ein Friedensfest}).
\end{itemize}

\begin{exoresolu}[Compréhension globale]
\textbf{Énoncé :} Réponds en allemand avec une phrase complète : \all{Was machen die Streitschlichter in der Schule?}
\tcblower
\textbf{Solution détaillée :} On cherche dans le deuxième paragraphe les verbes qui décrivent leur action : \all{sie hören zu, sie vermitteln, sie helfen}. Réponse complète : \all{Die Streitschlichter hören zu, sie vermitteln und sie helfen, Probleme gewaltfrei zu lösen.} « Les médiateurs écoutent, servent d'arbitres et aident à résoudre les problèmes sans violence. » Attention à la place du verbe : \all{Die Streitschlichter} (sujet) + \all{hören} (verbe en 2e position).
\end{exoresolu}

\courssub{Lecture détaillée : les idées principales de chaque paragraphe}

Relisons le texte paragraphe par paragraphe et résumons l'idée essentielle de chacun :

\begin{enumerate}
\item \textbf{Paragraphe 1 — le problème :} il y a beaucoup de conflits dans le quotidien scolaire (disputes, insultes, parfois violence), et les élèves ne savent pas comment les résoudre.
\item \textbf{Paragraphe 2 — la solution :} douze élèves travaillent comme médiateurs ; ils écoutent, arbitrent et aident à résoudre les problèmes sans violence. Le conseil de Frau Berger : \all{Redet miteinander statt übereinander!} « Parlez les uns avec les autres au lieu de parler les uns des autres ! »
\item \textbf{Paragraphe 3 — les résultats :} l'école est plus calme, les élèves se respectent davantage, il y a moins de violence.
\item \textbf{Paragraphe 4 — l'avenir :} l'année prochaine, les élèves organiseront une fête de la paix avec de la musique et des débats sur la tolérance et le dialogue.
\item \textbf{Paragraphe 5 — la conclusion :} la paix ne commence pas en politique, mais chez chacun de nous : à l'école, en famille, dans la communauté.
\end{enumerate}

\begin{popfigure}
\begin{center}
\begin{tikzpicture}[
  box/.style={draw=popblue, thick, rounded corners, fill=popblueL, align=center, minimum width=3.4cm, minimum height=1cm, font=\small},
  arr/.style={-{Stealth[length=3mm]}, thick, popdark}]
\node[box, align=center] (intro) at (0,0){\textbf{Introduction}\\le problème :\\les conflits};
\node[box, fill=poporangeL, draw=poporange, align=center] (dev) at (0,-2){\textbf{Développement}\\la solution :\\la médiation};
\node[box, fill=popgreenL, draw=popgreen, align=center] (conc) at (0,-4){\textbf{Conclusion}\\l'avenir :\\l'espoir};
\draw[arr] (intro) -- (dev);
\draw[arr] (dev) -- (conc);
\end{tikzpicture}
\end{center}
\legende{Structure du texte : du problème (les conflits) à la solution (la médiation), puis vers l'avenir (l'espoir).}
\end{popfigure}

\courssub{Relevé dans le texte : le champ lexical de la paix}

Soulignons dans le texte tous les mots qui appartiennent au champ lexical de la paix et du conflit, et relevons aussi les formes verbales remarquables (impératif et futur), que nous étudierons en grammaire.

\begin{center}
\begin{tabularx}{\linewidth}{|X|X|X|}\hline
\rowcolor{popblueL} \textbf{Catégorie} & \textbf{Mots relevés dans le texte} & \textbf{Français} \\ \hline
La paix (Noms/Adj.) & \all{der Frieden, das Friedensfest, gewaltfrei, die Toleranz, der Dialog} & la paix, la fête de la paix, non violent, la tolérance, le dialogue \\ \hline
Le conflit (Noms/Adj.) & \all{der Konflikt, der Streit, die Beleidigung, die Gewalt, wütend} & le conflit, la dispute, l'insulte, la violence, furieux \\ \hline
La médiation (Noms/Verbes) & \all{die Streitschlichter, zuhören, vermitteln, lösen, die Lösung, respektieren} & les médiateurs, écouter, arbitrer, résoudre, la solution, respecter \\ \hline
L'impératif (Formes) & \all{Redet miteinander! Vergesst nie!} & Parlez entre vous ! N'oubliez jamais ! \\ \hline
Le futur (Formes) & \all{wir werden organisieren, wir werden einladen, wir werden diskutieren, werden ... kommen} & nous organiserons, nous inviterons, nous discuterons, ... viendront \\ \hline
\end{tabularx}
\end{center}

Observe les deux dernières lignes : l'\textbf{impératif} sert à donner des conseils et des ordres (\all{Redet!} — « Parlez ! »), et le \textbf{futur} sert à annoncer des projets (\all{wir werden organisieren} — « nous organiserons »). Ce sont exactement les deux points de grammaire de cette leçon : le texte les contient déjà !

\begin{popfigure}
\begin{center}
\begin{tikzpicture}[mindmap, grow cyclic,
  every node/.style={concept, align=center, font=\small},
  root concept/.append style={concept color=popblue, font=\small\bfseries, minimum size=2.6cm},
  level 1 concept/.append style={level distance=3.4cm, sibling angle=90, minimum size=2.1cm}]
\node[root concept]{Friedens\-kultur}
  child[concept color=poporange]{ node[align=center]{Schule\\Streit\-schlichter} }
  child[concept color=popgreen]{ node[align=center]{Familie\\Dialog} }
  child[concept color=poppurple]{ node[align=center]{Gemeinde\\Toleranz} }
  child[concept color=popgold]{ node[align=center]{Politik\\Ver\-söhnung} };
\end{tikzpicture}
\end{center}
\legende{Carte mentale : la culture de la paix (\all{Friedenskultur}) vit dans quatre espaces — l'école, la famille, la communauté, la politique.}
\end{popfigure}

\courssub{Lecture orale : prononciation, intonation, fluidité}

Lisons maintenant le texte à voix haute. Quelques conseils de prononciation pour ce texte :

\begin{itemize}
\item \all{Frieden} se prononce « fri-denn » ; \all{Streitschlichter} « chtraït-chli-hteur » (le \all{sch} = « ch » français [ʃ], le \all{ch} après \all{i} = [ç], fricative palatale sourde).
\item \all{gewaltfrei} : « gué-valt-fraï » — le \all{w} allemand se prononce [v].
\item \all{Redet miteinander statt übereinander!} : phrase à l'impératif, donc \textbf{intonation descendante et ferme}, comme un conseil énergique : « ré-det mit-aï-nan-deur chtat ü-beur-aï-nan-deur ! »
\item Les phrases au futur (\all{Wir werden ein Friedensfest organisieren.}) se lisent avec une intonation \textbf{montante puis descendante}, qui exprime l'enthousiasme et le projet.
\item Respecte les pauses aux virgules et aux points : la fluidité vient du rythme, pas de la vitesse.
\end{itemize}

\begin{experience}[Lecture à voix haute en classe]
Travail en binômes : l'élève A lit les paragraphes 1 et 2, l'élève B lit les paragraphes 3 à 5. Chacun écoute l'autre et note deux points positifs (prononciation, intonation) et un point à améliorer. Puis la classe entière lit le conseil de Frau Berger en chœur : \all{Redet miteinander statt übereinander!}
\end{experience}

\begin{savaistu}[Les Streitschlichter : un modèle pour le Cameroun ?]
En Allemagne, de nombreux lycées forment des \all{Streitschlichter} (littéralement « apaiseurs de disputes ») : des élèves volontaires, formés à l'écoute active et à la médiation, qui aident leurs camarades à régler leurs conflits sans violence et sans punition. Ce modèle existe aussi en Autriche et en Suisse. Au Cameroun, les clubs de paix et les clubs UNESCO des lycées jouent un rôle proche : à Yaoundé ou à Douala, des élèves organisent des campagnes contre la violence scolaire et pour le vivre-ensemble. La paix, partout, commence par des gestes simples : écouter, parler, pardonner.
\end{savaistu}

\begin{exercice}[À toi de jouer ! — Compréhension]
\begin{enumerate}
\item Réponds en allemand avec une phrase complète : \all{Wo gibt es viele Konflikte? Wer hilft den Schülern? Was werden die Schüler nächstes Jahr organisieren?}
\item Relève dans le texte trois mots transparents (mots apparentés au français) et donne leur équivalent français.
\item Résume le texte en deux phrases en français.
\end{enumerate}
\end{exercice}

\begin{corrige}[Exercice « À toi de jouer ! »]
\begin{enumerate}
\item \all{In der Schule in Stuttgart gibt es viele Konflikte. Die Streitschlichter helfen den Schülern. Nächstes Jahr werden sie ein Friedensfest organisieren.} « À Stuttgart, il y a beaucoup de conflits dans l'école. Les médiateurs aident les élèves. L'année prochaine, ils organiseront une fête de la paix. »
\item Exemples : \all{der Konflikt} = le conflit ; \all{die Toleranz} = la tolérance ; \all{der Dialog} = le dialogue ; \all{organisieren} = organiser ; \all{diskutieren} = discuter ; \all{das Problem} = le problème.
\item Exemple de résumé : « Dans un lycée de Stuttgart, douze élèves médiateurs aident leurs camarades à résoudre leurs conflits sans violence. L'année prochaine, ils organiseront une fête de la paix, car la paix commence avec chacun de nous. »
\end{enumerate}
\end{corrige}

\begin{aretenir}[L'essentiel de cette section]
\begin{itemize}
\item Le texte \all{Frieden beginnt bei uns} suit une structure claire : \textbf{problème} (les conflits) → \textbf{solution} (la médiation) → \textbf{avenir} (l'espoir, la fête de la paix).
\item Le champ lexical de la paix : \all{der Frieden, die Toleranz, der Dialog, die Versöhnung, gewaltfrei} ; le champ lexical du conflit : \all{der Konflikt, der Streit, die Gewalt}.
\item On lit un texte allemand d'abord \textbf{globalement} (qui ? où ? quoi ?), sans dictionnaire, en s'aidant des mots transparents, puis \textbf{en détail} avec le glossaire.
\item Le texte contient déjà les deux points de grammaire de la leçon : l'\textbf{impératif} (\all{Redet! Vergesst!}) pour le conseil, et le \textbf{futur} (\all{wir werden organisieren}) pour les projets.
\end{itemize}
\end{aretenir}

% ============================================================
% SECTION 4 : GRAMMAIRE 1 — L'IMPÉRATIF ET LE CONSEIL
% ============================================================

\coursec{Grammaire 1 : L'impératif et le conseil (Der Imperativ)}

\courssub{Observation dans le texte}

Reprenons les deux phrases impératives du texte, prononcées par Frau Berger :
\all{„Redet miteinander statt übereinander!“}
\all{„Vergesst nie: Zuhören ist der erste Schritt zum Frieden.“}

\begin{exemplebox}[Analyse des formes]
\begin{itemize}
\item \all{Redet} $\leftarrow$ infinitif \all{reden} (parler). Terminaison \textbf{-et} $\rightarrow$ forme du \textbf{2e personne du pluriel (ihr)}.
\item \all{Vergesst} $\leftarrow$ infinitif \all{vergessen} (oublier). Terminaison \textbf{-t} (verbe fort, racine \all{vergess-}) $\rightarrow$ forme du \textbf{2e personne du pluriel (ihr)}.
\item Le verbe est en \textbf{première position}. Il n'y a \textbf{pas de sujet} exprimé (le « vous » est implicite).
\item \all{miteinander} (ensemble/les uns avec les autres) et \all{übereinander} (les uns des autres) sont des adverbes pronominaux placés après le verbe.
\item \all{nie} (jamais) est un adverbe de négation placé après le verbe (avant le deux-points).
\end{itemize}
\end{exemplebox}

\begin{propriete}[Formation de l'impératif (Présent)]
L'impératif existe pour trois personnes (du, ihr, Sie) et la 1re personne du pluriel (wir). Le verbe se place \textbf{en tête de phrase}.

\begin{center}
\begin{tabular}{|l|c|c|c|c|}\hline
\rowcolor{popblueL} \textbf{Personne} & \textbf{Règle de formation} & \textbf{Exemple : \all{reden}} & \textbf{Exemple : \all{arbeiten}} & \textbf{Exemple : \all{sein}} \\ \hline
\all{du} (tu) & Radical du présent (sans \textbf{-st}), \textbf{-e} facultatif (souvent omis) & \all{red(e)!} & \all{arbeite!} (\textbf{-e} obligatoire) & \all{sei!} (irrégulier) \\ \hline
\all{ihr} (vous [familier]) & Forme identique au présent \all{ihr} & \all{redet!} & \all{arbeitet!} & \all{seid!} (irrégulier) \\ \hline
\all{Sie} (Vous [formel]) & Infinitif + \all{Sie} (sujet après le verbe) & \all{reden Sie!} & \all{arbeiten Sie!} & \all{seien Sie!} \\ \hline
\all{wir} (nous) & Infinitif + \all{wir} (sujet après le verbe) & \all{reden wir!} & \all{arbeiten wir!} & \all{seien wir!} \\ \hline
\end{tabular}
\end{center}

\textbf{Règles particulières :}
\begin{itemize}
\item \textbf{Verbes à particule séparable :} La particule va \textbf{à la fin} de la phrase. \all{zu-hören} $\rightarrow$ \all{Hör zu!} / \all{Hört zu!} / \all{Hören Sie zu!}
\item \textbf{Verbes pronominaux (réfléchis) :} Le pronom réfléchi suit le verbe (sauf \all{wir}/\all{Sie} où il suit le sujet). \all{sich streiten} $\rightarrow$ \all{Streit dich!} / \all{Streitet euch!} / \all{Streiten Sie sich!}
\item \textbf{Négation :} \all{nicht} (ou \all{nie, niemals}) se place \textbf{avant la particule séparable} ou \textbf{à la fin} s'il n'y en a pas. \all{Redet nicht übereinander!} / \all{Vergesst das nicht!}
\end{itemize}
\end{propriete}

\begin{attention}[Pièges fréquents pour les francophones]
\begin{enumerate}
\item \textbf{Oublier le \textbf{-e} du \all{du} pour les radicaux en \textbf{-d, -t, -m, -n} (après consonne) :} \all{arbeite!}, \all{warte!}, \all{atme!}, \all{zeichne!} (mais \all{rede!}, \all{kaufe!} $\rightarrow$ \all{kauf!} courant à l'oral).
\item \textbf{Confondre \all{ihr} et \all{Sie} :} \all{ihr} = tutoiement pluriel (camarades, famille) ; \all{Sie} = vouvoiement (adultes, inconnus, professeurs). \textbf{Au lycée, entre élèves : \all{ihr}. Avec le prof : \all{Sie}.}
\item \textbf{Placer le sujet avant le verbe :} \textbf{Interdit} : \all{*Ihr redet!} (c'est un énoncé). \textbf{Correct} : \all{Redet ihr!} (non standard) ou simplement \all{Redet!}.
\item \textbf{Oublier la particule séparable :} \all{*Hör!} au lieu de \all{Hör zu!} (écoute !).
\end{enumerate}
\end{attention}

\begin{exemplebox}[Impératif et conseils (Rat geben)]
Pour donner un \textbf{conseil} amical, on utilise souvent :
\begin{itemize}
\item L'impératif \all{ihr} (entre pairs) : \all{Redet miteinander!}
\item L'impératif \all{Sie} (politesse) : \all{Reden Sie miteinander!}
\item La construction \textbf{\all{Man sollte} + infinitif} (on devrait) — plus doux : \all{Man sollte miteinander reden.}
\item La construction \textbf{\all{Du könntest} + infinitif} (tu pourrais) — suggestion : \all{Du könntest zuhören.}
\end{itemize}
\textit{Traduction :} « Parlez-vous les uns aux autres ! » / « Vous devriez vous parler. » / « Tu pourrais écouter. »
\end{exemplebox}

\begin{exoresolu}[Transformation à l'impératif]
\textbf{Énoncé :} Mets les phrases suivantes à l'impératif pour la personne indiquée.
\begin{enumerate}
\item \all{Du hörst zu.} (du) $\rightarrow$ \all{Hör zu!}
\item \all{Ihr vergesst das nicht.} (ihr) $\rightarrow$ \all{Vergesst das nicht!}
\item \all{Sie organisieren das Fest.} (Sie) $\rightarrow$ \all{Organisieren Sie das Fest!}
\item \all{Wir diskutieren darüber.} (wir) $\rightarrow$ \all{Diskutieren wir darüber!}
\item \all{Du streitest dich nicht.} (du) $\rightarrow$ \all{Streit dich nicht!}
\end{enumerate}
\tcblower
\textbf{Solution :}
1. \all{Hör zu!} (radical \all{hör-}, particule \all{zu} à la fin). 
2. \all{Vergesst das nicht!} (forme \all{ihr} = présent, \all{nicht} à la fin). 
3. \all{Organisieren Sie das Fest!} (infinitif + \all{Sie}). 
4. \all{Diskutieren wir darüber!} (infinitif + \all{wir}). 
5. \all{Streit dich nicht!} (radical \all{streit-}, pronom \all{dich}, \all{nicht} à la fin).
\end{exoresolu}

\begin{exercice}[Entraînement : L'impératif de la paix]
\begin{enumerate}
\item Conjugue les verbes à l'impératif \textbf{ihr} (conseils entre élèves) : \all{zuhören, vermitteln, lösen, respektieren, streiten}.
\item Transforme en impératif \textbf{Sie} (conseils à un adulte) : \all{Du hilfst mir. Du wartest. Du fängst an.}
\item Traduis en allemand (impératif \textbf{ihr}) : « Écoutez ! N'oubliez pas la tolérance ! Parlez ensemble ! Ne vous disputez pas ! »
\end{enumerate}
\end{exercice}

\begin{corrige}[Exercice « L'impératif de la paix »]
\begin{enumerate}
\item \all{Hört zu! Vermittelt! Löst! Respektiert! Streitet nicht!} (attention : \all{streiten} $\rightarrow$ \all{streitet}, mais négation \all{streitet nicht!})
\item \all{Helfen Sie mir! Warten Sie! Fangen Sie an!} (\all{anfangen} séparable $\rightarrow$ \all{fangen ... an}).
\item \all{Hört zu! Vergesst die Toleranz nicht! Redet miteinander! Streitet euch nicht!} (\all{sich streiten} pronominal $\rightarrow$ \all{streitet euch nicht}).
\end{enumerate}
\end{corrige}

% ============================================================
% SECTION 5 : GRAMMAIRE 2 — LE FUTUR (FUTUR I)
% ============================================================

\coursec{Grammaire 2 : Le futur (Das Futur I)}

\courssub{Observation dans le texte}

Reprenons les phrases du dernier paragraphe :
\all{Nächstes Jahr werden wir ein Friedensfest organisieren.}
\all{Wir werden alle Klassen einladen.}
\all{Wir werden Musik machen.}
\all{Wir werden über Toleranz und Dialog diskutieren.}
\all{Vielleicht werden auch andere Schulen zu uns kommen.}

\begin{exemplebox}[Analyse de la structure]
\begin{itemize}
\item \all{werden} est l'\textbf{auxiliaire du futur} (conjugué au présent). Il se place en \textbf{2e position} (phrase déclarative).
\item L'\textbf{infinitif du verbe de sens} (\all{organisieren, einladen, machen, diskutieren, kommen}) va \textbf{à la fin de la phrase} (position de clôture).
\item \all{Nächstes Jahr}, \all{Vielleicht} sont des compléments de temps/de modalité en \textbf{1re position} (thématisation).
\item \all{werden} se conjugue : \all{ich werde, du wirst, er/sie/es wird, wir werden, ihr werdet, sie/Sie werden}.
\end{itemize}
\end{exemplebox}

\begin{propriete}[Formation du Futur I (\all{Futur I})]
\textbf{Forme :} \all{werden} (présent) + \textbf{Infinitif du verbe principal} (à la fin).

\begin{center}
\begin{tabular}{|l|l|l|}\hline
\rowcolor{popblueL} \textbf{Personne} & \textbf{\all{werden} (auxiliaire)} & \textbf{Exemple : \all{organisieren}} \\ \hline
\all{ich} & \all{werde} & \all{werde organisieren} \\ \hline
\all{du} & \all{wirbst} & \all{wirbst organisieren} \\ \hline
\all{er / sie / es} & \all{wird} & \all{wird organisieren} \\ \hline
\all{wir} & \all{werden} & \all{werden organisieren} \\ \hline
\all{ihr} & \all{werdet} & \all{werdet organisieren} \\ \hline
\all{sie / Sie} & \all{werden} & \all{werden organisieren} \\ \hline
\end{tabular}
\end{center}

\textbf{Verbes à particule séparable :} L'infinitif reste \textbf{entier} à la fin (particule collée). \all{an-fangen} $\rightarrow$ \all{wir werden anfangen}. \all{ein-laden} $\rightarrow$ \all{wir werden einladen}.
\textbf{Verbes pronominaux :} Pronom réfléchi avant l'infinitif. \all{sich versöhnen} $\rightarrow$ \all{wir werden uns versöhnen}.
\textbf{Négation :} \all{nicht} avant l'infinitif (ou avant le complément d'objet). \all{Wir werden nicht streiten.}
\end{propriete}

\begin{propriete}[Emplois du Futur I (niveau A2/B1)]
\begin{enumerate}
\item \textbf{Action future certaine ou planifiée} (souvent avec indicateur de temps : \all{morgen, nächstes Jahr, bald, in Zukunft}) : \all{Nächstes Jahr werden wir ein Friedensfest organisieren.}
\item \textbf{Intention, promesse, décision prise au moment de parler} : \all{Ich werde dir helfen!} (Je t'aiderai ! — promesse).
\item \textbf{Supposition, probabilité sur le présent ou le futur} (souvent avec \all{wohl, wahrscheinlich, vielleicht}) : \all{Er wird wohl krank sein.} (Il est probablement malade.) $\rightarrow$ \textit{Attention : ce sens est plus avancé (B1), à reconnaître en lecture.}
\end{enumerate}
\end{propriete}

\begin{attention}[Futur I vs Présent + indicateur de temps]
\cle{Très important} : En allemand courant, le \textbf{présent + indicateur de temps} remplace souvent le Futur I pour exprimer un projet sûr.
\begin{center}
\begin{tabularx}{\linewidth}{|X|X|}\hline
\rowcolor{popblueL} \textbf{Futur I (formel, écrit, promesse, insistance)} & \textbf{Présent + temps (courant, oral, planning)} \\ \hline
\all{Wir werden nächstes Jahr ein Fest organisieren.} & \all{Wir organisieren nächstes Jahr ein Fest.} \\ \hline
\all{Ich werde morgen nach Yaoundé fahren.} & \all{Ich fahre morgen nach Yaoundé.} \\ \hline
\end{tabularx}
\end{center}
\textbf{Conseil :} À l'écrit (examen, rédaction), utilise le \textbf{Futur I} pour montrer que tu maîtrises la structure. À l'oral, le présent est naturel.
\end{attention}

\begin{attention}[Faux ami : \all{werden} $\neq$ « vouloir »]
\all{werden} comme auxiliaire = « will / shall » (marque du futur).
\all{werden} comme verbe principal = « devenir » (\all{Er wird Arzt} = Il devient médecin).
\all{wollen} = « vouloir » (\all{Ich will Frieden} = Je veux la paix).
\textbf{Ne confonds pas :} \all{Ich werde Frieden schaffen} (Je créerai la paix — futur) vs \all{Ich will Frieden schaffen} (Je veux créer la paix — intention).
\end{attention}

\begin{exoresolu}[Conjugaison au Futur I]
\textbf{Énoncé :} Conjugue les verbes entre parenthèses au Futur I.
\begin{enumerate}
\item \all{Die Schüler \_\_\_\_\_\_\_\_ (organisieren) ein Friedensfest.}
\item \all{Ich \_\_\_\_\_\_\_\_ (nicht / streiten) mehr.}
\item \all{Wir \_\_\_\_\_\_\_\_ (uns / versöhnen).}
\item \all{Was \_\_\_\_\_\_\_\_ du \_\_\_\_\_\_\_\_ (machen) nächstes Jahr?}
\end{enumerate}
\tcblower
\textbf{Solution :}
1. \all{werden organisieren} (3e pers. pl. \all{sie} = \all{werden}). 
2. \all{werde nicht streiten} (1re pers. sg. \all{ich werde}, \all{nicht} avant infinitif). 
3. \all{werden uns versöhnen} (1re pers. pl. \all{wir werden}, pronom \all{uns} avant infinitif). 
4. \all{wirst ... machen} (2e pers. sg. \all{du wirst}, infinitif \all{machen} à la fin).
\end{exoresolu}

\begin{exercice}[Entraînement : Projets pour la paix]
\begin{enumerate}
\item Conjugue au Futur I : \all{ich (planen), du (teilnehmen), wir (diskutieren), ihr (einladen), sie (kommen)}.
\item Traduis en allemand (Futur I) :
\begin{itemize}
\item L'année prochaine, nous organiserons une grande fête.
\item Nous inviterons tous les élèves de l'école.
\item Nous ne nous disputerons plus.
\item Est-ce que vous (Sie) viendrez aussi ?
\end{itemize}
\item Complète avec \all{werden} + infinitif : \all{Nächstes Jahr \_\_\_\_\_\_\_\_ die Schule ruhiger \_\_\_\_\_\_\_\_ (sein).} (Indice : \all{sein} à la fin).
\end{enumerate}
\end{exercice}

\begin{corrige}[Exercice « Projets pour la paix »]
\begin{enumerate}
\item \all{ich werde planen, du wirst teilnehmen, wir werden diskutieren, ihr werdet einladen, sie werden kommen}.
\item \all{Nächstes Jahr werden wir ein großes Fest organisieren. Wir werden alle Schüler der Schule einladen. Wir werden nicht mehr streiten. Werden Sie auch kommen?}
\item \all{Nächstes Jahr wird die Schule ruhiger sein.} (\all{die Schule} sg. $\rightarrow$ \all{wird}, \all{sein} à la fin).
\end{enumerate}
\end{corrige}

% ============================================================
% SECTION 6 : MÉTHODE DU COMMENTAIRE ET PRODUCTION GUIDÉE
% ============================================================

\coursec{Méthode du commentaire et production guidée}

\courssub{Méthode : Commenter un texte et produire un court paragraphe}

\begin{methode}[Écrire un résumé commenté (Zusammenfassung mit Stellungnahme) — Niveau A2]
\begin{enumerate}
\item \textbf{Introduction (1 phrase) :} Titre, source, thème. \all{Der Text „Frieden beginnt bei uns“ aus einer Schülerzeitung handelt von der Gewaltprävention an einer Schule in Stuttgart.}
\item \textbf{Résumé du contenu (3-4 phrases, Présent) :} Structure : Problème $\rightarrow$ Solution $\rightarrow$ Résultats $\rightarrow$ Avenir. Utilise des connecteurs : \all{Zuerst / Zunächst ... Dann ... Schließlich / Am Ende ...}
\item \textbf{Opinion personnelle / Lien avec son contexte (2-3 phrases, Futur I ou Présent + \all{meiner Meinung nach}) :} \all{Meiner Meinung nach ist das Projekt sehr gut. In meiner Schule in Yaoundé könnten wir auch Streitschlichter ausbilden. Wir werden nächstes Jahr ein Friedensfest planen.}
\item \textbf{Conclusion (1 phrase) :} \all{Frieden beginnt wirklich bei uns.}
\end{enumerate}
\end{methode}

\begin{attention}[Critères de réussite pour l'écrit (Bac / Contrôle)]
\begin{itemize}
\item \textbf{Ordre des mots (Wortstellung) :} Verbe conjugué en 2e position (phrase déclarative), infinitif/particule à la fin. \checkmark
\item \textbf{Majuscules :} Tous les noms (\all{der Frieden, die Schule, wir}). \checkmark
\item \textbf{Accords :} Sujet + verbe (\all{die Schüler helfen}, pas \all{*die Schüler hilft}). \checkmark
\item \textbf{Cas :} \all{helfen} + Dativ (\all{den Schülern}), \all{respektieren} + Akkusativ (\all{einander}). \checkmark
\item \textbf{Orthographe :} Umlaute (\all{ä, ö, ü}), \all{ß} (après voyelle longue/diphtongue : \all{Straße, groß, weiß}). \checkmark
\end{itemize}
\end{attention}

\courssub{Production écrite guidée : Mon projet pour la paix}

\begin{experience}[Atelier d'écriture : « Mon engagement pour la paix »]
\textbf{Consigne :} Rédige un court paragraphe (60-80 mots) en allemand pour le journal de ton lycée. Tu es délégué(e) de classe. Utilise :
\begin{itemize}
\item Le vocabulaire de la leçon (\all{Frieden, Konflikt, Toleranz, Dialog, Streitschlichter, gewaltfrei}).
\item L'\textbf{impératif} (\all{ihr} ou \all{Sie}) pour donner 2 conseils à tes camarades.
\item Le \textbf{Futur I} pour annoncer 2 actions concrètes que tu vas organiser cette année.
\end{itemize}
\textbf{Aide au démarrage (Gerüst) :}
\all{Liebe Mitschülerinnen und Mitschüler,}
\all{In unserer Schule gibt es manchmal \_\_\_\_\_\_\_\_ (Konflikte). Das ist nicht gut.}
\all{\_\_\_\_\_\_\_\_ (Impératif: Zuhören / Reden) \_\_\_\_\_\_\_\_!}
\all{\_\_\_\_\_\_\_\_ (Impératif: Nicht streiten / Respektieren) \_\_\_\_\_\_\_\_!}
\all{Deshalb \_\_\_\_\_\_\_\_ ich \_\_\_\_\_\_\_\_ (Futur: ein Projekt starten / Streitschlichter ausbilden).}
\all{Auch \_\_\_\_\_\_\_\_ wir \_\_\_\_\_\_\_\_ (Futur: ein Friedensfest feiern / über Toleranz sprechen).}
\all{Denn Frieden beginnt bei uns!}
\all{Dein/Euer \_\_\_\_\_\_\_\_ (Name)}
\end{experience}

\begin{exemplebox}[Production type (Corrigé modèle)]
\all{Liebe Mitschülerinnen und Mitschüler,}
\all{In unserer Schule gibt es manchmal Streit und Beleidigungen. Das ist nicht gut.}
\all{Hört einander zu! Redet miteinander statt übereinander!}
\all{Deshalb werde ich ein Friedensprojekt starten. Auch werden wir nächstes Jahr ein Friedensfest feiern und über Toleranz diskutieren.}
\all{Denn Frieden beginnt bei uns!}
\all{Euer Paul}
\end{exemplebox}

\begin{exercice}[Expression écrite — Évaluation formative]
Choisis \textbf{un} des deux sujets et rédige 60-80 mots en allemand.
\begin{enumerate}
\item \textbf{Sujet A (Conseils) :} Tu es médiateur (\all{Streitschlichter}). Écris une « Lettre ouverte aux élèves » (\all{Offener Brief an die Schüler}) pour expliquer ton rôle et donner 3 conseils impératifs pour éviter la violence. Utilise : \all{zuhören, vermitteln, gewaltfrei lösen, respektieren}.
\item \textbf{Sujet B (Projets) :} Tu es président du Club Paix. Écris un article pour le journal scolaire annonçant le programme de l'année prochaine (Futur I) : fête de la paix, atelier dialogue, invitation d'autres écoles, campagne « École sans violence ».
\end{enumerate}
\end{exercice}

\begin{corrige}[Propositions de correction — Sujets A et B]
\textbf{Sujet A (Modèle) :}
\all{Offener Brief an die Schüler}
\all{Liebe Mitschüler, ich bin Streitschlichter. Meine Aufgabe ist zuhören und vermitteln. Bitte: Hört einander zu! Redet miteinander, nicht übereinander! Löst Probleme gewaltfrei! Respektiert euch! Nur so schaffen wir Frieden. Euer Max}

\textbf{Sujet B (Modèle) :}
\all{Unser Friedensprogramm nächstes Jahr}
\all{Nächstes Jahr werden wir ein großes Friedensfest organisieren. Wir werden Workshops zum Dialog anbieten. Andere Schulen werden wir einladen. Wir werden eine Kampagne „Schule ohne Gewalt“ starten. Wir werden für Toleranz und Versöhnung kämpfen. Frieden beginnt bei uns!}
\end{corrige}

\begin{aretenir}[Bilan de la leçon ALA26 — La promotion de la culture de la paix]
\begin{itemize}
\item \textbf{Vocabulaire clé :} \all{der Frieden, der Konflikt, die Toleranz, der Dialog, die Versöhnung, die Streitschlichter, gewaltfrei, die Lösung, das Friedensfest}. (Articles + Pluriels maîtrisés).
\item \textbf{Grammaire 1 — Impératif :} Formes \all{du} (radical), \all{ihr} (présent), \all{Sie/wir} (infinitif + pronom). Place verbe = 1re. Particule séparable à la fin. Usage : conseils, ordres, consignes.
\item \textbf{Grammaire 2 — Futur I :} \all{werden} (présent) + Infinitif (fin). Emplois : projets, promesses, prédictions. Alternative orale : Présent + indicateur de temps.
\item \textbf{Compétences validées :} Comprendre un texte authentique (global/détaillé), extraire le lexique thématique, appliquer l'impératif et le futur dans des phrases et un paragraphe cohérent, produire un texte argumentatif court (conseils / projets).
\end{itemize}
\end{aretenir}

\coursec{Compréhension du texte}

Après avoir découvert le texte support \emph{Frieden beginnt bei uns} dans la section précédente, nous allons maintenant l'exploiter en profondeur. La compréhension de l'écrit ne se limite pas à saisir le sens global : elle exige une lecture active, structurée, capable de restituer l'information avec ses propres mots, de justifier ses réponses par le texte et de synthétiser l'essentiel. C'est précisément ce que demandent les épreuves du Baccalauréat (épreuve écrite de compréhension et expression). Nous allons progresser par étapes : globale, détaillée, personnelle, reformulation et résumé.

\courssub{Le texte de référence : \emph{Frieden beginnt bei uns}}

Voici le texte intégral qui servira de base à tous les exercices de cette section. Lisez-le une première fois sans dictionnaire, puis une seconde fois en vous aidant du glossaire ci-dessous. (Les numéros de ligne Z. sont indiqués pour les citations.)

\begin{textebox}[Texte support : Frieden beginnt bei uns]
An der Goethe-Schule in Yaoundé gibt es seit diesem Schuljahr ein neues Projekt: die Streitschlichter. Zwölf Schülerinnen und Schüler der Klassen 10 und 11 wurden dafür speziell ausgebildet. Ihre Aufgabe ist es, bei Konflikten zwischen Mitschülern zu vermitteln. „Wir hören beiden Seiten zu und suchen gemeinsam eine Lösung“, erklärt Amadou, ein Streitschlichter der 11. Klasse. „Wir entscheiden nicht, wer Recht hat, wir helfen nur beim Reden.“

Frau Berger, die betreuende Lehrerin, betont: „Gewalt löst keine Probleme. Wer sprechen lernt, muss nicht schlagen.“ Die Streitschlichter sind in den Pausen im Raum 12 zu finden. Jeder Konflikt wird vertraulich behandelt. „Am Anfang war ich skeptisch“, sagt Kevin, ein Schüler der 9. Klasse. „Aber nach einem Streit mit meinem Freund haben wir uns dank der Streitschlichter wieder vertragen. Das hat mir sehr geholfen.“

Nächstes Jahr will die Schule eine Projektwoche „Kultur des Friedens“ organisieren. Geplant sind Workshops zu Toleranz, interkulturellem Dialog und Gewaltprävention. Auch Eltern und Vereine aus dem Viertel sollen eingebunden werden. „Frieden beginnt im Kleinen“, sagt die Schulleiterin, Frau Mengue. „Wenn wir hier in der Schule respektvoll miteinander umgehen, tragen wir zu einer friedlicheren Gesellschaft bei.“
\end{textebox}

\begin{definition}[Glossaire du texte \emph{Frieden beginnt bei uns}]
\begin{itemize}
    \item \cle{das Schuljahr, die Schuljahre} : l'année scolaire
    \item \cle{das Projekt, die Projekte} : le projet
    \item \cle{der Streitschlichter, die Streitschlichter} : le médiateur (scolaire), celui qui aide à résoudre un conflit
    \item \cle{ausbilden} (bildet aus, hat ausgebildet) : former, éduquer (verbe à particule séparable : \all{ausbilden})
    \item \cle{der Konflikt, die Konflikte} : le conflit
    \item \cle{vermitteln} : médiatiser, faire de la médiation
    \item \cle{beide Seiten} : les deux parties / les deux camps
    \item \cle{die Lösung, die Lösungen} : la solution
    \item \cle{entscheiden} : décider
    \item \cle{recht haben} : avoir raison
    \item \cle{die betreuende Lehrerin} : l'enseignante encadrante / référente
    \item \cle{betonen} : souligner, insister sur
    \item \cle{die Gewalt} : la violence
    \item \cle{lösen} (löst, hat gelöst) : résoudre
    \item \cle{die Pause, die Pausen} : la récréation
    \item \cle{vertraulich} : confidentiel
    \item \cle{behandeln} : traiter
    \item \cle{skeptisch} : sceptique
    \item \cle{sich vertragen} (verträgt sich, hat sich vertragen) : se réconcilier, faire la paix (verbe à particule inséparable $\to$ pas de \emph{ge-} au Partizip II)
    \item \cle{dank} (+ Génitif) : grâce à
    \item \cle{die Projektwoche, die Projektwochen} : la semaine de projet
    \item \cle{die Kultur des Friedens} : la culture de la paix
    \item \cle{der Workshop, die Workshops} : l'atelier
    \item \cle{die Toleranz} : la tolérance
    \item \cle{der interkulturelle Dialog} : le dialogue interculturel
    \item \cle{die Gewaltprävention} : la prévention de la violence
    \item \cle{einbinden} (bindet ein, hat eingebunden) : impliquer, associer (verbe à particule séparable)
    \item \cle{der Viertel, die Viertel} : le quartier
    \item \cle{die Schulleiterin, die Schulleiterinnen} : la directrice (d'école)
    \item \cle{respektvoll} : respectueux
    \item \cle{der Umgang, die Umgänge} : le comportement, les rapports (avec qqn)
    \item \cle{beitragen zu} (+ Datif) : contribuer à
    \item \cle{die Gesellschaft, die Gesellschaften} : la société
\end{itemize}
\end{definition}

\courssub{Compréhension globale}

La compréhension globale vise à identifier le \cle{sujet}, les \cle{acteurs}, le \cle{lieu}, le \cle{moment} et l'\cle{intention} du texte. On répond en une ou deux phrases complètes.

\begin{exemplebox}[Questions globales et réponses modèles]
\textbf{1. Worum geht es in diesem Text?} (De quoi parle ce texte ?)\\
\all{Es geht um ein neues Projekt an der Goethe-Schule in Yaoundé: Schüler werden zu Streitschlichtern ausgebildet, um Konflikte friedlich zu lösen.} \\
« Il s'agit d'un nouveau projet à la Goethe-Schule de Yaoundé : des élèves sont formés comme médiateurs pour résoudre les conflits pacifiquement. »

\textbf{2. Wer sind die Streitschlichter?} (Qui sont les médiateurs ?)\\
\all{Die Streitschlichter sind zwölf Schülerinnen und Schüler der Klassen 10 und 11, die speziell ausgebildet wurden, um bei Streit zu vermitteln.} \\
« Les médiateurs sont douze élèves de 10e et 11e classes, spécialement formés pour médiatiser lors de disputes. »

\textbf{3. Was ist das Ziel des Projekts?} (Quel est l'objectif du projet ?)\\
\all{Das Ziel ist, Gewalt durch Gespräch zu ersetzen und eine Kultur des Friedens in der Schule zu fördern.} \\
« L'objectif est de remplacer la violence par le dialogue et de promouvoir une culture de la paix dans l'école. »
\end{exemplebox}

\begin{methode}[Comment répondre à une question de compréhension globale]
\begin{enumerate}
    \item \textbf{Lisez la question} et identifiez le mot-clé (Wer, Was, Worum, Wo, Wann, Warum).
    \item \textbf{Repérez l'information} dans le texte (souvent dans le premier ou le dernier paragraphe).
    \item \textbf{Formulez une phrase complète} en allemand : Sujet + Verbe conjugué (2e position) + Compléments.
    \item \textbf{Reprenez le vocabulaire du texte} mais changez la structure si possible (reformulation).
    \item \textbf{Vérifiez} : la réponse répond-elle exactement à la question posée ?
\end{enumerate}
\end{methode}

\courssub{Compréhension détaillée}

Ici, il faut aller chercher des informations précises, souvent implicites ou dispersées dans le texte. Les réponses doivent être justifiées par une \cle{citation} (numéro de ligne ou copier-coller court).

\begin{exemplebox}[Questions détaillées et réponses justifiées]
\textbf{1. Was machen die Streitschlichter genau?} (Que font exactement les médiateurs ?)\\
\all{Sie hören beiden Seiten zu und helfen, gemeinsam eine Lösung zu finden, ohne zu entscheiden, wer Recht hat.} (Z. 4--5)\\
« Ils écoutent les deux parties et aident à trouver une solution ensemble, sans décider qui a raison. »

\textbf{2. Was wird die Schule nächstes Jahr organisieren?} (Qu'est-ce que l'école organisera l'an prochain ?)\\
\all{Die Schule will eine Projektwoche „Kultur des Friedens“ mit Workshops zu Toleranz, interkulturellem Dialog und Gewaltprävention organisieren.} (Z. 13--15)\\
« L'école veut organiser une semaine de projet "Culture de la paix" avec des ateliers sur la tolérance, le dialogue interculturel et la prévention de la violence. »

\textbf{3. Was sagt Frau Berger zur Gewalt?} (Que dit Mme Berger sur la violence ?)\\
\all{Frau Berger betont: „Gewalt löst keine Probleme. Wer sprechen lernt, muss nicht schlagen.“} (Z. 7--8)\\
« Mme Berger souligne : "La violence ne résout pas les problèmes. Qui apprend à parler n'a pas besoin de frapper." »
\end{exemplebox}

\begin{exercice}[Vrai ou Faux ? Justifiez par une citation du texte (numéro de ligne)]
Déterminez si les affirmations suivantes sont \all{Richtig} (Vrai) ou \all{Falsch} (Faux). Citez la ligne qui justifie votre choix.
\begin{enumerate}
    \item Die Streitschlichter entscheiden, wer schuld ist.
    \item Kevin hat die Streitschlichter nie benutzt.
    \item Die Projektwoche findet in diesem Jahr statt.
    \item Eltern und Vereine sollen bei der Projektwoche mitmachen.
    \item Frau Mengue glaubt, dass Frieden in der Schule keine Auswirkung auf die Gesellschaft hat.
\end{enumerate}
\end{exercice}

\begin{corrige}[Exercice Vrai/Faux]
\begin{enumerate}
    \item \textbf{Falsch.} Z. 5: „Wir entscheiden nicht, wer Recht hat…“
    \item \textbf{Falsch.} Z. 10--11: „Aber nach einem Streit mit meinem Freund haben wir uns dank der Streitschlichter wieder vertragen.“
    \item \textbf{Falsch.} Z. 13: „Nächstes Jahr will die Schule eine Projektwoche … organisieren.“
    \item \textbf{Richtig.} Z. 14--15: „Auch Eltern und Vereine aus dem Viertel sollen eingebunden werden.“
    \item \textbf{Falsch.} Z. 16--17: „Wenn wir hier in der Schule respektvoll miteinander umgehen, tragen wir zu einer friedlicheren Gesellschaft bei.“
\end{enumerate}
\end{corrige}

\courssub{Questions à réponse personnelle (Persönliche Stellungnahme)}

Ces questions vérifient votre capacité à exprimer votre opinion en allemand en réutilisant le vocabulaire du thème. Il n'y a pas de « bonne » réponse unique, mais la forme doit être correcte.

\begin{exemplebox}[Exemples de réponses personnelles]
\textbf{1. Gibt es Konflikte in deiner Schule? Wie sehen sie aus?} (Y a-t-il des conflits dans ton école ? À quoi ressemblent-ils ?)\\
\all{Ja, es gibt manchmal Konflikte, zum Beispiel wegen Missverständnissen oder beim Fußballspielen auf dem Hof. Meistens sind es kleine Streitereien, die schnell vorbei sind.} \\
« Oui, il y a parfois des conflits, par exemple à cause de malentendus ou en jouant au foot dans la cour. Ce sont surtout de petites disputes qui finissent vite. »

\textbf{2. Wie kann man Konflikte in der Schule am besten lösen?} (Comment résoudre au mieux les conflits à l'école ?)\\
\all{Am besten ist es, ruhig zu reden und zuzuhören. Ein Streitschlichter oder ein Lehrer kann helfen, wenn man allein keine Lösung findet. Gewalt ist nie eine Lösung.} \\
« Le mieux est de parler calmement et d'écouter. Un médiateur ou un prof peut aider si on ne trouve pas de solution seul. La violence n'est jamais une solution. »

\textbf{3. Was bedeutet „Kultur des Friedens“ für dich persönlich?} (Que signifie "culture de la paix" pour toi personnellement ?)\\
\all{Für mich bedeutet das, dass man andere respektiert, auch wenn sie anders sind. Man soll Probleme mit Worten lösen, nicht mit Fäusten. Toleranz und Dialog sind wichtig.} \\
« Pour moi, cela signifie qu'on respecte les autres, même s'ils sont différents. On doit résoudre les problèmes par les mots, pas par les poings. Tolérance et dialogue sont importants. »
\end{exemplebox}

\courssub{Entraînement à la reformulation (Umformulieren)}

Au Bac, \textbf{recopier mot pour mot une phrase du texte est pénalisé}. Il faut prouver que vous avez compris en reformulant : changer la structure syntaxique, utiliser des synonymes, passer de l'actif au passif, transformer le discours direct en indirect.

\begin{methode}[La reformulation en 3 étapes]
\begin{enumerate}
    \item \textbf{Identifiez l'information cible} dans le texte (soulignez-la).
    \item \textbf{Changez le point de vue ou la structure} :
    \begin{itemize}
        \item Discours direct $\to$ Discours indirect (\all{„Wir helfen“} $\to$ \all{Sie sagten, dass sie helfen.})
        \item Actif $\to$ Passif (\all{Schüler ausbilden} $\to$ \all{Schüler werden ausgebildet})
        \item Nominalisation (\all{Sie vermitteln} $\to$ \all{Ihre Vermittlung ist wichtig})
        \item Synonymes (\all{streiten} $\to$ \all{sich streiten / einen Konflikt haben})
    \end{itemize}
    \item \textbf{Vérifiez la grammaire} (cas, conjugaison, ordre des mots) et le sens.
\end{enumerate}
\end{methode}

\begin{exoresolu}[Reformulation guidée]
\textbf{Consigne :} Reformulez la phrase suivante en une phrase complète différente, sans changer le sens.
\\ \textit{Original (Z. 7--8):} \all{Frau Berger betont: „Gewalt löst keine Probleme. Wer sprechen lernt, muss nicht schlagen.“}

\textbf{Étape 1 :} Information cible = La violence ne résout rien ; parler évite de frapper.
\textbf{Étape 2 :} Passage au discours indirect + nominalisation.
\textbf{Étape 3 :} Rédaction.

\all{Frau Berger betont, dass Gewalt keine Probleme löst und dass man nicht schlagen muss, wenn man sprechen kann.}
\\ (Mme Berger souligne que la violence ne résout pas les problèmes et qu'on n'a pas besoin de frapper si on sait parler.)

\textbf{Autre possibilité (passif + nominale) :}
\all{Laut Frau Berger werden Probleme nicht durch Gewalt gelöst, sondern durch das Sprechenlernen vermieden.}
\\ (Selon Mme Berger, les problèmes ne sont pas résolus par la violence, mais évités par l'apprentissage du parler.)
\end{exoresolu}

\begin{attention}[Piège du Bac : Copier-coller = 0 point]
\begin{itemize}
    \item Si la consigne demande \all{„in eigenen Worten“} (avec vos mots) ou \all{„zusammenfassen“} (résumer), toute phrase copiée intégralement du texte \textbf{ne rapporte aucun point}.
    \item Même sans consigne explicite, l'usage veut que vous \textbf{reformuliez}.
    \item \textbf{Astuce :} Gardez les \cle{mots-clés techniques} (Streitschlichter, Gewaltprävention, interkultureller Dialog) mais changez les \cle{verbes}, la \cle{structure de phrase} et les \cle{connecteurs}.
\end{itemize}
\end{attention}

\courssub{Le résumé guidé (Zusammenfassung)}

Le résumé doit condenser l'essentiel du texte en 3 à 4 phrases, dans un ordre logique (chronologique ou thématique), en utilisant des \cle{connecteurs} (Konjunktionen / Adverbien) pour marquer l'enchaînement des idées.

\begin{propriete}[Connecteurs temporels et logiques pour le résumé]
\begin{center}
\begin{tabularx}{\linewidth}{|X|X|X|}
\hline
\rowcolor{popblueL} \textbf{Allemand} & \textbf{Français} & \textbf{Fonction / Exemple} \\
\hline
\cle{zuerst} & d'abord, en premier & Introduction du sujet / 1re action : \all{Zuerst gibt es das Streitschlichter-Projekt.} \\
\hline
\cle{dann} / \cle{anschließend} & ensuite, puis & Suite chronologique : \all{Dann werden die Schüler ausgebildet.} \\
\hline
\cle{außerdem} / \cle{zusätzlich} & de plus, en outre & Ajout d'information : \all{Außerdem gibt es eine Projektwoche.} \\
\hline
\cle{deshalb} / \cle{deswegen} & c'est pourquoi, donc & Conséquence / Cause-effet : \all{Deshalb gibt es weniger Streit.} \\
\hline
\cle{schließlich} / \cle{zuletzt} & finalement, enfin & Conclusion / Dernière étape : \all{Schließlich will die Schule die Eltern einbinden.} \\
\hline
\cle{weil} / \cle{da} & parce que, comme & Cause (verbe à la fin) : \all{Das Projekt gelingt, weil alle mitmachen.} \\
\hline
\cle{obwohl} & bien que, quoique & Concession (verbe à la fin) : \all{Obwohl Kevin skeptisch war, half es ihm.} \\
\hline
\end{tabularx}
\end{center}
\end{propriete}

\begin{experience}[Atelier résumé guidé : À vous de jouer !]
\textbf{Consigne :} Écrivez un résumé du texte en \textbf{4 phrases} en utilisant \textbf{au moins 4 connecteurs différents} du tableau ci-dessus. Suivez ce plan :
\begin{enumerate}
    \item Phrase 1 (Sujet + lancement du projet) : \textit{zuerst} ...
    \item Phrase 2 (Rôle des médiateurs + exemple Kevin) : \textit{dann / außerdem} ...
    \item Phrase 3 (Réaction de l'encadrante + principe) : \textit{weil / deshalb} ...
    \item Phrase 4 (Projet futur + portée) : \textit{schließlich} ...
\end{enumerate}
\textbf{Aide lexicale :} \all{das Streitschlichter-Projekt, ausbilden, vermitteln, sich vertragen, die Projektwoche, die Gewaltprävention, einbinden, beitragen zu.}
\end{experience}

\begin{exoresolu}[Proposition de résumé corrigé]
\all{Zuerst wurde an der Goethe-Schule in Yaoundé ein Streitschlichter-Projekt eingeführt. Dann wurden zwölf Schüler ausgebildet, um bei Konflikten zu vermitteln; außerdem halfen sie Kevin, sich mit seinem Freund zu vertragen. Deshalb betont Frau Berger, dass Gewalt keine Probleme löst. Schließlich plant die Schule für nächstes Jahr eine Projektwoche zur Gewaltprävention, um die Kultur des Friedens auch mit Eltern und Vereinen zu stärken.}
\\
\textit{Traduction :} D'abord, un projet de médiateurs a été introduit à la Goethe-Schule de Yaoundé. Ensuite, douze élèves ont été formés pour médiatiser lors de conflits ; de plus, ils ont aidé Kevin à se réconcilier avec son ami. C'est pourquoi Mme Berger souligne que la violence ne résout pas les problèmes. Finalement, l'école prévoit pour l'an prochain une semaine de projet sur la prévention de la violence, afin de renforcer la culture de la paix aussi avec les parents et les associations.
\end{exoresolu}

\begin{popfigure}
\begin{tikzpicture}[
    node distance=1.2cm and 1.5cm,
    box/.style={draw=popblue, thick, rounded corners, fill=popblueL, text width=3.2cm, align=center, font=\small},
    arrow/.style={-Stealth, thick, popdark}
]
\node[box, align=center] (globale){\textbf{Compréhension\\ Globale}\\\textit{Worum geht es?\\ Wer? Was? Ziel?}};
\node[box, right=of globale, align=center] (detailee){\textbf{Compréhension\\ Détaillée}\\\textit{Précisions,\\ Citations,\\ Vrai/Faux}};
\node[box, right=of detailee, align=center] (personnelle){\textbf{Réponse\\ Personnelle}\\\textit{Meinung,\\ Erfahrung,\\ Vokabeltransfer}};
\node[box, below=of detailee, align=center] (resume){\textbf{Résumé /\\ Reformulation}\\\textit{Zusammenfassung,\\ Connecteurs,\\ Eigene Worte}};

\draw[arrow] (globale) -- (detailee);
\draw[arrow] (detailee) -- (personnelle);
\draw[arrow] (detailee) -- (resume);
\draw[arrow] (personnelle) -- (resume);
\draw[arrow] (globale) |- (resume);
\end{tikzpicture}
\legende{Organigramme : progression des types de questions de compréhension (globale $\to$ détaillée $\to$ personnelle $\to$ résumé/reformulation).}
\end{popfigure}

\begin{aretenir}[Check-list pour l'examen de compréhension écrite]
\begin{itemize}
    \item \textbf{Lisez le texte 2 fois} : 1x global, 1x détaillé (surlignez mots-clés, chiffres, noms, connecteurs).
    \item \textbf{Analysez la consigne} : Allemand ou français ? Phrase complète ou mot unique ? Citation exigée ?
    \item \textbf{Répondez en phrases complètes} (Sujet + Verbe 2e position), reformulées (\emph{eigene Worte}).
    \item \textbf{Justifiez} (Vrai/Faux) par le numéro de ligne ou la citation exacte entre guillemets allemands „…“.
    \item \textbf{Résumé} : 3--4 phrases, ordre logique, \textbf{4 connecteurs minimum} (\all{zuerst, dann, deshalb, schließlich}...), pas de détails superflus.
    \item \textbf{Relisez-vous} : Orthographe (majuscules noms, Umlaute, ß), grammaire (cas, conjugaison), sens.
\end{itemize}
\end{aretenir}

\coursec{Lexique thématique : Frieden und Konflikt}

Après avoir compris le texte \all{Frieden beginnt bei uns} dans son ensemble et dans le détail, nous allons maintenant nous arrêter sur son vocabulaire. Un texte ne se comprend vraiment que si l'on maîtrise les mots qui le composent : c'est le lexique qui porte le sens. Ici, deux univers s'opposent : le monde de la \cle{paix} (\all{der Frieden}) et celui du \cle{conflit} (\all{der Konflikt}). Entre les deux, des verbes d'action montrent le chemin : écouter, discuter, pardonner, se réconcilier. Apprenons ce vocabulaire méthodiquement, avec les articles, les pluriels et les familles de mots, comme l'exige le programme de Première.

\courssub{Observer d'abord : un court texte repère}

\begin{textebox}[Ein Streit im Klassenzimmer]
In der Klasse 1 A am Gymnasium in Yaoundé gibt es heute ein Problem: Junior und Sandrine haben einen Streit. Junior hat laut Musik gehört, und Sandrine konnte nicht lernen. Sie schreit: „Du hast keinen Respekt!“ Er antwortet: „Du bist intolerant!“ Die anderen Schüler haben Angst vor Gewalt. Aber die Lehrerin bleibt ruhig. Sie sagt: „Hört zu! Wir lösen diesen Konflikt gewaltfrei. Redet miteinander!“ Die beiden Schüler setzen sich und diskutieren. Junior versteht: Sandrine hat morgen eine Prüfung. Sandrine versteht: Junior wollte nicht stören. Am Ende verzeihen sie sich und versöhnen sich. Die Lehrerin lächelt: „Seht ihr? Der Dialog ist wichtig für den Frieden. Toleranz und Verständigung sind besser als Krieg und Gewalt.“ Nach der Stunde schreibt die Klasse an die Tafel: „Wir müssen Konflikte gewaltfrei lösen. Frieden beginnt bei uns!“
\end{textebox}

\textbf{Glossaire :} der Streit (-e) : la dispute ; laut : bruyamment ; schreien : crier ; die Angst (Ängste) : la peur ; die Gewalt (-) : la violence ; gewaltfrei : sans violence, non violent ; zuhören : écouter ; miteinander reden : parler ensemble ; stören : déranger ; sich verzeihen : se pardonner ; sich versöhnen : se réconcilier ; die Tafel (-n) : le tableau.

\textbf{Questions de repérage :} 1. Wer hat einen Streit? 2. Wie löst die Lehrerin den Konflikt? 3. Was schreibt die Klasse an die Tafel?

Observez : presque tous les mots importants du chapitre apparaissent dans ce petit texte. Nous allons maintenant les classer, les définir et les apprendre avec leurs articles et leurs pluriels.

\courssub{Le champ lexical de la paix}

\begin{definition}[Die Verständigung]
\all{Die Verständigung (-en)} désigne le fait de se comprendre, de s'entendre par le dialogue, de communiquer pour trouver un accord. Le verbe associé est \all{sich verständigen} (s'entendre, communiquer, se mettre d'accord) : \all{Die beiden Länder verständigen sich auf Deutsch.} « Les deux pays communiquent en allemand. » Attention : ce n'est pas seulement « comprendre » (\all{verstehen}), c'est parvenir à un accord mutuel.
\end{definition}

Voici les noms essentiels du champ de la paix. Apprenez chaque nom \textbf{avec son article et son pluriel}, comme une seule unité :

\begin{center}
\begin{tabularx}{\linewidth}{|l|l|X|}
\hline
\rowcolor{popblueL} Deutsch & Plural & Français \\
\hline
der Frieden & (ohne Plural) & la paix \\
\hline
die Versöhnung & die Versöhnungen & la réconciliation \\
\hline
der Dialog & die Dialoge & le dialogue \\
\hline
die Toleranz & (ohne Plural) & la tolérance \\
\hline
die Verständigung & die Verständigungen & l'entente, la communication \\
\hline
die Freundschaft & die Freundschaften & l'amitié \\
\hline
der Respekt & (ohne Plural) & le respect \\
\hline
die Solidarität & (ohne Plural) & la solidarité \\
\hline
\end{tabularx}
\end{center}

\begin{propriete}[La méthode des couleurs pour les articles]
Pour mémoriser le genre des noms allemands, utilisez un code couleur : \all{der} = \textbf{bleu} (masculin), \all{die} = \textbf{rouge} (féminin), \all{das} = \textbf{vert} (neutre). Écrivez vos listes de vocabulaire en couleur, ou visualisez mentalement chaque nom dans sa couleur. Le genre est \cle{arbitraire} : il ne se devine pas, il s'apprend avec le mot. Ne dites jamais « Frieden » seul, dites toujours \all{der Frieden}.
\end{propriete}

Remarquez une régularité utile : les noms en \all{-ung} (\all{die Versöhnung}, \all{die Verständigung}), en \all{-schaft} (\all{die Freundschaft}) et en \all{-tät} (\all{die Solidarität}) sont \textbf{toujours féminins} : \all{die}. C'est une aide précieuse pour l'examen.

\begin{exemplebox}[Le champ de la paix en phrases]
\all{Der Dialog ist wichtig für den Frieden.} « Le dialogue est important pour la paix. »

\all{Toleranz und Respekt sind die Basis der Freundschaft.} « La tolérance et le respect sont la base de l'amitié. »

\all{Die Solidarität zwischen den Schülern ist groß.} « La solidarité entre les élèves est grande. »

\all{Nach dem Streit haben sie sich versöhnt.} « Après la dispute, ils se sont réconciliés. »
\end{exemplebox}

\courssub{Le champ lexical du conflit}

\begin{center}
\begin{tabularx}{\linewidth}{|l|l|X|}
\hline
\rowcolor{popblueL} Deutsch & Plural & Français \\
\hline
der Krieg & die Kriege & la guerre \\
\hline
der Konflikt & die Konflikte & le conflit \\
\hline
die Gewalt & (ohne Plural) & la violence \\
\hline
der Streit & die Streite & la dispute, la querelle \\
\hline
das Problem & die Probleme & le problème \\
\hline
die Angst & die Ängste & la peur \\
\hline
\end{tabularx}
\end{center}

Notez le pluriel avec tréma (Umlaut) : \all{die Angst} $\rightarrow$ \all{die Ängste}. C'est un pluriel classique des noms féminins courts en allemand, comme \all{die Hand} $\rightarrow$ \all{die Hände}, \all{die Nacht} $\rightarrow$ \all{die Nächte}. Remarquez aussi que le pluriel de \all{der Streit} (\all{die Streite}) est rare ; dans la langue courante, on emploie souvent \all{die Streitereien} pour « les disputes, les querelles ».

\begin{attention}[Der Krieg ou der Streit ?]
Ne confondez pas \all{der Krieg} (la guerre, entre États ou groupes armés) et \all{der Streit} (la dispute, la querelle entre personnes). Un conflit scolaire, une dispute entre amis ou entre voisins, c'est \all{ein Streit}, jamais \all{ein Krieg}. Dire \all{Wir haben einen Krieg in der Klasse} pour une dispute de cour de récréation serait ridicule en allemand. Nuance intermédiaire : \all{der Konflikt} est le terme général et abstrait, utilisable à tous les niveaux (conflit familial, conflit social, conflit international).
\end{attention}

\begin{exemplebox}[Le champ du conflit en phrases]
\all{Der Krieg zerstört Städte und Familien.} « La guerre détruit les villes et les familles. »

\all{Gewalt ist nie eine Lösung.} « La violence n'est jamais une solution. »

\all{Die Angst vor Gewalt ist groß.} « La peur de la violence est grande. »

\all{Wir müssen Konflikte gewaltfrei lösen.} « Nous devons résoudre les conflits sans violence. »
\end{exemplebox}

\courssub{Carte mentale : deux pôles opposés}

\begin{popfigure}
\begin{tikzpicture}[mindmap, grow cyclic, every node/.style={concept}, concept color=popblue!40, root concept/.append style={font=\Large\bfseries}, level 1/.append style={level distance=3.2cm, sibling angle=60}, level 2/.append style={level distance=2.2cm, sibling angle=45}]
\node[root concept]{Frieden und Konflikt}
  child[concept color=popgreen!50]{ node {FRIEDEN}
    child { node {der Dialog} }
    child { node {die Toleranz} }
    child { node {die Versöhnung} }
  }
  child[concept color=poporange!50]{ node {KONFLIKT}
    child { node {der Krieg} }
    child { node {die Gewalt} }
    child { node {der Streit} }
  }
  child[concept color=popgold!60]{ node {LÖSUNG}
    child { node {miteinander reden} }
    child { node {zuhören} }
    child { node {verzeihen} }
  };
\end{tikzpicture}
\legende{Carte mentale du lexique : les deux pôles opposés (paix / conflit) et le pont qui les relie : la solution par le dialogue.}
\end{popfigure}

Cette carte mentale résume la logique du chapitre : face au conflit, il existe toujours un chemin — \all{die Lösung} (la solution) — qui passe par l'écoute, la parole et le pardon.

\courssub{Les verbes d'action : agir pour la paix}

\begin{center}
\begin{tabularx}{\linewidth}{|l|X|X|}
\hline
\rowcolor{popblueL} Verb & Français & Beispiel \\
\hline
lösen & résoudre & Wir lösen das Problem. \\
\hline
vermitteln & servir de médiateur & Der Lehrer vermittelt im Streit. \\
\hline
zuhören & écouter & Hör mir bitte zu! \\
\hline
diskutieren & discuter & Wir diskutieren ruhig. \\
\hline
respektieren & respecter & Respektiere die anderen! \\
\hline
helfen (+ Dat.) & aider & Ich helfe meinem Freund. \\
\hline
verzeihen / vergeben (+ Dat.) & pardonner & Sie verzeiht ihm. \\
\hline
sich versöhnen & se réconcilier & Sie versöhnen sich nach dem Streit. \\
\hline
\end{tabularx}
\end{center}

Points de conjugaison à retenir :

\begin{itemize}
\item \all{helfen} est un verbe fort : \all{ich helfe, du hilfst, er hilft} ; au Perfekt : \all{er hat geholfen}. Il se construit avec le \textbf{datif} : \all{Ich helfe dir.} « Je t'aide. »
\item \all{verzeihen} et \all{vergeben} se construisent aussi avec le datif : \all{Ich verzeihe dir.} « Je te pardonne. » Ce sont des verbes forts : \all{er verzeiht, er hat verziehen} ; \all{er vergibt, er hat vergeben}.
\item \all{sich versöhnen} est réfléchi : \all{ich versöhne mich, du versöhnst dich...} ; avec une personne : \all{sich versöhnen mit} : \all{Er versöhnt sich mit seinem Bruder.} « Il se réconcilie avec son frère. »
\item \all{zuhören} est un verbe à \cle{particule séparable} : \all{Ich höre dir zu.} « Je t'écoute. » Au Perfekt : \all{Ich habe dir zugehört.}
\end{itemize}

\begin{propriete}[Les expressions utiles à réutiliser à l'écrit et à l'oral]
\begin{itemize}
\item \all{Frieden schließen} : conclure la paix. \all{Die beiden Länder schließen Frieden.} « Les deux pays concluent la paix. »
\item \all{einen Konflikt lösen} : résoudre un conflit. \all{Wir lösen den Konflikt mit Worten.} « Nous résolvons le conflit par les mots. »
\item \all{miteinander reden} : parler ensemble. \all{Redet miteinander, nicht übereinander!} « Parlez ensemble, ne parlez pas les uns des autres ! »
\item \all{für den Frieden kämpfen} : lutter pour la paix. \all{Viele Menschen kämpfen für den Frieden.} « Beaucoup de gens luttent pour la paix. »
\item \all{gewaltfrei} : non violent, sans violence. \all{eine gewaltfreie Lösung} « une solution non violente ».
\end{itemize}
\end{propriete}

\courssub{Les familles de mots : un mot en cache plusieurs}

L'allemand construit des familles de mots très régulières. Connaître un nom, c'est souvent connaître l'adjectif et le verbe de la même famille :

\begin{center}
\begin{tabularx}{\linewidth}{|X|X|X|}
\hline
\rowcolor{popblueL} Nomen & Adjektiv / Verb & Beispiel \\
\hline
der Frieden & friedlich (Adj.) & eine friedliche Lösung (une solution pacifique) \\
\hline
die Toleranz & tolerant (Adj.) & Sei tolerant! (Sois tolérant !) \\
\hline
die Verständigung & sich verständigen (Verb) & Sie verständigen sich auf Deutsch. \\
\hline
der Krieg & Krieg führen (Ausdruck) & Die Länder führen Krieg. (Les pays font la guerre.) \\
\hline
die Gewalt & gewaltfrei / gewalttätig & gewaltfrei demonstrieren (manifester sans violence) \\
\hline
die Freundschaft & befreundet (Adj.) & Wir sind befreundet. (Nous sommes amis.) \\
\hline
\end{tabularx}
\end{center}

\begin{exemplebox}[De la famille à la phrase]
\all{Der Konflikt ist friedlich gelöst.} « Le conflit est résolu pacifiquement. »

\all{Ein toleranter Mensch respektiert andere Meinungen.} « Une personne tolérante respecte les opinions des autres. »

\all{Nach langem Krieg schließen die Länder endlich Frieden.} « Après une longue guerre, les pays concluent enfin la paix. »
\end{exemplebox}

\begin{attention}[Der Frieden : un nom faible !]
\all{Der Frieden} est un \cle{nom faible} : il prend la terminaison \all{-n} au datif et à l'accusatif singulier. On dit donc : \all{für den Frieden} (pour la paix), \all{den Frieden lieben} (aimer la paix), \all{mit dem Frieden} (avec la paix). À l'examen, n'écrivez jamais « der Friede » dans cet usage : la forme correcte du nominatif moderne est \all{der Frieden}. Déclinaison complète : nominatif \all{der Frieden}, accusatif \all{den Frieden}, datif \all{dem Frieden}, génitif \all{des Friedens}.
\end{attention}

\begin{savaistu}[Die Streitschlichter : les médiateurs des écoles allemandes]
Le mot \all{Streitschlichter} (littéralement « celui qui apaise les disputes », de \all{der Streit} + \all{schlichten}, apaiser) est typiquement allemand et intraduisible en un seul mot français : il désigne les \textbf{médiateurs scolaires}, des élèves formés pour aider leurs camarades à résoudre leurs conflits sans violence. De nombreuses écoles en Allemagne ont leurs \all{Streitschlichter} ; en Autriche et en Suisse, on parle plutôt de \all{Mediation} dans les écoles. Et au Cameroun ? Dans vos lycées, les délégués de classe et les conseillers d'orientation jouent souvent ce rôle de médiateur : \all{Sie vermitteln im Streit.}
\end{savaistu}

\courssub{Exercice résolu et entraînement}

\begin{exoresolu}[Complétez avec le bon mot du lexique]
Énoncé : complétez les phrases avec le mot qui convient : \all{der Dialog, die Gewalt, der Streit, die Versöhnung, tolerant, der Frieden}.

1. Nach dem langen \_\_\_\_\_\_\_\_ zwischen den zwei Freunden kommt endlich die \_\_\_\_\_\_\_\_.

2. \_\_\_\_\_\_\_\_ ist nie eine Lösung; wir müssen gewaltfrei handeln.

3. Alle Menschen in Kamerun wünschen sich \_\_\_\_\_\_\_\_ für das Land.

4. Sei \_\_\_\_\_\_\_\_ und respektiere die Meinung deines Nachbarn!
\tcblower
Solution détaillée :

1. \all{Nach dem langen \textbf{Streit} zwischen den zwei Freunden kommt endlich die \textbf{Versöhnung}.} « Après la longue dispute entre les deux amis arrive enfin la réconciliation. » Entre amis, on parle de \all{Streit} (dispute), pas de \all{Krieg} ; et la suite logique d'une dispute est la réconciliation.

2. \all{\textbf{Die Gewalt} ist nie eine Lösung; wir müssen gewaltfrei handeln.} « La violence n'est jamais une solution ; nous devons agir sans violence. » Le mot \all{gewaltfrei} dans la même phrase confirme le choix.

3. \all{Alle Menschen in Kamerun wünschen sich \textbf{Frieden} für das Land.} « Tous les habitants du Cameroun souhaitent la paix pour le pays. » Attention : \all{sich wünschen} est réfléchi et se construit avec l'accusatif ; le nom faible prend la forme \all{Frieden} — ici la forme est identique au nominatif, mais au masculin avec article on écrirait \all{den Frieden}.

4. \all{Sei \textbf{tolerant} und respektiere die Meinung deines Nachbarn!} « Sois tolérant et respecte l'opinion de ton voisin ! » Il fallait ici un adjectif de la famille de \all{die Toleranz}.
\end{exoresolu}

\begin{exercice}[À vous de jouer]
1. Donnez l'article et le pluriel de : \all{Konflikt, Angst, Freundschaft, Problem, Dialog}.

2. Traduisez en français : \all{Wir müssen Konflikte gewaltfrei lösen.} / \all{Nach dem Streit haben sie sich versöhnt.}

3. Traduisez en allemand : « Le dialogue est important pour la paix. » / « Écoute-moi, s'il te plaît ! » / « Elle aide son amie. »

4. Trouvez le mot de la famille : \all{der Frieden} (adjectif ?) ; \all{die Toleranz} (adjectif ?) ; \all{die Verständigung} (verbe ?).
\end{exercice}

\begin{corrige}[Exercice « À vous de jouer »]
1. \all{der Konflikt, die Konflikte} ; \all{die Angst, die Ängste} ; \all{die Freundschaft, die Freundschaften} ; \all{das Problem, die Probleme} ; \all{der Dialog, die Dialoge}.

2. « Nous devons résoudre les conflits sans violence. » / « Après la dispute, ils se sont réconciliés. »

3. \all{Der Dialog ist wichtig für den Frieden.} (attention au nom faible : \all{für den Frieden}) / \all{Hör mir bitte zu!} (verbe séparable : \all{zu} en fin de phrase) / \all{Sie hilft ihrer Freundin.} (datif : \all{ihrer Freundin}).

4. \all{der Frieden} $\rightarrow$ \all{friedlich} ; \all{die Toleranz} $\rightarrow$ \all{tolerant} ; \all{die Verständigung} $\rightarrow$ \all{sich verständigen}.
\end{corrige}

\begin{aretenir}[L'essentiel du lexique]
\begin{itemize}
\item Deux champs opposés : la paix (\all{der Frieden, der Dialog, die Toleranz, die Versöhnung, die Verständigung}) et le conflit (\all{der Krieg, der Konflikt, der Streit, die Gewalt, die Angst}).
\item Chaque nom s'apprend avec son \textbf{article} (méthode des couleurs : der = bleu, die = rouge, das = vert) et son \textbf{pluriel}.
\item \all{Der Frieden} est un nom faible : \all{für den Frieden}.
\item Une dispute entre personnes = \all{der Streit}, jamais \all{der Krieg}.
\item Les verbes clés : \all{lösen, zuhören (séparable), helfen (+ datif), verzeihen (+ datif), sich versöhnen (mit)}.
\item Les familles de mots multiplient votre vocabulaire : \all{der Frieden} $\rightarrow$ \all{friedlich} ; \all{die Toleranz} $\rightarrow$ \all{tolerant}.
\end{itemize}
\end{aretenir}

Ce vocabulaire est désormais votre boîte à outils. Dans la prochaine section, nous verrons comment l'activer avec la grammaire : l'impératif, pour donner des conseils et des consignes de paix (\all{Redet miteinander!}), puis le futur, pour imaginer le monde de demain.

\coursec{Grammaire 1 : l'impératif et le conseil}

Après avoir exploré le lexique de la paix et du conflit (\emph{Frieden, Konflikt, Versöhnung, Dialog, Toleranz}), nous disposons maintenant des mots pour agir. Mais comment formuler une recommandation, une consigne ou un appel à l'action en allemand ? C'est la fonction de \textbf{l'impératif} (\cle{der Imperativ}) et de ses alternatives modales. Observons d'abord comment le texte support \emph{Frieden beginnt bei uns} utilise cette forme pour nous interpeller.

\courssub{Observation à partir du texte}

Dans le texte lu précédemment, nous avons rencontré cette phrase forte :

\begin{textebox}[Extrait du texte « Frieden beginnt bei uns »]
„Redet miteinander statt übereinander!“
\end{textebox}

\all{Redet miteinander statt übereinander!} « Parlez les uns \emph{avec} les autres, et non les uns \emph{sur} les autres ! »

\begin{aretenir}[Analyse de la forme]
\begin{itemize}
    \item \textbf{Forme :} \all{Redet} $\rightarrow$ verbe \all{reden} (parler), 2\textsuperscript{e} personne du pluriel (\cle{ihr}), sans le pronom \all{ihr}.
    \item \textbf{Fonction :} c'est un \textbf{conseil}, une \textbf{exhortation} adressée à un groupe (la classe, les jeunes, les citoyens). L'auteur n'ordonne pas de manière autoritaire ; il invite à un changement de comportement pour instaurer la paix.
    \item \textbf{Structure :} \all{miteinander} (les uns avec les autres) s'oppose à \all{übereinander} (les uns sur les autres) ; \all{statt} signifie ici « au lieu de, plutôt que ».
\end{itemize}
\end{aretenir}

\courssub{Formation de l'impératif : les trois formes de l'adresse}

Contrairement au français qui possède trois formes d'impératif (\emph{parle / parlons / parlez}), l'allemand distingue \textbf{trois formes} selon la personne à qui l'on s'adresse (familier singulier, familier pluriel, politesse), mais \textbf{aucune forme en « nous »}. Cette distinction est \textbf{obligatoire}.

\begin{propriete}[Règle de formation de l'impératif]
\begin{center}
\begin{tabular}{|l|l|l|l|}
\hline
\rowcolor{popblueL}
\textbf{Personne} & \textbf{Règle de formation} & \textbf{\all{kommen}} & \textbf{\all{helfen}} \\ \hline
\cle{du} (tu) & Radical du présent (+ \textbf{-e} souvent omis à l'oral) & \all{Komm(e)!} & \all{Hilf!} (chgt \emph{e $\rightarrow$ i}) \\ \hline
\cle{ihr} (vous familier) & Forme du présent en \textbf{-t}, sans le pronom \all{ihr} & \all{Kommt!} & \all{Helft!} \\ \hline
\cle{Sie} (vous de politesse) & Infinitif + \textbf{Sie} (verbe puis pronom) & \all{Kommen Sie!} & \all{Helfen Sie!} \\ \hline
\end{tabular}
\end{center}

\textbf{Points clés :}
\begin{itemize}
    \item \textbf{Verbes à changement de voyelle (\emph{e $\rightarrow$ i/ie}) :} le changement se maintient à la forme \cle{du} : \all{lesen $\rightarrow$ Lies!}, \all{sprechen $\rightarrow$ Sprich!}, \all{geben $\rightarrow$ Gib!}, \all{helfen $\rightarrow$ Hilf!}.
    \item \textbf{Verbes en \emph{-eln} :} le \emph{e} du radical reste : \all{handeln $\rightarrow$ Handle!}, \all{lächeln $\rightarrow$ Lächle!}.
    \item \textbf{Verbes en \emph{a $\rightarrow$ ä} au présent :} \textbf{pas} de tréma à l'impératif \cle{du} : \all{fahren $\rightarrow$ Fahr!} (et non *\all{Fähr!}), \all{schlafen $\rightarrow$ Schlaf!}.
    \item \textbf{Verbe \all{sein} (être) :} totalement irrégulier : \all{Sei!} / \all{Seid!} / \all{Seien Sie!}.
    \item \textbf{Verbes à particule séparable :} la particule se détache et va \textbf{à la fin} : \all{zuhören $\rightarrow$ Hör zu! / Hört zu! / Hören Sie zu!}.
\end{itemize}
\end{propriete}

\begin{popfigure}
\begin{tikzpicture}[
    node distance=0.9cm and 0.6cm,
    box/.style={draw=popblue, thick, rounded corners, fill=popblueL, text width=3.2cm, align=center, font=\small},
    arrow/.style={-Stealth, thick, popdark},
    lab/.style={font=\bfseries\small, popdark}
]
\node[box, align=center] (inf){\all{reden}\\(infinitif)};
\node[box, below left=1.2cm and 1.6cm of inf, align=center] (du){\cle{du :} \all{Rede!}\\radical + \emph{e}};
\node[box, below=1.2cm of inf, align=center] (ihr){\cle{ihr :} \all{Redet!}\\présent sans \all{ihr}};
\node[box, below right=1.2cm and 1.6cm of inf, align=center] (sie){\cle{Sie :} \all{Reden Sie!}\\infinitif + \all{Sie}};
\draw[arrow] (inf) -- (du);
\draw[arrow] (inf) -- (ihr);
\draw[arrow] (inf) -- (sie);
\node[box, below=2.6cm of ihr, fill=poporangeL, draw=poporange, align=center] (sep){Verbe séparable \all{zuhören} :\\ \all{Hör zu!} / \all{Hört zu!} / \all{Hören Sie zu!}\\ la particule \all{zu} va à la fin};
\end{tikzpicture}
\legende{Schéma de formation : de l'infinitif aux trois formes de l'impératif, avec le cas des verbes à particule séparable.}
\end{popfigure}

\begin{center}
\begin{tabular}{|l|l|l|l|}
\hline
\rowcolor{popblueL}
\textbf{Verbe} & \textbf{\cle{du}} & \textbf{\cle{ihr}} & \textbf{\cle{Sie}} \\ \hline
\all{kommen} (venir) & \all{Komm(e)!} & \all{Kommt!} & \all{Kommen Sie!} \\ \hline
\all{helfen} (aider) & \all{Hilf!} & \all{Helft!} & \all{Helfen Sie!} \\ \hline
\all{lesen} (lire) & \all{Lies!} & \all{Lest!} & \all{Lesen Sie!} \\ \hline
\all{zuhören} (écouter) & \all{Hör zu!} & \all{Hört zu!} & \all{Hören Sie zu!} \\ \hline
\all{sein} (être) & \all{Sei!} & \all{Seid!} & \all{Seien Sie!} \\ \hline
\end{tabular}
\end{center}

\courssub{Emplois et nuances : consigne, conseil, invitation}

L'impératif ne sert pas qu'à commander. Dans le contexte de la « culture de la paix », il exprime surtout l'\textbf{invitation}, le \textbf{conseil bienveillant} ou l'\textbf{encouragement}.

\begin{exemplebox}[Emplois de l'impératif — exemples traduits]
\begin{itemize}
    \item \textbf{Consigne, instruction (affichage, règlement) :}\\
    \all{Seien Sie pünktlich!} « Soyez ponctuels ! » (forme polie)\\
    \all{Schreibt den Text ab!} « Recopiez le texte ! » (le professeur à la classe)
    \item \textbf{Conseil, encouragement (bienveillance) :}\\
    \all{Hilf deinem Freund!} « Aide ton ami ! » (\all{deinem Freund} : datif, car \all{helfen} + datif)\\
    \all{Vertraut einander!} « Faites-vous confiance mutuellement ! »
    \item \textbf{Invitation :}\\
    \all{Komm mit zum Friedenstreffen!} « Viens avec nous à la rencontre pour la paix ! »
    \item \textbf{Adoucissement avec \cle{bitte} et \cle{mal} (particules modales) :}\\
    \all{Hör mal zu!} « Écoute donc un peu ! » (\cle{mal} adoucit et familiarise)\\
    \all{Redet bitte miteinander!} « Parlez-vous, s'il vous plaît ! » (\cle{bitte} marque la politesse)\\
    \all{Sei nicht so streng!} « Ne sois pas si sévère ! »
\end{itemize}
\end{exemplebox}

\begin{attention}[Erreur fréquente n°1 : le \textbf{-st} fantôme du \cle{du}]
Les francophones gardent souvent la marque du présent de la 2\textsuperscript{e} personne du singulier (\all{-st}).
\begin{center}
\begin{tabular}{|l|l|}
\hline
\rowcolor{poporangeL}
\textbf{Faux} & \textbf{Correct} \\ \hline
*\all{Hörst zu!} & \all{Hör zu!} \\ \hline
*\all{Liest den Brief!} & \all{Lies den Brief!} \\ \hline
*\all{Kommst her!} & \all{Komm her!} \\ \hline
\end{tabular}
\end{center}
\cle{Règle d'or :} l'impératif \cle{du} \textbf{ne prend jamais de -st}. C'est le radical nu (+ \emph{e} facultatif).
\end{attention}

\begin{attention}[Erreur fréquente n°2 : la particule séparable oubliée]
Avec les verbes à particule séparable (\all{anrufen, zuhören, mitmachen, aufhören}), la particule \textbf{va à la fin de la phrase}, comme au présent.
\begin{center}
\begin{tabular}{|l|l|}
\hline
\rowcolor{poporangeL}
\textbf{Faux} & \textbf{Correct} \\ \hline
*\all{Zuhör!} & \all{Hör zu!} \\ \hline
*\all{Anruf mich!} & \all{Ruf mich an!} \\ \hline
*\all{Aufhör!} & \all{Hör auf!} \\ \hline
\end{tabular}
\end{center}
\end{attention}

\courssub{Le conseil avec les verbes modaux : une alternative nuancée}

L'impératif peut paraître direct, voire brutal, dans un dialogue personnel ou une lettre de conseil. L'allemand (comme le français) utilise alors les \textbf{verbes modaux} au \textbf{Konjunktiv II} (\all{sollten}, \all{könnten}) pour \textbf{nuancer} : suggestion, recommandation, possibilité.

\begin{propriete}[Conseil modal : \all{sollen} / \all{können} au Konjunktiv II + infinitif]
\begin{center}
\begin{tabular}{|l|l|l|}
\hline
\rowcolor{popblueL}
\textbf{Structure} & \textbf{Exemple allemand} & \textbf{Traduction française} \\ \hline
\cle{Du solltest} + Inf. & \all{Du solltest mehr zuhören.} & Tu devrais écouter davantage. \\ \hline
\cle{Man sollte} + Inf. & \all{Man sollte Konflikte gewaltfrei lösen.} & On devrait résoudre les conflits sans violence. \\ \hline
\cle{Du könntest} + Inf. & \all{Du könntest mit ihm reden.} & Tu pourrais parler avec lui. \\ \hline
\cle{Wir könnten} + Inf. & \all{Wir könnten einen Brief schreiben.} & Nous pourrions écrire une lettre. \\ \hline
\end{tabular}
\end{center}
\textbf{Nuance :} \all{sollten} = recommandation de bon sens ; \all{könnten} = proposition ouverte, possibilité ; \all{man sollte} = conseil général impersonnel (très fréquent dans les articles de presse et les éditoriaux). L'infinitif se place \textbf{à la fin} de la phrase.
\end{propriete}

\begin{methode}[Choisir : impératif ou modal ?]
\begin{enumerate}
    \item \textbf{Contexte public, slogan, affiche, règle, urgence} $\rightarrow$ \textbf{impératif} : \all{Seien Sie tolerant!}, \all{Redet miteinander!}
    \item \textbf{Dialogue entre amis, lettre de conseil, avis personnel, suggestion douce} $\rightarrow$ \textbf{modal} (\all{solltest} / \all{könntest}) : \all{Du solltest mit deinem Bruder reden.}
    \item \textbf{Conseil général, « on » impersonnel, rédaction formelle} $\rightarrow$ \textbf{\all{Man sollte} + infinitif} : \all{Man sollte Brücken bauen statt Mauern.}
    \item \textbf{Exhortation collective (« nous »)} $\rightarrow$ \textbf{\all{Lasst uns} + infinitif} : \all{Lasst uns Frieden schließen!} « Faisons la paix ! »
\end{enumerate}
\end{methode}

\courssub{Comparaison français / allemand : pronoms et formes}

\begin{center}
\begin{tabularx}{\linewidth}{|X|X|X|}
\hline
\rowcolor{popblueL}
\textbf{Aspect} & \textbf{Français} & \textbf{Allemand} \\ \hline
Pronom sujet à l'impératif & Absent (\emph{Parle !}, \emph{Parlez !}) & Absent pour \cle{du} / \cle{ihr} ; \textbf{présent et obligatoire} pour \cle{Sie} : \all{Kommen Sie!} \\ \hline
Nombre de formes & 3 (\emph{parle, parlons, parlez}) & 3 (\cle{du} / \cle{ihr} / \cle{Sie}) ; pas de forme en « nous » : on utilise \all{Lasst uns} + infinitif \\ \hline
Distinction de familiarité & \emph{tu} / \emph{vous} & \cle{du} / \cle{ihr} (familier) contre \cle{Sie} (poli) : choix social crucial \\ \hline
Verbe à particule & — & Particule en fin de phrase : \all{Hör zu!} \\ \hline
Adoucissement & \emph{s'il te plaît}, \emph{donc} & \cle{bitte}, \cle{mal}, \cle{doch} : \all{Komm mal!}, \all{Hilf mir bitte!} \\ \hline
Conseil nuancé & \emph{tu devrais}, \emph{on pourrait} & \all{Du solltest}, \all{Man könnte} : même logique modale \\ \hline
\end{tabularx}
\end{center}

\courssub{Exercice résolu : de l'intention à la forme correcte}

\begin{exoresolu}[Mettre au bon mode : impératif (\cle{du}/\cle{ihr}/\cle{Sie}) ou modal (\all{sollten}/\all{könnten})]
\textbf{Consigne :} transformez l'intention entre parenthèses en phrase allemande correcte, en choisissant la forme d'adresse appropriée au contexte.
\begin{enumerate}
    \item [Contexte : affiche dans un couloir d'école] (Règle : être tolérant !)
    \item [Contexte : d'un ami à un ami] (Conseil : tu devrais l'appeler.)
    \item [Contexte : le professeur à la classe] (Consigne : écoutez bien maintenant !)
    \item [Contexte : article de journal sur la culture de la paix] (Recommandation générale : on devrait construire des ponts.)
    \item [Contexte : une grande sœur à son petit frère] (Invitation : viens, nous faisons la paix !)
\end{enumerate}
\tcblower
\begin{enumerate}
    \item \all{Seien Sie tolerant!} — forme polie \cle{Sie} pour une affiche publique ; verbe \all{sein} irrégulier.
    \item \all{Du solltest ihn anrufen.} — modal \all{solltest} pour un conseil doux entre amis ; \all{anrufen} séparable, mais ici à l'infinitif en fin de phrase, la particule reste soudée.
    \item \all{Hört jetzt gut zu!} — impératif \cle{ihr} pour s'adresser au groupe ; \all{zuhören} séparable $\rightarrow$ \all{zu} à la fin.
    \item \all{Man sollte Brücken bauen.} — impersonnel \all{man} + \all{sollte} pour une recommandation générale.
    \item \all{Komm, wir schließen Frieden!} — impératif \cle{du} familier ; \all{wir schließen} exprime l'hortatif « nous ».
\end{enumerate}
\end{exoresolu}

\courssub{Entraînement : à vous de jouer !}

\begin{exercice}[Impératif et conseils pour la paix]
\begin{enumerate}
    \item Mettez les verbes entre parenthèses à l'impératif de la personne indiquée :
    \begin{itemize}
        \item \cle{du} : \all{\_\_\_\_\_\_\_\_ (lesen)} den Artikel über Gewaltfreiheit!
        \item \cle{ihr} : \all{\_\_\_\_\_\_\_\_ (helfen)} den Opfern von Mobbing!
        \item \cle{Sie} : \all{\_\_\_\_\_\_\_\_ (zugeben)} Ihre Fehler!
        \item \cle{du} : \all{\_\_\_\_\_\_\_\_ (sein)} mutig!
    \end{itemize}
    \item Choisissez la meilleure formulation (impératif ou modal) pour ces situations :
    \begin{itemize}
        \item Une pancarte à l'entrée d'une salle de dialogue : « \_\_\_\_\_\_\_\_ (Parlez calmement) ! »
        \item Un courriel d'un élève à son camarade en conflit : « \_\_\_\_\_\_\_\_ (Tu devrais t'excuser). »
        \item Un slogan pour la Journée de la paix : « \_\_\_\_\_\_\_\_ (Construisons la paix !) » — attention, il n'existe pas d'impératif en « nous » : utilisez \all{Lasst uns...} + infinitif.
    \end{itemize}
    \item Traduisez en allemand :
    \begin{itemize}
        \item Écoute-moi bien ! (familier, avec insistance)
        \item Aidez-vous les uns les autres ! (forme polie, affiche)
        \item On ne devrait pas juger trop vite. (conseil général)
    \end{itemize}
\end{enumerate}
\end{exercice}

\begin{corrige}[Exercice : Impératif et conseils pour la paix]
\begin{enumerate}
    \item \all{Lies} (changement \emph{e $\rightarrow$ ie}) ; \all{Helft} ; \all{Geben Sie Ihre Fehler zu!} (\all{zugeben} est séparable : \all{zu} à la fin) ; \all{Sei} (irrégulier).
    \item \all{Sprechen Sie ruhig!} ou \all{Redet ruhig!} (affiche $\rightarrow$ \cle{Sie} ou \cle{ihr} selon le public visé) ; \all{Du solltest dich entschuldigen.} (\all{sich entschuldigen} : verbe réfléchi, \all{dich} = accusatif) ; \all{Lasst uns Frieden bauen!} (hortatif \all{Lasst uns} + infinitif).
    \item \all{Hör mir mal gut zu!} ; \all{Helfen Sie einander!} ; \all{Man sollte nicht zu schnell urteilen.}
\end{enumerate}
\end{corrige}

\courssub{Texte d'application : « Ein Brief an die Schülersprecher »}

Le texte suivant mobilise le lexique de la paix et les structures grammaticales étudiées ici (impératif, modaux, particules séparables). Il sert de base à la compréhension détaillée et à la production guidée.

\begin{textebox}[Texte original : « Ein Brief an die Schülersprecher » — lettre ouverte d'une enseignante aux délégués des élèves]
\all{Liebe Schülersprecher,}

\all{der Frieden an unserer Schule beginnt nicht im Büro der Direktion, sondern auf dem Pausenhof, in den Klassenräumen und in den Chats. Deshalb appelliere ich an euch: Redet miteinander, nicht übereinander! Wenn ein Konflikt entsteht, hört zuerst zu, bevor ihr urteilt. Seid ehrlich, aber seid auch fair.}

\all{Manchmal ist es schwer, den ersten Schritt zu machen. Ihr könntet Vermittler sein. Man sollte Brücken bauen statt Mauern. Helft den Jüngeren, Streit gewaltfrei zu lösen. Seid Vorbilder für Toleranz und Respekt!}

\all{Ich bitte euch: Nehmt diese Aufgabe ernst. Kommt ins Gespräch, schreibt einen Plan, handelt jetzt. Der Frieden wartet nicht.}

\all{Mit freundlichen Grüßen}

\all{Frau Dr. Ngo Bitjek}
\end{textebox}

\begin{definition}[Glossaire du texte (mots nouveaux ou difficiles)]
\begin{itemize}
    \item \all{der Schülersprecher, die -} : le délégué des élèves
    \item \all{der Pausenhof, die -höfe} : la cour de récréation
    \item \all{appellieren an (+ Akk.)} : lancer un appel à
    \item \all{entstehen} (entstand, ist entstanden) : naître, éclater (conflit)
    \item \all{urteilen} : juger, porter un jugement
    \item \all{der Vermittler, die -} : le médiateur
    \item \all{die Mauer, die -n} : le mur
    \item \all{das Vorbild, die -er} : le modèle, l'exemple à suivre
    \item \all{die Aufgabe, die -n} : la tâche, la mission
    \item \all{ernst nehmen} : prendre au sérieux
    \item \all{handeln} : agir, passer à l'acte
    \item \all{gewaltfrei} : non violent, sans violence
\end{itemize}
\end{definition}

\begin{exercice}[Questions de compréhension et repérage grammatical sur le texte]
\begin{enumerate}
    \item Relevez \textbf{toutes les formes d'impératif} du texte et indiquez pour chacune la personne (\cle{du}, \cle{ihr} ou \cle{Sie}) et le verbe à l'infinitif.
    \item Relevez la \textbf{structure de conseil modal} (\all{könnten} / \all{sollten}) et traduisez-la.
    \item Identifiez \textbf{un verbe à particule séparable} employé à l'impératif et montrez où se place la particule.
    \item À qui s'adresse l'enseignante : à une seule personne ou à un groupe ? Justifiez avec la forme des impératifs.
    \item Traduisez en français le dernier paragraphe (\all{Ich bitte euch...}).
\end{enumerate}
\end{exercice}

\begin{corrige}[Exercice : Questions de compréhension et repérage grammatical]
\begin{enumerate}
    \item \all{Redet} (\cle{ihr}, \all{reden}) ; \all{hört ... zu} (\cle{ihr}, \all{zuhören}) ; \all{Seid} (\cle{ihr}, \all{sein}) ; \all{seid} (\cle{ihr}, \all{sein}) ; \all{Helft} (\cle{ihr}, \all{helfen}) ; \all{Seid} (\cle{ihr}, \all{sein}) ; \all{Nehmt} (\cle{ihr}, \all{nehmen}) ; \all{Kommt} (\cle{ihr}, \all{kommen}) ; \all{schreibt} (\cle{ihr}, \all{schreiben}) ; \all{handelt} (\cle{ihr}, \all{handeln}).
    \item \all{Ihr könntet Vermittler sein.} « Vous pourriez être médiateurs. » ; \all{Man sollte Brücken bauen statt Mauern.} « On devrait construire des ponts plutôt que des murs. »
    \item \all{hört zuerst zu} : la particule \all{zu} se détache du verbe \all{zuhören} et va en fin de phrase.
    \item À un groupe (les délégués) : tous les impératifs sont à la forme \cle{ihr} (\all{redet, hört, seid, helft, nehmt...}), et le pronom \all{euch} confirme le pluriel familier.
    \item « Je vous demande : prenez cette tâche au sérieux. Entrez dans le dialogue, écrivez un plan, agissez maintenant. La paix n'attend pas. »
\end{enumerate}
\end{corrige}

\begin{savaistu}[Culture : la « Streitschlichtung », la médiation par les pairs en Allemagne]
Dans de nombreuses écoles allemandes, il existe des programmes de \all{Streitschlichtung} (médiation des conflits) : des élèves formés, les \all{Streitschlichter} (médiateurs), aident leurs camarades à résoudre leurs disputes \textbf{sans violence} et, dans un premier temps, \textbf{sans intervention des adultes}. C'est une application concrète du principe : \all{„Redet miteinander statt übereinander!“} Au Cameroun aussi, des clubs de paix et des initiatives de médiation scolaire se développent dans les lycées, notamment à Yaoundé et à Douala, avec l'appui d'organisations de la société civile. La paix, cela s'apprend !
\end{savaistu}

\coursec{Grammaire 2 : le futur}

Dans la section précédente, nous avons appris à donner des conseils et des consignes grâce à l'impératif : \all{Sprich mit deinem Nachbarn !}, \all{Sei tolerant !} Ces formes servent à agir sur le présent immédiat. Mais notre texte \all{Frieden beginnt bei uns} ne se contente pas de conseils : il annonce aussi des projets, des promesses, des actions à venir. Pour cela, l'allemand dispose du \cle{futur}. C'est l'objet de cette nouvelle section : comment parler de l'avenir en allemand, et — surprise agréable — pourquoi les Allemands utilisent souvent... le présent pour parler du futur.

\courssub{Observation : que remarque-t-on dans le texte ?}

Reprenons la phrase clé de notre texte support :

\all{Nächstes Jahr werden wir ein Friedensfest organisieren.} « L'année prochaine, nous organiserons une fête de la paix. »

Observons attentivement sa construction :

\begin{itemize}
\item \all{werden} est conjugué (\all{werden wir}) et occupe la \textbf{deuxième position} de la phrase, comme tout verbe conjugué en allemand ;
\item \all{organisieren} est à l'\textbf{infinitif} et se trouve à la \textbf{toute dernière place} de la phrase ;
\item entre les deux, on trouve tous les compléments : \all{ein Friedensfest}.
\end{itemize}

C'est exactement la même structure en « cadre » (\all{Satzklammer}) que nous connaissons déjà avec les verbes à particule séparable et avec le \all{Perfekt} : le verbe conjugué ouvre le cadre en position 2, et l'élément non conjugué le ferme à la fin.

\begin{exemplebox}[Autres exemples observés]
\all{Wir werden Frieden schließen.} « Nous conclurons la paix. »

\all{Du wirst sehen, es wird besser.} « Tu verras, ça ira mieux. »

\all{Ich werde Medizin studieren.} « J'étudierai la médecine. »

\all{Die Schüler werden einen Dialog organisieren.} « Les élèves organiseront un dialogue. »
\end{exemplebox}

\courssub{Formation du futur I : werden + infinitif}

\begin{propriete}[Règle du futur I (Futur I)]
Le futur I se forme avec :
\begin{enumerate}
\item l'auxiliaire \cle{werden}, conjugué au présent, en \textbf{deuxième position} ;
\item le verbe principal à l'\textbf{infinitif}, à la \textbf{dernière position} de la phrase.
\end{enumerate}
\textbf{Dans une subordonnée} (introduite par \all{weil}, \all{dass}, \all{wenn}...), l'ordre change : l'infinitif reste avant l'auxiliaire, et \all{werden} conjugué passe à la toute fin :

\all{..., weil wir ein Fest organisieren werden.} « ..., parce que nous organiserons une fête. »

\all{Ich hoffe, dass du mich besuchen wirst.} « J'espère que tu me rendras visite. »
\end{propriete}

Voici la conjugaison complète de \all{werden} au présent, avec l'infinitif final :

\begin{center}
\begin{tabular}{|l|l|l|}\hline
\rowcolor{popblueL} \textbf{Personne} & \textbf{Futur I} & \textbf{Français} \\ \hline
ich & ich werde ... organisieren & j'organiserai \\ \hline
du & du wirst ... organisieren & tu organiseras \\ \hline
er/sie/es & er/sie/es wird ... organisieren & il/elle organisera \\ \hline
wir & wir werden ... organisieren & nous organiserons \\ \hline
ihr & ihr werdet ... organisieren & vous organiserez \\ \hline
sie/Sie & sie/Sie werden ... organisieren & ils/vous organiseront \\ \hline
\end{tabular}
\end{center}

\begin{attention}[Attention à la place de l'infinitif]
L'erreur la plus fréquente des francophones est de calquer l'ordre français (« je vais organiser une fête ») et de coller l'infinitif juste après \all{werden} :

✗ \all{Ich werde organisieren ein Fest.}

✓ \all{Ich werde ein Fest organisieren.}

En allemand, l'infinitif ferme toujours la phrase. Tout ce qui est complément (objet, lieu, temps) se place \emph{entre} \all{werden} et l'infinitif.
\end{attention}

\courssub{Les emplois du futur I}

Le futur I s'emploie dans quatre situations principales :

\begin{enumerate}
\item \textbf{La prévision} (\all{Vorhersage}) : on annonce ce qui va probablement arriver.

\all{Es wird morgen regnen.} « Il va pleuvoir demain. »

\all{Der Konflikt wird bald enden.} « Le conflit finira bientôt. »

\item \textbf{Le projet, l'intention} : on présente ce qu'on a décidé de faire.

\all{Ich werde nach dem Abitur Medizin studieren.} « Après le bac, j'étudierai la médecine. »

\all{Wir werden nächstes Jahr ein Friedensfest organisieren.} « Nous organiserons une fête de la paix l'année prochaine. »

\item \textbf{La promesse} : on s'engage fermement.

\all{Ich werde dir immer helfen.} « Je t'aiderai toujours. »

\all{Wir werden den Frieden schützen.} « Nous protégerons la paix. »

\item \textbf{L'hypothèse, la supposition} (souvent avec \all{wohl}, \all{sicher}, \all{vielleicht}) :

\all{Du wirst sehen, es wird besser.} « Tu verras, ça ira mieux. »
\end{enumerate}

\courssub{Le présent à valeur de futur : l'usage préféré des Allemands}

Voici un point capital pour les francophones. En français, on exige presque toujours le futur morphologique : « Demain, j'irai à Berlin. » En allemand, quand un \textbf{complément de temps} indique clairement l'avenir, on utilise très souvent le \textbf{présent} :

\all{Morgen fahre ich nach Berlin.} « Demain, je vais à Berlin. » (littéralement : « demain je pars »)

\all{Nächste Woche beginnt der Kurs.} « La semaine prochaine, le cours commencera. »

\all{Morgen diskutieren wir über Toleranz.} « Demain, nous discuterons de la tolérance. »

\begin{propriete}[Présent + marqueur temporel = futur]
Si la phrase contient un marqueur temporel d'avenir (\all{morgen}, \all{bald}, \all{nächstes Jahr}...), l'allemand emploie volontiers le \textbf{présent}. Le futur I n'est alors pas nécessaire ; il ajoute une nuance de promesse, de prévision forte ou d'insistance.
\end{propriete}

Les marqueurs temporels du futur à connaître :

\begin{center}
\begin{tabular}{|l|l|}\hline
\rowcolor{popgreenL} \textbf{Deutsch} & \textbf{Français} \\ \hline
morgen & demain \\ \hline
übermorgen & après-demain \\ \hline
bald & bientôt \\ \hline
später & plus tard \\ \hline
nächste Woche / nächstes Jahr & la semaine prochaine / l'année prochaine \\ \hline
in der Zukunft & à l'avenir \\ \hline
in zwei Tagen & dans deux jours \\ \hline
\end{tabular}
\end{center}

\begin{center}
\begin{tabularx}{\linewidth}{|X|X|X|}\hline
\rowcolor{popblueL} \textbf{Français} & \textbf{Deutsch} & \textbf{Remarque} \\ \hline
Demain, je viendrai. & Ich komme morgen. & présent + \all{morgen} : le plus naturel \\ \hline
Je t'aiderai, je te le promets. & Ich werde dir helfen. & futur I : promesse ferme \\ \hline
Il pleuvra sans doute. & Es wird wohl regnen. & futur I : prévision, supposition \\ \hline
L'année prochaine, nous voyagerons au Cameroun. & Nächstes Jahr reisen wir nach Kamerun. & présent + marqueur temporel \\ \hline
\end{tabularx}
\end{center}

\begin{savaistu}[Comment parlent vraiment les Allemands ?]
À l'oral, les Allemands réservent le futur I aux promesses solennelles et aux prévisions fortes. Pour un simple programme ou un rendez-vous, le présent suffit largement : \all{Ich komme morgen.} « Je viens demain. » Un Allemand qui dirait \all{Ich werde morgen kommen} insisterait : « je viendrai, compte sur moi ! » En traduction, choisis donc la forme selon le sens : programme simple → présent ; promesse ou prédiction → futur I.
\end{savaistu}

\courssub{La frise des temps : situer le futur I}

\begin{popfigure}
\begin{tikzpicture}[scale=1.0]
\draw[-{Stealth[length=3mm]}, thick, popdark] (0,0) -- (12.5,0);
\foreach \x/\name/\col in {2/{Vergangenheit\\(Perfekt)}/popgreen, 6/{Gegenwart\\(Präsens)}/popblue, 10/{Zukunft\\(Futur I)}/poporange} {
  \fill[\col] (\x,0) circle (0.12);
  \node[align=center, font=\small\bfseries, text=\col] at (\x,0.9) {\name};
}
\node[align=center, font=\footnotesize, text=popgreen] at (2,-0.9) {Wir haben ein Fest\\organisiert.};
\node[align=center, font=\footnotesize, text=popblue] at (6,-0.9) {Wir organisieren\\ein Fest.};
\node[align=center, font=\footnotesize, text=poporange] at (10,-0.9) {Wir werden ein Fest\\organisieren.};
\draw[-{Stealth[length=2mm]}, poppurple, thick] (6.6,-1.7) -- (9.4,-1.7);
\node[align=center, font=\footnotesize, text=poppurple] at (8,-2.2) {morgen, bald, nächstes Jahr};
\end{tikzpicture}
\legende{Frise des temps : le futur I projette l'action après le présent ; un marqueur temporel permet aussi d'utiliser le présent.}
\end{popfigure}

\courssub{Piège : werden, auxiliaire du futur ou verbe « devenir » ?}

\begin{attention}[Deux werden différents]
\all{Werden} est aussi un verbe lexical qui signifie \textbf{« devenir »}. Comment les distinguer ? Regarde s'il y a un \textbf{infinitif à la fin de la phrase} :

\all{Er wird Lehrer.} « Il devient professeur. » (pas d'infinitif → verbe « devenir », au présent)

\all{Er wird Lehrer werden.} « Il deviendra professeur. » (deux \all{werden} : auxiliaire du futur + infinitif \all{werden} !)

\all{Sie wird Ärztin.} « Elle devient médecin. » ≠ \all{Sie wird Medizin studieren.} « Elle étudiera la médecine. »
\end{attention}

\begin{attention}[Ne pas confondre werden (futur) et wollen (vouloir)]
Faux ami classique pour les francophones, car les deux se traduisent parfois par le futur en français :

\all{Ich werde helfen.} « J'aiderai. » (futur, simple constat ou promesse)

\all{Ich will helfen.} « Je veux aider. » (volonté, modal \all{wollen})

Dans « je veux réussir », n'écris jamais ✗ \all{Ich werde Erfolg haben wollen} pour dire « je veux » : dis simplement \all{Ich will Erfolg haben.}
\end{attention}

\begin{aretenir}[Complément Bac : le futur I de supposition]
Au baccalauréat, tu peux rencontrer le futur I avec \all{wohl} ou \all{sicher} pour exprimer une \textbf{supposition sur le présent} (et non une action future) :

\all{Er wird wohl krank sein.} « Il est probablement malade. » (littéralement : « il sera sans doute malade »)

\all{Sie wird jetzt sicher zu Hause sein.} « Elle est sûrement à la maison maintenant. »

Astuce de traduction : rends ce futur allemand par « probablement / sans doute + présent » en français.
\end{aretenir}

\courssub{Texte d'application : Unsere Pläne für den Frieden}

\begin{textebox}[Unsere Pläne für den Frieden]
Nächstes Jahr wird unsere Schule in Yaoundé ein großes Friedensfest organisieren. Wir werden viele Gäste einladen: Lehrer, Eltern und Schüler aus anderen Schulen. Zuerst werden wir einen Dialog zwischen den Klassen führen. Dann werden wir über Toleranz und Verständigung diskutieren. Ein Schüler sagt: „Ich werde ein Theaterstück über Versöhnung schreiben.“ Eine Schülerin ergänzt: „Und ich werde ein Lied für den Frieden singen.“ Unser Deutschlehrer verspricht: „Ich werde euch immer helfen.“ Am Nachmittag werden wir zusammen Fußball spielen, denn Sport verbindet die Menschen. Am Abend werden wir ein Friedensmahl vorbereiten. Wir werden Frieden schließen mit allen Konflikten in unserer Klasse. Morgen beginnt die Vorbereitung, und übermorgen werden wir die ersten Plakate malen. In der Zukunft werden wir jedes Jahr dieses Fest wiederholen. Frieden beginnt bei uns — und er wird bei uns bleiben!
\end{textebox}

\textbf{Glossaire} : der Gast, die Gäste (l'invité) — einladen (inviter) — führen (mener) — das Theaterstück, die Theaterstücke (la pièce de théâtre) — die Versöhnung (la réconciliation) — ergänzen (ajouter) — versprechen (promettre) — verbinden (relier, unir) — das Friedensmahl (le repas de la paix) — vorbereiten (préparer) — die Vorbereitung (la préparation) — das Plakat, die Plakate (l'affiche) — malen (peindre) — wiederholen (répéter, renouveler) — bleiben (rester).

\begin{enumerate}
\item \textbf{Questions de compréhension :}
\begin{enumerate}
\item Que fera l'école l'année prochaine ?
\item Qui les élèves inviteront-ils ?
\item Que promet le professeur d'allemand ?
\item Quand commence la préparation ?
\item Souligne dans le texte tous les verbes au futur I et repère la place de l'infinitif.
\end{enumerate}
\end{enumerate}

\courssub{Exercices}

\begin{exoresolu}[Mets les phrases au futur I]
Énoncé : mets les phrases suivantes au futur I. a) \all{Wir organisieren ein Fest.} b) \all{Ich helfe meinem Freund.} c) \all{Ihr diskutiert über Toleranz.}
\tcblower
Solution détaillée : on conjugue \all{werden} à la deuxième position et on place l'infinitif à la fin.

a) \all{Wir werden ein Fest organisieren.} (sujet \all{wir} → \all{werden} ; \all{organisieren} à la fin)

b) \all{Ich werde meinem Freund helfen.} (\all{ich werde} ; attention, \all{helfen} + datif : \all{meinem Freund})

c) \all{Ihr werdet über Toleranz diskutieren.} (\all{ihr werdet} ; le complément \all{über Toleranz} reste dans le cadre)
\end{exoresolu}

\begin{exoresolu}[Futur I ou présent ?]
Énoncé : choisis la forme la plus naturelle et justifie. a) « Demain, je vais à Douala. » b) « Je t'aiderai toujours, je te le promets ! »
\tcblower
Solution détaillée :

a) \all{Morgen fahre ich nach Douala.} — Il y a un marqueur temporel (\all{morgen}) et c'est un simple programme : le présent est la forme la plus naturelle en allemand.

b) \all{Ich werde dir immer helfen.} — C'est une promesse solennelle : le futur I s'impose. Note le datif \all{dir} après \all{helfen}.
\end{exoresolu}

\begin{exercice}[À toi de jouer]
\begin{enumerate}
\item Conjugue au futur I : \all{du (schließen) Frieden} ; \all{er (studieren) Medizin} ; \all{sie plural (malen) Plakate}.

\item Traduis en allemand : a) « L'année prochaine, nous organiserons un dialogue. » b) « Demain, nous discutons de la paix. » (utilise le présent) c) « Il deviendra professeur. »

\item Mets la phrase à la subordonnée avec \all{weil} : \all{Wir werden ein Fest organisieren.} → \all{Ich freue mich, weil ...}
\end{enumerate}
\end{exercice}

\begin{corrige}[Exercice « À toi de jouer »]
\begin{enumerate}
\item \all{Du wirst Frieden schließen.} — \all{Er wird Medizin studieren.} — \all{Sie werden Plakate malen.}

\item a) \all{Nächstes Jahr werden wir einen Dialog organisieren.} b) \all{Morgen diskutieren wir über den Frieden.} c) \all{Er wird Lehrer werden.} (deux \all{werden} : auxiliaire + infinitif « devenir »)

\item \all{Ich freue mich, weil wir ein Fest organisieren werden.} (dans la subordonnée, l'auxiliaire conjugué \all{werden} passe à la toute fin, après l'infinitif.)
\end{enumerate}
\end{corrige}

\begin{aretenir}[L'essentiel du futur]
\begin{itemize}
\item Futur I = \all{werden} conjugué (2\textsuperscript{e} position) + infinitif (dernière position) : \all{Ich werde ein Fest organisieren.}
\item Dans une subordonnée : \all{..., weil wir ein Fest organisieren werden.}
\item Emplois : prévision, projet, promesse, hypothèse.
\item Avec un marqueur temporel (\all{morgen}, \all{bald}, \all{nächstes Jahr}), l'allemand préfère souvent le présent : \all{Morgen fahre ich nach Berlin.}
\item \all{Er wird Lehrer.} = « il devient professeur » (pas d'infinitif) ≠ \all{Ich will helfen.} = « je veux aider » (\all{wollen}, pas le futur).
\item Niveau Bac : \all{Er wird wohl krank sein.} = « il est probablement malade » (supposition sur le présent).
\end{itemize}
\end{aretenir}

\coursec{Méthode du commentaire et production guidée}

Après avoir maîtrisé l'impératif pour formuler des conseils concrets et le futur I pour exprimer des engagements durables, nous passons maintenant à la compétence ultime du programme : \cle{l'exploitation autonome d'un texte} et \cle{la production écrite/orale structurée}. Cette section vous donne la méthode officielle pour réussir le commentaire de texte au Baccalauréat et vous guide pas à pas pour rédiger un paragraphe argumenté et mener un dialogue de médiation, en réinvestissant tout le lexique et la grammaire du chapitre.

\courssub{Méthode du commentaire de texte au Baccalauréat}

Le commentaire de texte n'est pas un résumé : c'est une \cle{analyse argumentée} qui montre que vous avez compris le sens global, la structure interne, le vocabulaire clé et que vous êtes capable de prendre position. Au Bac (épreuve écrite et orale), on attend cinq étapes rigoureuses.

\begin{methode}[Les 5 étapes du commentaire de texte (Bac)]
\textbf{Étape 1 — Présenter le texte} : Nature (article, rapport, dialogue, extrait de roman), thème général, source (titre, auteur, date, provenance), public visé. \textit{Formule type :} \all{Der Text „Frieden beginnt bei uns“ ist ein Zeitungsartikel / ein Bericht über… Es geht um…}
\textbf{Étape 2 — Dégager l'idée principale} : Quel est le message central en une phrase ? \textit{Formule :} \all{Die Hauptidee des Textes ist… / Der Autor möchte zeigen, dass…}
\textbf{Étape 3 — Analyser la structure et les arguments} : Découper le texte en parties logiques (introduction, développement, conclusion). Relever les connecteurs logiques (\all{zuerst, dann, außerdem, aber, deshalb, schließlich}). Citer 2--3 arguments précis avec des exemples du texte.
\textbf{Étape 4 — Relever le lexique thématique et les procédés} : Lister les mots-champs du thème (ici : \cle{Frieden, Konflikt, Gewalt, Dialog, Toleranz, Versöhnung}). Noter les figures de style (répétitions, métaphores, questions rhétoriques) et les temps de verbe dominants (Präsens pour les faits, Futur pour les perspectives, Impératif pour les appels à l'action).
\textbf{Étape 5 — Donner un avis personnel argumenté} : Exprimer votre opinion en la reliant au texte ET à votre contexte (Cameroun, votre école). \textit{Formules :} \all{Meiner Meinung nach… / Ich finde, dass… / Meines Erachtens ist… wichtig, weil…}
\end{methode}

\begin{popfigure}
\begin{tikzpicture}[
    node distance=1.2cm and 1.5cm,
    box/.style={draw=popblue, thick, rounded corners, fill=popblueL, text width=3.2cm, align=center, minimum height=1.8cm, font=\small},
    arrow/.style={-Stealth, thick, popdark}
]
\node[box, align=center] (p1){1. Präsentation\\Typ, Thema, Quelle};
\node[box, right=of p1, align=center] (p2){2. Hauptidee\\Kernaussage in einem Satz};
\node[box, right=of p2, align=center] (p3){3. Struktur \& Argumente\\Gliederung, Beweise, Konnektoren};
\node[box, right=of p3, align=center] (p4){4. Lexik \& Verfahren\\Wortfelder, Stilfiguren, Zeiten};
\node[box, right=of p4, align=center] (p5){5. Persönliche Meinung\\Bezug auf Text + eigener Kontext};
\draw[arrow] (p1) -- (p2);
\draw[arrow] (p2) -- (p3);
\draw[arrow] (p3) -- (p4);
\draw[arrow] (p4) -- (p5);
\end{tikzpicture}
\legende{Organigramme : les 5 étapes du commentaire de texte (flux logique obligatoire).}
\end{popfigure}

\begin{exemplebox}[Formules d'opinion personnelle (avec traduction)]
Pour l'étape 5, variez les introductions d'opinion. Le verbe est en \textbf{2e position} après l'adverbe de liaison.
\begin{center}
\begin{tabularx}{\linewidth}{|X|X|X|}\hline
\rowcolor{popblueL} \textbf{Allemand} & \textbf{Français} & \textbf{Structure} \\ \hline
\all{Meiner Meinung nach ist Dialog der Schlüssel.} & « Selon moi, le dialogue est la clé. » & Adv. + V2 + Sujet \\ \hline
\all{Ich finde, dass Toleranz sehr wichtig ist.} & « Je trouve que la tolérance est très importante. » & \all{Ich finde,} + \cle{dass}-phrase (verbe à la fin) \\ \hline
\all{Meines Erachtens sollten wir mehr zuhören.} & « À mon avis, nous devrions plus écouter. » & Génitif + \all{sollten} (Konj. II) + infinitif \\ \hline
\all{Ich bin der Ansicht, dass Gewalt nie löst.} & « Je suis d'avis que la violence ne résout rien. » & \all{der Ansicht sein, dass} \\ \hline
\all{Für mich beginnt Frieden im Klassenzimmer.} & « Pour moi, la paix commence dans la salle de classe. » & Préposition + pronom + V2 \\ \hline
\end{tabularx}
\end{center}
\end{exemplebox}

\begin{textebox}[Commentaire modèle — « Frieden beginnt bei uns »]
\all{Der Text „Frieden beginnt bei uns“ handelt von einem Mediationsprojekt an einer Schule in Stuttgart. Die Hauptidee ist: Man kann Konflikte ohne Gewalt lösen, wenn man miteinander redet. Der Text ist in drei Teile gegliedert: Erstens wird das Projekt vorgestellt (Streitschlichter-Ausbildung), zweitens werden konkrete Beispiele für gelöste Streitfälle genannt, drittens fordert der Autor die Ausweitung auf alle Schulen. Der Wortschatz gehört zum Themenfeld „Frieden und Konflikt“: \textbf{der Dialog, die Toleranz, die Verständigung, der Streit, die Gewalt, die Mediation}. Meiner Meinung nach ist dieses Projekt auch für Kamerun sehr interessant, denn in unseren Schulen in Yaoundé oder Douala gibt es ebenfalls oft Streit unter Schülern. Ich finde, dass wir auch Streitschlichter ausbilden sollten, weil Gewalt keine Lösung ist. Frieden beginnt wirklich bei jedem von uns.}
\end{textebox}

\begin{attention}[Pièges fréquents au commentaire]
\begin{itemize}
    \item \textbf{Confondre résumé et commentaire :} Ne racontez pas l'histoire ; analysez \textit{comment} le texte construit son sens.
    \item \textbf{Oublier les citations :} Chaque argument de l'étape 3 doit s'appuyer sur un mot ou une phrase du texte (\all{„Zitat“}, Zeile X).
    \item \textbf{Opinion trop vague :} « \all{Ich finde es gut} » ne suffit pas. Justifiez : « \all{Ich finde es gut, WEIL…} » (subordonnée \all{weil} + verbe à la fin).
    \item \textbf{Ordre des mots dans l'opinion :} Après \all{Meiner Meinung nach} $\rightarrow$ \textbf{verbe en 2e position} (\all{ist}, \all{gibt}, \all{sollen}). Après \all{Ich finde, dass} $\rightarrow$ \textbf{verbe À LA FIN}.
\end{itemize}
\end{attention}

\courssub{Plan type de production écrite : De la situation à l'avenir}

Toute production écrite argumentée (article, discours, lettre ouverte, essai) suit une architecture en trois temps qui mobilise les temps et modes vus dans ce chapitre.

\begin{center}
\begin{tabularx}{\linewidth}{|X|X|X|}\hline
\rowcolor{popblueL} \textbf{Partie} & \textbf{Fonction communicative} & \textbf{Outils grammaticaux obligatoires (Chap. 5)} \\ \hline
\textbf{Einleitung (Situation)} & Constater un problème, poser le décor, annoncer le sujet. & \textbf{Präsens} (faits actuels), \textbf{Präteritum} (récit passé si anecdote), connecteurs (\all{heutzutage, in unserer Schule, es gibt…}). \\ \hline
\textbf{Hauptteil (Vorschläge)} & Proposer des solutions concrètes, donner des conseils, exprimer l'obligation/la nécessité. & \textbf{Imperativ} (ordres/conseils directs), \textbf{Modale Verben} (\all{sollen, müssen, können, wollen}), \textbf{Konjunktiv II} (\all{würde, könnte, sollte}) pour la politesse/hypothèse. \\ \hline
\textbf{Schluss (Zukunft / Appell)} & Projeter les résultats, lancer un appel à l'action, ouvrir sur l'avenir. & \textbf{Futur I} (\all{werden} + infinitif) pour les prédictions/engagements, \textbf{Imperativ} (dernier appel), \all{lasst uns…}. \\ \hline
\end{tabularx}
\end{center}

\begin{popfigure}
\begin{tikzpicture}[
    node distance=1.0cm and 0.8cm,
    main/.style={draw=poporange, thick, rounded corners, fill=poporangeL, text width=2.8cm, align=center, minimum height=1.6cm, font=\small\bfseries},
    sub/.style={draw=popmuted, rounded corners, fill=white, text width=2.8cm, align=center, minimum height=1.2cm, font=\footnotesize},
    arrow/.style={-Stealth, thick, popdark}
]
\node[main, align=center] (sit){SITUATION\\Präsens / Fakten};
\node[sub, below=of sit, align=center] (sitd){„In unserer Schule\\gibt es oft Streit.“};
\node[main, right=of sit, align=center] (rat){RATSCHLÄGE\\Imperativ / sollen / müssen};
\node[sub, below=of rat, align=center] (ratd){„Hört zu!\\Wir sollen tolerant sein.“};
\node[main, right=of rat, align=center] (zuk){ZUKUNFT\\Futur I / werden + Inf.};
\node[sub, below=of zuk, align=center] (zukd){„Wir werden ein\\Projekt starten.“};
\node[main, right=of zuk, align=center] (faz){FAZIT / APPELL\\Imperativ / Lasst uns…};
\node[sub, below=of faz, align=center] (fazd){„Schaffen wir\\Frieden zusammen!“};
\draw[arrow] (sit) -- (rat);
\draw[arrow] (rat) -- (zuk);
\draw[arrow] (zuk) -- (faz);
\draw[arrow, dashed, popgreen] (sit.south) -- (sitd.north);
\draw[arrow, dashed, popgreen] (rat.south) -- (ratd.north);
\draw[arrow, dashed, popgreen] (zuk.south) -- (zukd.north);
\draw[arrow, dashed, popgreen] (faz.south) -- (fazd.north);
\end{tikzpicture}
\legende{Schéma du paragraphe de production : Situation $\rightarrow$ 2 Conseils (Impératif/\all{sollen}) $\rightarrow$ 2 Projets (Futur I) $\rightarrow$ Conclusion.}
\end{popfigure}

\courssub{Production écrite guidée : « Was können wir für den Frieden in unserer Schule tun? »}

\noindent\textbf{Consigne :} Rédigez un paragraphe cohérent de \textbf{6 à 8 phrases} en allemand. Vous devez impérativement inclure :
\begin{enumerate}
    \item Une phrase d'introduction (Situation, Präsens).
    \item \textbf{Deux impératifs} (2e personne pluriel \all{ihr} ou forme polie \all{Sie} — choisissez et restez cohérent).
    \item \textbf{Un conseil avec \all{sollen}} (ex. \all{Wir sollen…} ou \all{Man soll…}).
    \item \textbf{Deux phrases au Futur I} (\all{werden} + infinitif en fin de phrase).
    \item Une phrase de conclusion (appel ou vision d'avenir).
\end{enumerate}
Utilisez au moins 5 mots du lexique thématique (\cle{Frieden, Konflikt, Gewalt, Dialog, Toleranz, Versöhnung, Mediation, Streitschlichter}).

\begin{exoresolu}[Production modèle corrigée et annotée]
\textbf{Énoncé :} Rédiger le paragraphe demandé ci-dessus.

\textbf{Solution détaillée :}
\all{In unserer Schule in Yaoundé gibt es oft Konflikte unter den Schülern.} \textit{(Situation : Präsens, contexte camerounais naturel.)}\\
\all{Hört einander gut zu, bevor ihr urteilt!} \textit{(Impératif 1 : \all{hört} [hören], 2e pers. pl., verbe en 1re position, particule \all{zu} à la fin.)}\\
\all{Redet miteinander statt übereinander!} \textit{(Impératif 2 : \all{Redet} [reden], structure \all{statt} + Dativ.)}\\
\all{Wir sollen mehr Toleranz und Verständigung zeigen.} \textit{(Conseil \all{sollen} + infinitif \all{zeigen} en FIN de phrase.)}\\
\all{Wir werden ein Mediationsprojekt mit Streitschlichtern starten.} \textit{(Futur I 1 : \all{werden} 2e pos., \all{starten} FIN.)}\\
\all{Die Schüler werden friedlicher miteinander umgehen.} \textit{(Futur I 2 : \all{werden} 2e pos., \all{umgehen} infinitif à la fin. Note : \all{umgehen} est à particule séparable \rightarrow conjugué : \all{wir gehen um}, mais infinitif après \all{werden} : \all{umgehen}.)}\\
\all{So schaffen wir gemeinsam echten Frieden.} \textit{(Conclusion : Präsens performatif / vision d'avenir.)}

\textbf{Vérification des contraintes :} 7 phrases. 2 Impératifs (\all{Hört}, \all{Redet}). 1 \all{sollen}. 2 Futur I (\all{werden starten}, \all{werden umgehen}). Lexique : \all{Konflikte, Toleranz, Verständigung, Mediation, Streitschlichter, Frieden} (6 mots). \checkmark
\end{exoresolu}

\begin{exercice}[Votre tour : Production écrite]
Rédigez votre propre paragraphe sur le même sujet (« Was können wir für den Frieden in unserer Schule tun? ») en respectant scrupuleusement les 5 contraintes grammaticales. Écrivez sur votre copie, puis utilisez la grille d'auto-évaluation (section 5) pour vous corriger.
\end{exercice}

\begin{corrige}[Exercice Production écrite]
\textbf{Points de vigilance pour l'autocorrection :}
\begin{itemize}
    \item \textbf{Impératif :} Verbe en 1re position ? Terminaison \all{-t} pour \all{ihr} (\all{hört}, \all{redet}) ? Particule séparée à la fin (\all{aufhören} $\rightarrow$ \all{Hört auf!}) ?
    \item \textbf{\all{sollen} :} Conjugué à la 2e position (\all{wir sollen}), infinitif \textbf{à la fin} de la phrase (\all{... zeigen}).
    \item \textbf{Futur I :} \all{werden} conjugué (2e pos.) + \textbf{infinitif du verbe principal TOUT À LA FIN} (après les compléments). Ex. : \all{Wir werden nächste Woche ein Projekt starten.} (Pas \all{*Wir werden starten ein Projekt.})
    \item \textbf{Cohérence :} Le « wir » (nous) engage le groupe ; le « ihr » (vous) s'adresse aux autres. Choisis une perspective.
\end{itemize}
\end{corrige}

\begin{attention}[Relecture active : La chasse aux verbes !]
Avant de rendre votre copie, \textbf{soulignez tous les verbes conjugués} (au crayon de couleur).
\begin{enumerate}
    \item Comptez : Ai-je mes 2 impératifs (verbe position 1) ? Mon \all{sollen} (position 2) ? Mes 2 \all{werden} (position 2) ?
    \item Vérifiez \textbf{chaque infinitif final} : Est-il bien tout à la fin de la phrase (après \all{starten}, \all{zeigen}, \all{umgehen}) ?
    \item Vérifiez les \textbf{majuscules} sur \textbf{TOUS les noms} (\all{Schule, Frieden, Konflikt, Schüler, Projekt}).
    \item Vérifiez les \textbf{Umlaute} : \all{für, Schüler, hören, hört, grün, müde}.
\end{enumerate}
Cette relecture systématique évite 80\% des erreurs de niveau A2/B1.
\end{attention}

\courssub{Expression orale guidée : Jeu de rôle « Mediation » (Streitschlichtung)}

L'oral du Bac évalue l'interaction. Le jeu de rôle « Médiateur / Élève en conflit » est un classique du thème « Citoyenneté / Paix ». Il mobilise l'impératif (conseils), les modaux (\all{sollen, müssen, können}) et le futur (engagements).

\begin{experience}[Jeu de rôle : La médiation au lycée]
\textbf{Situation :} Deux élèves (A et B) se disputent pour une place de banc / un ballon / une remarque blessante. Un troisième élève (C) est le \cle{Streitschlichter} (médiateur pair) formé par l'école.
\textbf{Rôles :}
\begin{itemize}
    \item \textbf{Élève A (En colère) :} Exprime son ressentiment (\all{Ich bin sauer, weil… / Er hat mir…}).
    \item \textbf{Élève B (Sur la défensive) :} Donne sa version (\all{Das stimmt nicht! Ich habe nur…}).
    \item \textbf{Médiateur C (Neutre, calme) :} Structure le dialogue : \all{Ruhig bitte! / Hört einander zu! / Was ist genau passiert? / Wie fühlen Sie sich? / Was können wir tun, damit es nicht wieder vorkommt? / Ihr sollt euch die Hand geben. / Ihr werdet euch in Zukunft respektieren.}
\end{itemize}
\textbf{Durée :} 3--4 minutes par groupe. Changement de rôles.
\textbf{Critères :} Utilisation de l'impératif (min. 2), \all{sollen} (min. 1), Futur I (min. 1), lexique paix/conflit, fluidité, politesse (\all{Sie}-form ou \all{du} selon convention de classe).
\end{experience}

\begin{center}
\begin{tabularx}{\linewidth}{|X|X|}\hline
\rowcolor{popblueL} \textbf{Formules utiles pour le Médiateur (C)} & \textbf{Traduction / Contexte} \\ \hline
\all{Beruhigen Sie sich bitte. / Bleibt ruhig!} & Calmez-vous. (Formel polie / Impératif familier) \\ \hline
\all{Lasst uns das Problem besprechen.} & Discutons du problème. (\all{lasst uns} + infinitif) \\ \hline
\all{Erzählen Sie mir / Erzähl mir, was passiert ist.} & Racontez-moi ce qui s'est passé. \\ \hline
\all{Wie fühlen Sie sich dabei? / Wie fühlst du dich dabei?} & Comment vous sentez-vous ? / Comment te sens-tu ? \\ \hline
\all{Verstehen Sie den anderen? / Verstehst du ihn?} & Comprenez-vous l'autre ? \\ \hline
\all{Was schlagen Sie vor? / Was schlägst du vor?} & Que proposez-vous ? \\ \hline
\all{Wir sollen / Ihr sollt eine Lösung finden.} & Nous devons / Vous devez trouver une solution. \\ \hline
\all{Ich denke, wir werden uns vertragen.} & Je pense qu'on va se réconcilier. (Futur I) \\ \hline
\all{Shake hands! / Gebt euch die Hand!} & Serrez-vous la main ! (Impératif) \\ \hline
\end{tabularx}
\end{center}

\courssub{Grille d'auto-évaluation (Check-liste pour le Bac)}

Utilisez cette grille après chaque production (écrite ou orale) pour valider vos acquis du Chapitre 5.

\begin{center}
\begin{tabularx}{\linewidth}{|X|c|c|c|X|}\hline
\rowcolor{popblueL} \textbf{Critère de réussite} & \textbf{Oui} & \textbf{En partie} & \textbf{Non} & \textbf{Preuve / Exemple dans ma production} \\ \hline
\textbf{Vocabulaire thème} : 3 mots min. (\all{Frieden, Konflikt, Toleranz, Dialog, Gewalt, Versöhnung, Mediation, Streitschlichter}) & & & & \textit{Ex: …} \\ \hline
\textbf{Grammaire : Impératif} : Forme correcte (Pos 1, terminaison), 2 occurrences min. & & & & \textit{Ex: Hört zu!} \\ \hline
\textbf{Grammaire : Futur I} : \all{werden} (Pos 2) + Infinitif (Pos FIN), 2 occurrences min. & & & & \textit{Ex: Wir werden… starten.} \\ \hline
\textbf{Grammaire : Conseil \all{sollen}} : Conjugué Pos 2 + Infinitif Pos FIN. & & & & \textit{Ex: Wir sollen… zuhören.} \\ \hline
\textbf{Cohérence / Structure} : Situation $\rightarrow$ Conseils $\rightarrow$ Avenir (Plan 3 parties). & & & & \\ \hline
\textbf{Orthographe / Majuscules} : Tous les noms en majuscule, Umlauts (\all{ä, ö, ü}), \all{ß}. & & & & \\ \hline
\textbf{Ponctuation / Guillemets} : Guillemets allemands „ “ pour citations. & & & & \\ \hline
\end{tabularx}
\end{center}

\begin{savaistu}[Le Cameroun et la culture de la paix : un atout pour l'oral du Bac]
Le Cameroun célèbre chaque année la \textbf{Journée internationale de la paix} (21 septembre) et organise des \textbf{« Semaines de la paix »} dans les établissements scolaires (concours de dessin, théâtre-forum, clubs « Jeunes Artisans de Paix »). Le Ministère de l'Éducation de Base (MINEDUB) et le MINESEC intègrent l'\textbf{Éducation à la paix} dans les programmes d'EMV (Éducation à la Vie).
\textbf{Pourquoi c'est un atout pour vous ?} L'épreuve orale du Bac (allemand LV2) porte souvent sur des thèmes de société (citoyenneté, environnement, vivre-ensemble). Pouvoir décrire en allemand une action concrète vécue dans \textit{votre} lycée (ex. : \all{„In meiner Schule in Douala haben wir einen Friedensclub gegründet. Wir werden nächstes Jahr eine Mediationsausbildung machen.“}) impressionne le jury : cela prouve votre capacité à \textbf{transférer la langue dans votre réalité}, compétence clé du CECRL niveau A2/B1. Préparez 3--4 phrases « prêtes à l'emploi » sur vos activités paix au Cameroun !
\end{savaistu}

\newpage

\coursec{Activité d'intégration}

\coursec{La promotion de la culture de la paix}

\courssub{Objectifs de l'activité}

À l'issue de cette activité d'intégration, l'élève sera capable de :

\begin{itemize}
\item réinvestir le \cle{vocabulaire thématique} de la paix et du conflit (\all{der Frieden}, \all{der Krieg}, \all{die Versöhnung}, \all{der Dialog}, \all{die Toleranz}, \all{die Verständigung}) dans une production authentique ;
\item employer correctement l'\cle{impératif} (2\textsuperscript{e} personne du singulier et du pluriel, forme \all{wir}, forme de politesse \all{Sie}) pour formuler des slogans et des appels à l'action ;
\item exprimer le \cle{conseil} avec les verbes modaux \all{sollen} (notamment au subjonctif II \all{sollten}) et \all{können} au présent ;
\item utiliser le \cle{futur I} (\all{werden} + infinitif) pour annoncer des projets concrets ;
\item présenter oralement, en deux minutes, un document collectif devant la classe, avec une prononciation soignée et un contact visuel avec l'auditoire ;
\item travailler en groupe, répartir les tâches et évaluer la production des autres selon des critères précis.
\end{itemize}

\courssub{Situation}

Le lycée bilingue de Nkol-Eton, à Yaoundé, participe à la „Woche des Friedens“ (la Semaine de la paix), organisée chaque année dans plusieurs établissements du Cameroun en partenariat avec le club allemand. À cette occasion, la direction lance un concours : chaque classe de Première doit concevoir une affiche et un court plaidoyer pour promouvoir la culture de la paix dans la cour de récréation, à la cantine et sur les réseaux sociaux du lycée. Les meilleures affiches seront exposées au CDI et le groupe gagnant lira son plaidoyer lors de la cérémonie de clôture.

Voici le message officiel du proviseur, affiché au tableau d'information :

\begin{textebox}[Ankündigung der Schulleitung]
Liebe Schülerinnen und Schüler,

unser Gymnasium feiert vom 12. bis 16. Mai die „Woche des Friedens“. Konflikte gibt es überall: in der Familie, in der Klasse, auf dem Sportplatz. Aber Gewalt ist keine Lösung! Deshalb laden wir alle Klassen ein: Gestaltet ein Plakat mit dem Motto „Frieden beginnt bei uns!“ und präsentiert eure Ideen in zwei Minuten vor der Klasse. Euer Plakat braucht: einen Slogan, drei gute Ratschläge und zwei konkrete Projekte für unser Gymnasium. Die beste Gruppe spricht bei der Abschlussfeier vor der ganzen Schule.

Viel Erfolg und viel Fantasie!

Der Schulleiter
\end{textebox}

\begin{definition}[Glossaire de la situation]
\begin{itemize}
\item \all{die Ankündigung, -en} : l'annonce
\item \all{die Gewalt} (pas de pluriel) : la violence
\item \all{die Lösung, -en} : la solution
\item \all{einladen (lädt ein, lud ein, hat eingeladen)} : inviter (verbe à particule séparable)
\item \all{gestalten} : concevoir, créer
\item \all{das Plakat, -e} : l'affiche
\item \all{der Ratschlag, -schläge} : le conseil
\item \all{die Abschlussfeier, -n} : la cérémonie de clôture
\end{itemize}
\end{definition}

\courssub{Consignes}

Travail en groupes de quatre élèves. Durée totale : 1 h de préparation en classe, 1 h de présentations et d'évaluation.

\textbf{Tâche 1 — Compréhension et brainstorming (10 minutes).} Lisez l'annonce du proviseur. Soulignez dans le texte trois mots du champ lexical de la paix et notez les trois éléments obligatoires de l'affiche. Puis, en groupe, listez au tableau de brouillon cinq problèmes ou conflits concrets de la vie du lycée (disputes à la cantine, moqueries, triche, bagarres au terrain de sport, malentendus entre camarades de classes différentes).

\textbf{Tâche 2 — Le slogan à l'impératif (10 minutes).} Rédigez un slogan court et percutant à l'impératif, 2\textsuperscript{e} personne du pluriel (et, si possible, une deuxième formule à la 2\textsuperscript{e} personne du singulier). Attention aux verbes à particule séparable et aux verbes forts. Exemples de formes attendues : \all{Redet miteinander!}, \all{Respektiert euch!}, \all{Hört zu!}, \all{Sprich mit deinem Nachbarn!}

\textbf{Tâche 3 — Trois conseils avec les modaux (10 minutes).} Formulez trois conseils pour résoudre les conflits sans violence, avec \all{sollen} (de préférence \all{sollten} pour la politesse) ou \all{können} au présent. Variez les sujets (\all{wir}, \all{man}, \all{ihr}, \all{jeder Schüler}). Chaque conseil doit contenir au moins un mot du vocabulaire thématique de la leçon (\all{der Dialog}, \all{die Toleranz}, \all{die Verständigung}, \all{die Versöhnung}, \all{gewaltfrei}...).

\textbf{Tâche 4 — Deux projets au futur I (10 minutes).} Annoncez deux actions concrètes que votre groupe ou votre classe réalisera cette année, au futur I : \all{werden} conjugué + infinitif en fin de phrase. Exemples : créer un club de la paix, organiser un match de football inter-classes, inviter un médiateur, décorer la salle de classe.

\textbf{Tâche 5 — Réalisation de l'affiche (15 minutes).} Disposez sur une feuille : le titre \all{„Frieden beginnt bei uns!“}, le slogan en gros caractères, les trois conseils, les deux projets, et un dessin symbolique simple (colombe, mains, drapeaux du Cameroun et de l'Allemagne). Vérifiez l'orthographe : majuscule à tous les noms, Umlaute, conjugaison correcte.

\textbf{Tâche 6 — Présentation orale et vote (2 minutes par groupe).} Chaque membre du groupe prend la parole au moins une fois : un élève présente le slogan, un autre les conseils, un troisième les projets, un quatrième conclut. La classe écoute et remplit la grille d'évaluation, puis vote pour la meilleure affiche et le meilleur plaidoyer.

\begin{methode}[Grille d'évaluation de la présentation]
\begin{enumerate}
\item Vocabulaire du thème : au moins cinq mots de la leçon correctement employés (2 points).
\item Grammaire : impératif correct, modaux bien conjugués, futur I avec infinitif final (4 points).
\item Communication orale : voix claire, regard vers la classe, répartition équilibrée des rôles (2 points).
\item Qualité visuelle et originalité de l'affiche (2 points).
\end{enumerate}
\end{methode}

\courssub{Ressources à mobiliser}

\begin{propriete}[Rappel 1 : l'impératif]
\begin{center}
\begin{tabular}{|l|l|l|}
\hline
\rowcolor{popblueL} Personne & Formation & Exemple \\
\hline
\all{du} & radical (sans \all{-st} ; voyelle modifiée pour verbes forts \all{e $\rightarrow$ i/ie}) & \all{Sprich lauter!} \\
\hline
\all{ihr} & radical + \all{-t} & \all{Redet miteinander!} \\
\hline
\all{Sie} & infinitif + \all{Sie} & \all{Hören Sie zu!} \\
\hline
\all{wir} & infinitif + \all{wir} & \all{Lösen wir Konflikte gewaltfrei!} \\
\hline
\end{tabular}
\end{center}
Verbes forts en \all{e $\rightarrow$ i/ie} au singulier : \all{sprechen $\rightarrow$ Sprich!}, \all{lesen $\rightarrow$ Lies!}, \all{helfen $\rightarrow$ Hilf!}, \all{geben $\rightarrow$ Gib!}, \all{nehmen $\rightarrow$ Nimm!}. Verbe à particule séparable : la particule va à la fin : \all{Hört zu!} (de \all{zuhören}), \all{Macht mit!} (de \all{mitmachen}).
\end{propriete}

\begin{propriete}[Rappel 2 : le conseil et le futur I]
\textbf{Le conseil} : \all{sollen} / \all{können} conjugués en 2\textsuperscript{e} position + infinitif en fin de phrase. 
\all{Wir sollten Konflikte gewaltfrei lösen.} « Nous devrions résoudre les conflits sans violence. » (forme du subjonctif II \all{sollten}, plus polie et recommandée pour un conseil). 
\all{Man kann immer zuerst sprechen.} « On peut toujours parler d'abord. » (présent de \all{können}, possibilité).

\textbf{Le futur I} : \all{werden} conjugué en 2\textsuperscript{e} position + infinitif en fin de phrase. 
\all{Wir werden einen Friedensclub gründen.} « Nous créerons un club de la paix. »
\all{Ich werde morgen hingehen.} « J'irai demain. » (verbe à particule : \all{hingehen} reste groupé à l'infinitif).
\end{propriete}

\begin{definition}[Boîte à mots utile : Frieden und Konflikt]
\begin{itemize}
\item \all{der Frieden} (pas de pluriel) : la paix ; \all{der Krieg, -e} : la guerre
\item \all{der Konflikt, -e} : le conflit ; \all{der Streit, -e} : la dispute
\item \all{die Versöhnung, -en} : la réconciliation ; \all{sich versöhnen} (versöhnt sich, versöhnte sich, hat sich versöhnt) : se réconcilier
\item \all{der Dialog, -e} : le dialogue ; \all{die Verständigung, -en} : l'entente, la compréhension mutuelle
\item \all{die Toleranz} (pas de pluriel) : la tolérance ; \all{tolerant} : tolérant
\item \all{der Respekt} (pas de pluriel) : le respect ; \all{respektieren} : respecter
\item \all{gewaltfrei} : non violent ; \all{der Friedensclub, -s} : le club de la paix
\item \all{miteinander reden} : se parler ; \all{zuhören} (hört zu, hörte zu, hat zugehört) : écouter ; \all{verzeihen} (verzeiht, verzieh, hat verziehen) : pardonner
\item \all{die Schlichtung, -en} : la médiation ; \all{der Vermittler, -} : le médiateur
\end{itemize}
\end{definition}

\begin{attention}[Erreurs fréquentes à éviter]
\begin{itemize}
\item \textbf{Futur I :} infinitif oublié ou mal placé. Jamais \all{Wir werden gründen einen Club}, toujours \all{Wir werden einen Club gründen.} (infinitif en \textbf{dernière position}).
\item \textbf{Impératif \all{ihr} :} terminaison \all{-t} obligatoire. Jamais \all{Rede miteinander!} (forme \all{du}) en s'adressant à un groupe, mais \all{Redet miteinander!}
\item \textbf{Verbe à particule séparable :} particule oubliée. \all{Hört zu!} et non \all{Hört!} (verbe \all{zuhören}) ; \all{Macht mit!} et non \all{Macht!} (verbe \all{mitmachen}).
\item \textbf{Majuscules :} tous les noms prennent une majuscule. \all{der Frieden}, \all{die Toleranz}, \all{der Dialog} — jamais \all{der frieden}.
\item \textbf{Genre et nombre :} \all{der Konflikt} (masculin, pluriel \all{die Konflikte}) — le français « conflit » est aussi masculin, pas de piège de genre, mais attention au pluriel allemand en \all{-e}. \all{die Gewalt}, \all{die Toleranz}, \all{die Versöhnung} sont féminins et n'ont pas de pluriel usuel.
\item \textbf{Modaux :} pas de \all{zu} devant l'infinitif. \all{Wir sollen helfen}, non \all{Wir sollen zu helfen}.
\end{itemize}
\end{attention}

\courssub{Proposition de production et corrigé commenté}

\begin{exoresolu}[Affiche et plaidoyer modèles du groupe « Les Colombes de Nkol-Eton »]
Énoncé : rédiger l'affiche complète (slogan, trois conseils, deux projets) et le plaidoyer oral de deux minutes.

\tcblower

\textbf{L'affiche :}

„FRIEDEN BEGINNT BEI UNS!“

Slogan : \all{Redet miteinander! Respektiert euch! Hört zu!}

Drei Ratschläge :
\begin{enumerate}
\item \all{Wir sollten Konflikte immer gewaltfrei lösen.}
\item \all{Man kann bei einem Streit zuerst den Dialog suchen.}
\item \all{Jeder Schüler soll Toleranz und Respekt zeigen.}
\end{enumerate}

Zwei Projekte :
\begin{enumerate}
\item \all{Wir werden einen Friedensclub an unserem Gymnasium gründen.}
\item \all{Wir werden ein Fußballspiel zwischen allen Klassen organisieren.}
\end{enumerate}

\textbf{Le plaidoyer oral :}

\all{Liebe Mitschüler! Gewalt ist keine Lösung — weder in der Familie noch in der Schule. Frieden beginnt bei uns, hier in Yaoundé, in unserer Klasse. Deshalb sagen wir: Redet miteinander! Respektiert euch! Hört zu! Wir sollten Konflikte gewaltfrei lösen und immer zuerst den Dialog suchen. Man kann auch einen Freund oder einen Lehrer um Hilfe bitten. In diesem Jahr werden wir einen Friedensclub gründen und ein Fußballspiel für den Frieden organisieren. Die Verständigung zwischen den Kulturen ist wichtig: Kamerun und Deutschland arbeiten schon lange zusammen. Macht mit! Gemeinsam sind wir stark. Danke für eure Aufmerksamkeit!}

« Chers camarades ! La violence n'est pas une solution — ni dans la famille ni à l'école. La paix commence avec nous, ici à Yaoundé, dans notre classe. C'est pourquoi nous disons : parlez-vous ! Respectez-vous ! Écoutez ! Nous devrions résoudre les conflits sans violence et toujours chercher d'abord le dialogue. On peut aussi demander de l'aide à un ami ou à un professeur. Cette année, nous créerons un club de la paix et nous organiserons un match de football pour la paix. L'entente entre les cultures est importante : le Cameroun et l'Allemagne collaborent depuis longtemps. Rejoignez-nous ! Ensemble, nous sommes forts. Merci de votre attention ! »

\textbf{Commentaire des choix de langue :}
\begin{itemize}
\item Le slogan utilise trois impératifs pluriels corrects : \all{Redet} (radical + \all{-t}), \all{Respektiert} (verbe régulier), et \all{Hört zu} avec la particule séparable \all{zu} placée en fin de segment (verbe \all{zuhören}).
\item Les conseils varient les sujets (\all{wir}, \all{man}, \all{jeder Schüler}) et les modaux (\all{sollten} subjonctif II, \all{kann} présent, \all{soll} présent) ; l'infinitif est toujours en fin de phrase (\all{lösen}, \all{suchen}, \all{zeigen}).
\item Les deux projets emploient le futur I : \all{werden} conjugué en 2\textsuperscript{e} position (\all{werden}), infinitif (\all{gründen}, \all{organisieren}) en dernière position.
\item Le plaidoyer réinvestit le vocabulaire de la leçon : \all{die Gewalt}, \all{die Lösung}, \all{der Dialog}, \all{die Verständigung}, \all{der Friedensclub}, et ancre le message dans le contexte camerounais (Yaoundé, coopération Cameroun-Allemagne).
\item La conclusion \all{Macht mit!} est un impératif à particule séparable (\all{mitmachen}), formule typique des campagnes germanophones.
\item La conjonction \all{weder ... noch} (ni ... ni) structure la négation de façon élégante.
\end{itemize}
\end{exoresolu}

\courssub{Bilan de l'activité}

Cette activité d'intégration a permis de mobiliser simultanément les quatre savoir-faire de la leçon : la \cle{compréhension} d'un document authentique (l'annonce du proviseur), l'emploi du \cle{lexique} de la paix et du conflit, l'application des \cle{règles de grammaire} (impératif, conseil modal, futur I) et la \cle{production} écrite et orale dans une situation de communication réaliste.

Pour vérifier que les objectifs sont atteints, chaque élève doit pouvoir répondre « oui » aux questions suivantes :

\begin{itemize}
\item Sais-je former l'impératif aux quatre personnes (\all{du}, \all{ihr}, \all{Sie}, \all{wir}), y compris avec un verbe fort (\all{e $\rightarrow$ i/ie}) ou à particule séparable ?
\item Sais-je donner un conseil avec \all{sollen} (idéalement \all{sollten}) ou \all{können} en plaçant l'infinitif en fin de phrase ?
\item Sais-je annoncer un projet au futur I avec \all{werden} + infinitif (infinitif en position finale) ?
\item Ai-je utilisé au moins cinq mots du vocabulaire thématique dans mon affiche ?
\item Ai-je pris la parole devant la classe avec une voix claire, un regard sur l'auditoire et un texte grammaticalement correct ?
\end{itemize}

\begin{aretenir}[L'essentiel]
Pour plaider en faveur de la paix en allemand, trois outils suffisent : l'\cle{impératif} pour mobiliser (\all{Redet miteinander!}), les \cle{modaux} pour conseiller (\all{Wir sollten den Dialog suchen.}) et le \cle{futur I} pour s'engager (\all{Wir werden einen Friedensclub gründen.}). Avec le vocabulaire de la paix, ces structures permettent de passer du message personnel au message citoyen : \all{Frieden beginnt bei uns!}
\end{aretenir}

\begin{savaistu}[Le Cameroun et la culture de la paix]
Le Cameroun, pays de plus de 250 langues, fait de la cohabitation pacifique de ses communautés un principe national : « vivre ensemble » est une valeur enseignée à l'école. En Allemagne, de nombreuses écoles décernent chaque année un \all{Friedenspreis} (prix de la paix) à des élèves qui s'engagent contre la violence et pour la tolérance. Le mot allemand \all{Frieden} vient d'une racine germanique ancienne (\all{fridu}) qui signifie « protection, amitié, absence d'hostilité » : la paix, c'est d'abord une relation positive entre les personnes — exactement le message de votre affiche.
\end{savaistu}

\newpage

\coursec{Méthodes et exercices résolus}

\courssub{Méthodes et exercices résolus}

\begin{methode}[Comprendre un texte : lecture globale puis détaillée]
\textbf{Objectif :} Extraire le sens global, puis les informations précises d'un texte d'actualité sur la paix.
\begin{enumerate}
\item \textbf{Lecture globale (1\textsuperscript{re} lecture) :} Lisez le texte sans vous arrêter au vocabulaire inconnu. Repérez : le \cle{titre}, la \cle{source} (si indiquée), le \cle{thème général}, le \cle{ton} (engagé, informatif, narratif). Répondez mentalement : « De quoi parle ce texte ? »
\item \textbf{Lecture détaillée (2\textsuperscript{e} lecture) :} Soulignez les \cle{mots-clés} du thème (Frieden, Gewalt, Dialog, etc.). Identifiez les \cle{connecteurs logiques} (deshalb, aber, weil, zuerst, schließlich) qui structurent l'argumentation.
\item \textbf{Vocabulaire inconnu :} Déduisez le sens par le contexte (famille de mot, transparent, antonyme). Notez les mots indispensables à la compréhension pour les regarder après.
\item \textbf{Réponse aux questions :} Citez le texte (numéro de ligne) pour justifier. Reformulez \emph{avec vos mots} (allemand) sans copier-coller de longues phrases.
\item \textbf{Synthèse :} Résumez en 3-4 phrases : Sujet + Problème + Solutions/Message final.
\end{enumerate}
\end{methode}

\begin{exoresolu}[Compréhension du texte « Frieden beginnt bei uns »]
\textbf{Énoncé :} Lisez le texte ci-dessous et répondez aux questions en allemand (phrases complètes).

\begin{textebox}[Texte support : Frieden beginnt bei uns]
(1) Frieden ist nicht nur das Fehlen von Krieg. (2) Frieden beginnt im Alltag, in der Familie, in der Schule und auf der Straße. (3) In Kamerun wie in Deutschland erleben junge Menschen manchmal Gewalt oder Streit. (4) Doch es gibt Wege aus der Spirale der Gewalt. (5) Der erste Schritt ist der Dialog. (6) Wer zuhört und die Meinung des anderen respektiert, praktiziert Toleranz. (7) In Yaoundé organisieren Schüler Friedenstage: sie singen, tanzen und diskutieren über Versöhnung. (8) Auch in Berlin gibt es Projekte gegen Rassismus und für Verständigung. (9) Wichtig ist: Jeder kann etwas tun. (10) Keine Beleidigungen, keine Schläge, sondern ein offenes Wort. (11) Frieden ist eine Aufgabe für uns alle – heute und morgen.
\end{textebox}

\begin{enumerate}
\item Was ist Frieden laut dem Text nicht nur? (Zitat)
\item Nennen Sie vier Orte, an denen Frieden beginnt.
\item Was ist der „erste Schritt“ zum Frieden?
\item Was machen die Schüler in Yaoundé an Friedenstagen? (Drei Aktivitäten)
\item Erklären Sie mit eigenen Worten: „Frieden ist eine Aufgabe für uns alle“.
\end{enumerate}

\tcblower
\textbf{Solution détaillée :}

\textbf{Analyse du texte :} Texte argumentatif, ton engagé, structure : Définition (l.1) $\rightarrow$ Contexte (l.2-3) $\rightarrow$ Solution/Dialogue (l.4-6) $\rightarrow$ Exemples Cameroun/Allemagne (l.7-8) $\rightarrow$ Appel à l'action (l.9-11).

\textbf{Repérage vocabulaire clé (avec article et pluriel) :}
\cle{das Fehlen} (absence, pas de pluriel), \cle{der Alltag} (quotidien, Pl. \textit{-e}), \cle{die Spirale der Gewalt} (spirale de la violence, Pl. \textit{-n}), \cle{die Versöhnung} (réconciliation, Pl. \textit{-en}), \cle{der Rassismus} (racisme, pas de pluriel usuel), \cle{die Verständigung} (entente/communication, Pl. \textit{-en}), \cle{die Beleidigung} (insulte, Pl. \textit{-en}), \cle{die Aufgabe} (tâche/mission, Pl. \textit{-n}).

\textbf{Réponses :}
\begin{enumerate}
\item \all{Frieden ist nicht nur das Fehlen von Krieg.} (L. 1) — \emph{La paix n'est pas seulement l'absence de guerre.}
\item \all{Frieden beginnt im Alltag, in der Familie, in der Schule und auf der Straße.} (L. 2) — \emph{Dans le quotidien, la famille, l'école et la rue.}
\item \all{Der erste Schritt ist der Dialog.} (L. 5) — \emph{Le premier pas est le dialogue.}
\item \all{Sie singen, tanzen und diskutieren über Versöhnung.} (L. 7) — \emph{Ils chantent, dansent et discutent de la réconciliation.}
\item \all{Das bedeutet, dass jeder Mensch Verantwortung hat, aktiv zum Frieden beizutragen, durch gewaltfreies Verhalten und Respekt im Alltag.} — \emph{Cela signifie que chaque être humain a la responsabilité de contribuer activement à la paix, par un comportement non-violent et le respect au quotidien.}
\end{enumerate}
\textbf{Conclusion :} Le texte montre que la paix est une construction quotidienne (Alltag) reliant le local (Yaoundé) et l'international (Berlin).
\end{exoresolu}

\begin{methode}[Travailler le lexique thématique : familles de mots, articles, pluriels]
\textbf{Objectif :} Maîtriser le vocabulaire de la paix et du conflit pour l'expression écrite/orale.
\begin{enumerate}
\item \textbf{Tableau de dérivation :} Pour chaque nom central, construisez la famille : \cle{Verbe} $\rightarrow$ \cle{Nom (article + Pluriel)} $\rightarrow$ \cle{Adjectif} $\rightarrow$ \cle{Antonyme}. Ex. : \all{versöhnen} (v) $\rightarrow$ \cle{die Versöhnung} (-, -en) $\rightarrow$ \all{versöhnlich} (adj) $\leftrightarrow$ \all{unversöhnlich}.
\item \textbf{Collocations (mots « amis ») :} Apprenez les verbes support : \all{Frieden \textbf{stiften / schließen / bewahren / gefährden}} ; \all{einen Konflikt \textbf{lösen / schlichten / vermeiden / eskalieren lassen}} ; \all{Toleranz \textbf{zeigen / üben / fordern}}.
\item \textbf{Prépositions :} Notez la régence : \all{sich \textbf{versöhnen} \textbf{mit} jdm (Dat)} ; \all{\textbf{Verständnis} \textbf{für} jdn/etwas (Akk) haben} ; \all{sich \textbf{einsetzen} \textbf{für} + Akk (paix, tolérance)}.
\item \textbf{Entraînement :} Transformez des phrases : Nominalisation (Verbe $\rightarrow$ Nom) / Verbalisation (Nom $\rightarrow$ Verbe). Ex. \all{Wir tolerieren uns} $\rightarrow$ \all{Wir üben Toleranz}.
\item \textbf{Révision espacée :} Réutilisez 5 mots du thème dans la production écrite de la leçon (résumé, conseil, futur).
\end{enumerate}
\end{methode}

\begin{exoresolu}[Lexique : Familles de mots et collocations]
\textbf{Énoncé :} Complétez le tableau. Donnez l'article, le pluriel (nom), la forme adjectivale et un antonyme si possible. Puis complétez les phrases avec la collocation correcte (verbe + nom/préposition).

\begin{center}
\begin{tabularx}{\linewidth}{|X|X|X|X|X|}
\hline
\rowcolor{popblueL} \textbf{Verbe} & \textbf{Nom (Art. Pl.)} & \textbf{Adjectif} & \textbf{Antonyme} & \textbf{Collocation type} \\
\hline
versöhnen & die Versöhnung, -en & versöhnlich & unversöhnlich & sich \textbf{mit jdm (Dat) versöhnen} \\
\hline
tolerieren & die Toleranz, -en & tolerant & intolerant / engstirnig & \textbf{Toleranz zeigen / üben} \\
\hline
verstehen & das Verständnis (kein Pl.) & verständnisvoll & unverständnisvoll & \textbf{Verständnis für jdn/etwas (Akk) haben} \\
\hline
streiten & der Streit, -e & streitbar / streitsüchtig & friedlich & \textbf{sich streiten um + Akk} \\
\hline
gewaltfrei leben & der Frieden (kein Pl.) & friedlich & gewalttätig / kriegerisch & \textbf{Frieden stiften / bewahren} \\
\hline
\end{tabularx}
\end{center}

\textbf{Phrases à compléter (conjuguez le verbe) :}
1. Die Schüler \_\_\_\_\_\_\_\_\_\_ (sich versöhnen) nach dem Streit \_\_\_\_\_\_\_\_\_\_ (miteinander).
2. Wir müssen \_\_\_\_\_\_\_\_\_\_ (Frieden stiften) in unseren Vierteln.
3. Die Lehrerin \_\_\_\_\_\_\_\_\_\_ (Verständnis haben) \_\_\_\_\_\_\_\_\_\_ (für die Probleme) der Jugendlichen.
4. Man soll Konflikte nicht \_\_\_\_\_\_\_\_\_\_ (eskalieren lassen), sondern \_\_\_\_\_\_\_\_\_\_ (lösen).

\tcblower
\textbf{Solution détaillée :}

\textbf{Tableau :} Voir ci-dessus. \textbf{Attention :} \cle{der Frieden} n'a pas de pluriel usuel (singularetantum). \cle{das Verständnis} pas de pluriel. \cle{streiten} $\rightarrow$ \cle{der Streit} (nom fort : der Streit, des Streites/Streits, die Streite).

\textbf{Phrases complétées :}
\begin{enumerate}
\item \all{Die Schüler \textbf{versöhnen sich} nach dem Streit \textbf{miteinander}.} (Präsens, verbe pronominal \all{sich versöhnen} + \all{miteinander} / \all{mit einander}).
\item \all{Wir \textbf{müssen Frieden stiften} in unseren Vierteln.} (Modal \all{müssen} + infinitif \all{stiften} à la fin).
\item \all{Die Lehrerin \textbf{hat Verständnis für die Probleme} der Jugendlichen.} (Présens \all{hat} + Akkusatif \all{die Probleme} après \all{für}).
\item \all{Man soll Konflikte nicht \textbf{eskalieren lassen}, sondern \textbf{lösen}.} (Structure \all{nicht ..., sondern ...} ; \all{lassen} + infinitif = passif de substitution).
\end{enumerate}

\textbf{Justifications grammaticales :} Cas (Dat après \all{mit}, Akk après \all{für}), ordre des mots (verbe à la fin après modal), verbes pronominaux (\all{sich versöhnen}).
\end{exoresolu}

\begin{methode}[Employer l'impératif pour conseiller, exhorter, interdire]
\textbf{Objectif :} Former et utiliser l'impératif (2\textsuperscript{e} pers. sg/pl, 1\textsuperscript{e} pers. pl, forme de politesse) dans un contexte d'éducation à la paix.
\begin{enumerate}
\item \textbf{Identifier le destinataire :} \cle{Du} (familier sg) $\rightarrow$ \cle{Ihr} (familier pl) $\rightarrow$ \cle{Sie} (politesse sg/pl) $\rightarrow$ \cle{Wir} (nous, exhortation collective).
\item \textbf{Former l'impératif :}
\begin{itemize}
\item \textbf{Du :} Radicale du présent (sans \textit{-e} final pour verbes réguliers, garde \textit{-e} si radical en \textit{-d, -t, -m, -n} ou verbes forts voyelle modifiée \textit{a$\rightarrow$ä, e$\rightarrow$i/ie}). \textbf{Pas de pronom.} Ex: \all{komm!}, \all{arbeite!}, \all{lies!}, \all{sei!}.
\item \textbf{Ihr :} Présent \textit{ihr} sans pronom. Ex: \all{kommt!}, \all{arbeitet!}, \all{lest!}, \all{seid!}.
\item \textbf{Sie / Wir :} Infinitif + pronom \textbf{après} le verbe. Ex: \all{Kommen Sie!}, \all{Seien Sie!}, \all{Gehen wir!}.
\end{itemize}
\item \textbf{Verbes séparables :} Particule à la fin. \all{Hör zu! / Hört zu! / Hören Sie zu! / Hören wir zu!}
\item \textbf{Négation :} \all{Nicht} (ou \all{kein}) placé avant le complément / à la fin. \all{Mach das nicht!}, \all{Keine Gewalt!}
\item \textbf{Adoucir (Conseil) :} Particules modales \all{mal, doch, bloß, nur}. \all{Hör mal zu!}, \all{Sei doch nicht so laut!}
\item \textbf{Alternatives « conseil » (plus polies) :} \all{Man sollte + Inf.} / \all{Du könntest + Inf.} / \all{Wäre es nicht besser, ...?} / \all{Ich rate dir, ... zu + Inf.}
\end{enumerate}
\end{methode}

\begin{exoresolu}[Impératif : Conseils pour la paix / Transformation]
\textbf{Énoncé :}
\begin{enumerate}
\item Mettez aux trois formes de l'impératif (du, ihr, Sie) : \all{zuhören}, \all{verzeihen}, \all{helfen}, \all{sein}, \all{aufräumen}.
\item Transformez les phrases à l'indicatif en impératif (forme \textit{du}) + particule modale \textit{mal/doch} si naturel :
\begin{itemize}
\item Du hörst nicht zu. $\rightarrow$ \_\_\_\_\_\_\_\_\_\_
\item Du bist nicht tolerant. $\rightarrow$ \_\_\_\_\_\_\_\_\_\_
\item Du hilfst den Opfern nicht. $\rightarrow$ \_\_\_\_\_\_\_\_\_\_
\end{itemize}
\item Traduisez en allemand (forme \textit{Sie} - politesse / contexte officiel) :
\begin{itemize}
\item « Écoutez votre voisin ! »
\item « Ne faites pas de violence ! »
\item « Réconcilions-nous ! » (forme \textit{wir})
\end{itemize}
\end{enumerate}

\tcblower
\textbf{Solution détaillée :}

\textbf{1. Tableau de conjugaison (Impératif) :}

\begin{center}
\begin{tabular}{|l|l|l|l|l|}
\hline
\rowcolor{popblueL} \textbf{Infinitif} & \textbf{du} & \textbf{ihr} & \textbf{Sie} & \textbf{Wir} \\
\hline
zuhören (sep.) & hör zu! & hört zu! & hören Sie zu! & hören wir zu! \\
\hline
verzeihen (stark, ei$\rightarrow$ie) & verzeih! / verzeihe! & verzeiht! & verzeihen Sie! & verzeihen wir! \\
\hline
helfen (stark, e$\rightarrow$i) & hilf! / hilfe! & helft! & helfen Sie! & helfen wir! \\
\hline
sein (irrégulier) & sei! & seid! & seien Sie! & seien wir! \\
\hline
aufräumen (sep., ä) & räum auf! & räumt auf! & räumen Sie auf! & räumen wir auf! \\
\hline
\end{tabular}
\end{center}

\textit{Note :} Pour \textit{du}, la finale \textit{-e} est optionnelle pour \textit{verzeihen, helfen} (souvent omise à l'oral, conservée à l'écrit formel). Pour \textit{aufräumen}, radical \textit{räum-} (changement \textit{ä}).

\textbf{2. Transformation (du + mal/doch) :}
\begin{itemize}
\item \all{Hör \textbf{mal} nicht zu!} (interdiction) ou \all{Hör \textbf{doch} zu!} (conseil positif « Écoute donc ! »).
\item \all{Sei \textbf{mal} tolerant!} / \all{Sei \textbf{doch} tolerant!} (verbe \textit{sein} $\rightarrow$ \textit{sei}).
\item \all{Hilf \textbf{mal} den Opfern!} / \all{Hilf \textbf{doch} den Opfern!} (Datif \textit{den Opfern}).
\end{itemize}

\textbf{3. Traduction (Sie / Wir) :}
\begin{itemize}
\item \all{Hören Sie Ihrem Nachbarn zu!} (Datif après \all{zuhören}).
\item \all{Üben Sie keine Gewalt aus!} (Collocation \all{Gewalt ausüben} ; \all{keine} + Akkusatif). \textit{Variante courante :} \all{Machen Sie keine Gewalt!}
\item \all{Versöhnen wir uns!} (Verbe pronominal \all{sich versöhnen}, pronom \all{uns} après verbe).
\end{itemize}
\end{exoresolu}

\begin{methode}[Exprimer l'avenir et la volonté : le Futur I (werden + Infinitif)]
\textbf{Objectif :} Construire des phrases au Futur I pour prédire, promettre, exprimer une intention ferme ou une hypothèse.
\begin{enumerate}
\item \textbf{Structure :} \cle{Subjekt} + \cle{werden (conjugué Präsens)} + \dots + \cle{Infinitif (à la fin de la phrase)}.
\item \textbf{Conjugaison de \textit{werden} (irrégulier) :} \all{ich werde, du wirst, er/sie/es wird, wir werden, ihr werdet, sie/Sie werden}. \textbf{Attention :} \textit{du wirst} (pas *werdst), \textit{er wird} (pas *werdet).
\item \textbf{Emplois (Niveau A2/B1) :}
\begin{itemize}
\item \textbf{Prédiction objective / Futur certain :} \all{Morgen wird es regnen.} / \all{Die Konferenz wird nächste Woche stattfinden.}
\item \textbf{Promesse / Engagement fort :} \all{Ich werde dir helfen.} (Plus fort que \all{Ich helfe dir.} ou \all{Ich will helfen.})
\item \textbf{Supposition sur le présent (avec \textit{wohl, sicher, wahrscheinlich}) :} \all{Er wird wohl schon zu Hause sein.}
\end{itemize}
\item \textbf{Marqueurs temporels :} \all{morgen, übermorgen, nächste Woche, in Zukunft, bald, demnächst, 2025}.
\item \textbf{Verbes séparables / Pronominaux :} L'infinitif reste groupé (particule + radical / \textit{sich} + verbe) à la fin. \all{Wir werden uns versöhnen.} / \all{Er wird das Zimmer aufräumen.}
\item \textbf{Négation :} \all{nicht} / \all{kein} avant l'infinitif (ou avant le complément). \all{Wir werden keine Waffen kaufen.}
\end{enumerate}
\end{methode}

\begin{exoresolu}[Futur I : Conjugaison, ordre des mots, traduction]
\textbf{Énoncé :}
\begin{enumerate}
\item Conjuguez \textit{werden} au présent pour toutes les personnes. Écrivez la forme du Futur I pour \textit{ich} avec les verbes : \all{streiten, zuhören, versöhnen, friedlich leben}.
\item Placez les éléments entre parenthèses au Futur I (ordre des mots !) :
\begin{itemize}
\item Die Jugendlichen \_\_\_\_\_\_\_\_\_\_ (sich / für Frieden / einsetzen).
\item Wir \_\_\_\_\_\_\_\_\_\_ (nächstes Jahr / ein Friedensfest / organisieren).
\item Mein Bruder \_\_\_\_\_\_\_\_\_\_ (nicht / mit mir / streiten).
\end{itemize}
\item Traduisez en allemand (Futur I) :
\begin{itemize}
\item « Demain, nous ferons la paix. »
\item « Les politiciens ne résoudront pas le conflit sans dialogue. »
\item « Est-ce que tu m'aideras ? » (familier)
\end{itemize}
\end{enumerate}

\tcblower
\textbf{Solution détaillée :}

\textbf{1. Conjugaison \textit{werden} + Futur I :}
\all{ich werde, du wirst, er/sie/es wird, wir werden, ihr werdet, sie/Sie werden.}
\begin{itemize}
\item \all{ich werde streiten}
\item \all{ich werde zuhören} (sep : infinitif \all{zuhören})
\item \all{ich werde mich versöhnen} (pron : \all{sich} $\rightarrow$ \all{mich} pour \textit{ich})
\item \all{ich werde friedlich leben}
\end{itemize}

\textbf{2. Structure de phrase (Verbe conjugué 2\textsuperscript{e} pos, Infinitif fin) :}
\begin{itemize}
\item \all{Die Jugendlichen \textbf{werden sich für Frieden einsetzen}.} (Verbe pronominal séparable \all{sich einsetzen für + Akk} : \all{sich} + \all{für Frieden} + \all{einsetzen}).
\item \all{Wir \textbf{werden nächstes Jahr ein Friedensfest organisieren}.} (Complément temps \all{nächstes Jahr} $\rightarrow$ Accusatif \all{ein Friedensfest} $\rightarrow$ Infinitif).
\item \all{Mein Bruder \textbf{wird nicht mit mir streiten}.} (\all{nicht} avant \all{mit mir} / ou \all{nicht streiten} ; ici négation de l'action).
\end{itemize}

\textbf{3. Traduction :}
\begin{itemize}
\item \all{Morgen \textbf{werden wir Frieden schließen / stiften}.} (\all{Frieden schließen} collocation fixe). \textit{Ordre : Adv de temps $\rightarrow$ werden $\rightarrow$ Sujet $\rightarrow$ Infinitif.}
\item \all{Die Politiker \textbf{werden den Konflikt ohne Dialog nicht lösen}.} (\all{ohne Dialog} $\rightarrow$ préposition \all{ohne} + Akk ; \all{nicht} avant \all{lösen}).
\item \all{\textbf{Wirst} du mir \textbf{helfen}?} (Question fermée : Verbe 1\textsuperscript{re} pos. \all{du} familier $\rightarrow$ \all{wirst}).
\end{itemize}

\textbf{Pièges évités :} \all{werden} vs \all{wollen} (volonté), place de l'infinitif (toujours tout à la fin), accord \textit{werden} (pas de \textit{-t} à \textit{er wird}), pronom réfléchi au cas accusatif (\all{mich, dich, sich}) pour \textit{ich, du, er/sie}.
\end{exoresolu}

\begin{methode}[Rédiger un résumé et une production courte guidée (Commentaire)]
\textbf{Objectif :} Produire un texte cohérent (80-100 mots) structurant : Introduction $\rightarrow$ Arguments/Exemples $\rightarrow$ Conclusion/Engagement personnel.
\begin{enumerate}
\item \textbf{Préparation (Brainstorming) :} Mots-clés du texte + idées personnelles. Groupez en : \cle{Problème} (Gewalt, Streit, Intoleranz) / \cle{Solutions} (Dialog, Toleranz, Versöhnung, Projekte) / \cle{Appel} (Jeder, Wir, Jetzt).
\item \textbf{Structure type (Plan) :}
\begin{itemize}
\item \textbf{Einleitung (1 phr.) :} Accroche + Sujet + Thèse. \textit{Ex: „Frieden ist mehr als nur kein Krieg.“}
\item \textbf{Hauptteil (3-4 phr.) :} Arguments + Exemples (Cameroun / Allemagne). Connecteurs : \all{Zuerst / Erstens}, \all{Außerdem}, \all{Ein Beispiel ist...}, \all{Dazu kommt}, \all{Deshalb}.
\item \textbf{Schluss (1-2 phr.) :} Synthèse + Engagement personnel (Ich / Wir + Futur I ou Impératif/Modal). \textit{Ex: „Ich werde mich für Toleranz einsetzen.“}
\end{itemize}
\item \textbf{Outils linguistiques obligatoires :}
\begin{itemize}
\item Au moins 3 mots du lexique thème (avec articles corrects).
\item Au moins 1 \textbf{Impératif} (conseil) ou \textbf{Modalverb} (\all{sollen, müssen, können}).
\item Au moins 1 \textbf{Futur I} (prédiction/engagement).
\item Connecteurs logiques (\all{weil, dass, obwohl, wenn, deshalb}).
\end{itemize}
\item \textbf{Rédaction :} Écrivez au brouillon. Vérifiez : Majuscules (Noms), Virgules (subordonnées), Ordre des mots (V2 / Verbe fin), Cas (Prépositions), Orthographe (Umlaute, ß).
\item \textbf{Relecture :} Lis à voix haute. Est-ce fluide ? Le message est-il clair ?
\end{enumerate}
\end{methode}

\begin{exoresolu}[Production écrite : « Mein Beitrag zum Frieden »]
\textbf{Énoncé :} Rédigez un court paragraphe (environ 80-100 mots) en allemand sur le thème : \textbf{„Was kann ich für den Frieden tun?“} (Que puis-je faire pour la paix ?). Utilisez le vocabulaire de la leçon, l'impératif (conseil), le futur I (engagement) et des connecteurs.

\tcblower
\textbf{Solution détaillée (Modèle de rédaction corrigé) :}

\textbf{Texte proposé :}
„Frieden beginnt bei mir. \textbf{Erstens} höre ich meinen Freunden zu, wenn sie Probleme haben. \textbf{Man soll} nicht sofort streiten, sondern \textbf{reden}. \textbf{Außerdem} engagiere ich mich in der Schul-AG „Friedenstag“. \textbf{Dort} organisieren wir Workshops für Toleranz und Versöhnung. \textbf{In Zukunft werde ich} keine Beleidigungen mehr benutzen und \textbf{werde ich} Konflikte friedlich lösen. \textbf{Sei auch du tolerant!} Gemeinsam \textbf{können wir} die Welt besser machen.“

\textbf{Analyse grammaticale (Check-list Bac) :}
\begin{itemize}
\item \textbf{Vocabulaire :} \cle{der Frieden}, \cle{das Problem}, \cle{der Streit}, \cle{die Toleranz}, \cle{die Versöhnung}, \cle{die Beleidigung}, \cle{der Konflikt}. Articles/Pluriels corrects.
\item \textbf{Impératif :} \all{Sei auch du tolerant!} (Forme \textit{du} de \textit{sein} + particule \textit{auch}).
\item \textbf{Futur I :} \all{In Zukunft werde ich ... benutzen / werde ich ... lösen.} (Inversion après complément de temps \textit{In Zukunft} $\rightarrow$ \textit{werde ich} ; 2\textsuperscript{e} phrase coordonnée \textit{und werde ich}).
\item \textbf{Modaux :} \all{Man soll nicht streiten} (Impersonnel, conseil général). \all{Wir können ... machen}.
\item \textbf{Connecteurs/Structure :} \all{Erstens} (énumération), \all{sondern} (correction après \textit{nicht}), \all{Außerdem} (addition), \all{Dort} (référence lieu), \all{In Zukunft} (temps), \all{Gemeinsam} (adverbe en tête $\rightarrow$ inversion \textit{können wir}).
\item \textbf{Cas/Prépositions :} \all{für Frieden} (Akk), \all{für Toleranz} (Akk), \all{mit Freunden} (Dat), \all{in der Schul-AG} (Dat), \all{keine Beleidigungen} (Akk négatif).
\item \textbf{Ordre des mots :} Verbe 2\textsuperscript{e} position respecté dans toutes les principales. Verbe à la fin dans subordonnées (ici \textit{wenn}-phrase : \all{... wenn sie Probleme haben}).
\end{itemize}
\textbf{Note :} Ce texte obtient la note maximale car il répond au sujet, structure logique, utilise les points de grammaire imposés sans erreur, vocabulaire riche et orthographe soignée.
\end{exoresolu}

\begin{attention}[Les 5 erreurs qui coûtent des points au Bac]
\begin{enumerate}
\item \textbf{Majuscules oubliées sur les noms :} \all{der frieden, die toleranz, das verständnis} $\rightarrow$ \textbf{TOUS les noms en MAJUSCULE} : \all{der Frieden, die Toleranz, das Verständnis}. (Règle d'or allemande).
\item \textbf{Confusion \textit{werden} (Futur) / \textit{wollen} (Vouloir) / \textit{werden} (Passif) :} \all{Ich will Frieden machen} (volonté) $\neq$ \all{Ich werde Frieden machen} (promesse/prédiction) $\neq$ \all{Der Frieden wird gemacht} (passif). \textit{Conjugaison : ich werde, du \textbf{wirst}, er \textbf{wird}} (pas *werdet).
\item \textbf{Impératif \textit{du} : forme erronée :} \all{*Du hör zu!} (pronom interdit), \all{*liese!} (verbe fort \textit{lesen} $\rightarrow$ \textit{lies!}), \all{*arbeite!} (radicale en \textit{-t} $\rightarrow$ garde \textit{-e} : \textit{arbeite!}), \all{*vergiss nicht!} (correct), mais \all{*vergess nicht!} (faute). \textbf{Séparable :} \all{*Hör zu nicht!} $\rightarrow$ \all{Hör nicht zu!}
\item \textbf{Calque français « Faire la paix » $\rightarrow$ *\textit{machen Frieden} :} \textbf{Collocation correcte :} \all{Frieden \textbf{schließen}} ou \all{Frieden \textbf{stiften}}. \textit{Versöhnen} = se réconcilier (verbe pronominal \all{sich versöhnen mit + Dat}), pas *\textit{versöhnen} transitif direct.
\item \textbf{Cas après prépositions (Pièges « Paix ») :} \all{sich versöhnen \textbf{mit} + \textbf{Datif}} (pas *\textit{mit den} Ami $\rightarrow$ \textit{mit dem} Freund). \all{Verständnis \textbf{für} + \textbf{Akkusativ}} (pas *\textit{für dem} Problem $\rightarrow$ \textit{für das} Problem). \all{sich einsetzen \textbf{für} + \textbf{Akkusativ}}.
\end{enumerate}
\textbf{Astuce de relecture :} Surlignez \textbf{tous les noms} (majuscules ?), \textbf{tous les verbes conjugués} (2\textsuperscript{e} place ? accord ?), \textbf{toutes les prépositions} (cas correct ?).
\end{attention}

\newpage

\coursec{Exercices}

\courssub{Je vérifie mes connaissances}

\begin{exercice}[QCM : vocabulaire de la paix]
Entoure la bonne réponse.
\begin{enumerate}
\item \all{die Versöhnung} signifie : a) la guerre \quad b) la réconciliation \quad c) la peur
\item Le pluriel de \all{der Krieg} est : a) die Kriegen \quad b) die Kriege \quad c) die Krieger
\item \all{die Verständigung} signifie : a) l'entente, la compréhension mutuelle \quad b) la punition \quad c) la frontière
\item Quel verbe exprime le pardon ? a) kämpfen \quad b) verzeihen \quad c) zerstören
\end{enumerate}
\end{exercice}

\begin{exercice}[Vrai ou faux ? Justifie ta réponse]
Réponds par \all{richtig} ou \all{falsch} et justifie en français ou en allemand.
\begin{enumerate}
\item \all{Frieden beginnt bei uns.} signifie : « La paix commence chez nous / avec nous. »
\item À l'impératif avec \all{Sie}, le verbe se place à la fin de la phrase.
\item \all{Ich werde morgen einen Dialog organisieren.} est une phrase au futur I.
\item \all{die Toleranz} est un nom masculin.
\end{enumerate}
\end{exercice}

\begin{exercice}[Vocabulaire : trouve le mot]
Complète chaque phrase avec le mot qui convient : \all{der Frieden, der Krieg, die Versöhnung, der Dialog, die Toleranz, die Verständigung}.
\begin{enumerate}
\item Nach dem Konflikt brauchen die beiden Familien \dots\ (réconciliation).
\item Ohne \dots\ gibt es keine Lösung : man muss miteinander sprechen !
\item \dots\ bedeutet: Man akzeptiert die anderen, auch wenn sie anders denken.
\item Nach vielen Jahren \dots\ wünschen sich alle Menschen endlich Frieden.
\end{enumerate}
\end{exercice}

\begin{exercice}[Conjugaison à compléter]
Complète avec la forme correcte du verbe entre parenthèses.
\begin{enumerate}
\item Impératif (du) : \all{\dots\ (zuhören) deinem Nachbarn!}
\item Impératif (ihr) : \all{\dots\ (sein) ruhig und tolerant!}
\item Impératif (Sie) : \all{\dots\ (lesen) Sie den Text, bitte!}
\item Futur I : \all{Wir \dots\ (organisieren) nächste Woche eine Debatte.}
\item Futur I : \all{Ich \dots\ (kämpfen) für den Frieden.}
\end{enumerate}
\end{exercice}

\courssub{Je m'entraîne}

\begin{exercice}[Mets les verbes à l'impératif]
Mets les verbes entre parenthèses à l'impératif à la personne indiquée.
\begin{enumerate}
\item (du / mitkommen) \all{\dots\ zur Friedenswoche \dots!}
\item (ihr / helfen) \all{\dots\ euren Freunden!}
\item (Sie / sein) \all{\dots\ geduldig, bitte!}
\item (du / lesen) \all{\dots\ den Text über die Versöhnung!}
\item (ihr / zuhören) \all{\dots\ dem Mediator!}
\end{enumerate}
\end{exercice}

\begin{exercice}[Conjugue au futur I]
Conjugue les verbes entre parenthèses au futur I.
\begin{enumerate}
\item \all{Wir \dots\ (organisieren) ein Friedensfest in der Schule.}
\item \all{Du \dots\ (lernen) viel über Toleranz.}
\item \all{Sie (pluriel) \dots\ (verzeihen) ihren Nachbarn.}
\item \all{Ich \dots\ (kämpfen) für den Frieden in meiner Gemeinde.}
\item \all{Ihr \dots\ (besuchen) nächstes Jahr Deutschland.}
\end{enumerate}
\end{exercice}

\begin{exercice}[Transforme : présent + marqueur temporel ou futur I]
Transforme les phrases : si elles sont au présent avec un marqueur temporel, mets-les au futur I ; si elles sont au futur I, mets-les au présent avec un marqueur temporel adapté.
\begin{enumerate}
\item \all{Morgen fahren wir nach Bonn.}
\item \all{Nächste Woche beginnt der Dialog.}
\item \all{Die Schüler werden ein Friedensprojekt starten.}
\item \all{Ich werde morgen mit meinem Nachbarn sprechen.}
\end{enumerate}
Exemple : \all{Morgen fahren wir nach Bonn.} $\rightarrow$ \all{Morgen werden wir nach Bonn fahren.}
\end{exercice}

\begin{exercice}[Appariements : mots et définitions]
Relie chaque mot allemand à sa définition française.
\begin{center}
\begin{tabularx}{\linewidth}{|X|X|}
\hline
\rowcolor{popblueL} Deutsch & Français \\
\hline
1. der Frieden & a) parler ensemble pour résoudre un problème \\
\hline
2. der Krieg & b) accepter les différences des autres \\
\hline
3. der Dialog & c) l'état sans guerre ni violence \\
\hline
4. die Toleranz & d) le conflit armé entre des groupes ou des pays \\
\hline
5. die Versöhnung & e) retrouver de bonnes relations après un conflit \\
\hline
\end{tabularx}
\end{center}
\end{exercice}

\courssub{Je me prépare au Bac}

\begin{exercice}[Texte de compréhension type Bac]
Lis le texte suivant, puis réponds aux questions.

\begin{textebox}[Versöhnung nach dem Konflikt]
In einem Dorf im Westen Kameruns gab es lange einen Konflikt zwischen zwei Familien. Sie stritten über ein Stück Land und sprachen zehn Jahre nicht miteinander. Die Kinder der beiden Familien durften nicht zusammen spielen, und im Dorf gab es oft Spannungen.

Dann kam ein neuer Lehrer in die Schule des Dorfes. Er organisierte einen Dialog zwischen den Familien. „Hört einander zu!“, sagte er. „Sprecht über eure Probleme, aber bleibt ruhig und respektvoll!“ Zuerst war es schwer, aber langsam begannen die Menschen, einander zu verstehen.

Nach drei Monaten fanden die Familien eine Lösung: Sie werden das Land jetzt gemeinsam nutzen. Am Ende des Treffens sagte der älteste Mann des Dorfes: „Die Versöhnung ist wichtiger als das Land. Der Frieden beginnt bei uns, in unseren Herzen.“

Nächstes Jahr wird das Dorf ein großes Friedensfest organisieren. Alle werden zusammen essen, singen und tanzen. Der Lehrer ist stolz: „Diese Gemeinde wird ein Beispiel für andere Dörfer sein.“
\end{textebox}

\textbf{Questions de compréhension} (réponds en allemand, avec des phrases complètes) :
\begin{enumerate}
\item Worüber stritten die beiden Familien ? (question globale)
\item Wie lange sprachen die Familien nicht miteinander ?
\item Wer organisierte den Dialog ?
\item Nenne zwei Ratschläge (conseils) des Lehrers.
\item Richtig oder falsch ? Justifie mit une phrase du texte :
\begin{enumerate}
\item Die Kinder durften zusammen spielen.
\item Die Familien werden das Land gemeinsam nutzen.
\end{enumerate}
\item Was wird das Dorf nächstes Jahr machen ?
\end{enumerate}

\textbf{Questions de langue} :
\begin{enumerate}
\item Relève dans le texte deux impératifs et donne leur infinitif.
\item Relève deux phrases au futur I et conjugue les verbes au présent.
\end{enumerate}
\end{exercice}

\begin{exercice}[Thème et version]
\textbf{A. Thème.} Traduis en allemand les phrases suivantes :
\begin{enumerate}
\item Écoute ton voisin !
\item Soyez tolérants ! (forme \all{ihr})
\item Nous organiserons un débat.
\item Tu devrais parler avec lui. (conseil avec \all{sollen})
\item La paix commence par le dialogue.
\end{enumerate}

\textbf{B. Version.} Traduis en français le paragraphe suivant :

\begin{textebox}[Ein Mediationsprojekt]
An einer Schule in München gibt es ein neues Projekt: Schüler helfen anderen Schülern bei Konflikten. „Kommt zu uns, wenn ihr ein Problem habt!“, sagen die Mediatoren. „Wir werden euch zuhören und eine Lösung finden.“ Nächstes Jahr werden zwanzig neue Schüler diese Ausbildung machen. Die Direktorin sagt: „Dieses Projekt wird unsere Schule verändern.“
\end{textebox}
\end{exercice}

\begin{exercice}[Production écrite guidée type Bac]
\textbf{Sujet :} \all{Was wirst du nach dem Bac tun, um den Frieden in deiner Gemeinde zu fördern ?}

Rédige un texte de \textbf{8 à 10 phrases} (environ 80 à 100 mots) en allemand. Tu utiliseras obligatoirement :
\begin{itemize}
\item le futur I (au moins 4 verbes différents) ;
\item au moins un impératif ou un conseil avec \all{sollen} ;
\item au moins 5 mots du vocabulaire du thème : \all{der Frieden, der Dialog, die Toleranz, die Versöhnung, die Verständigung, verzeihen, zuhören, helfen}.
\end{itemize}

\textbf{Pistes :} Parle de tes projets concrets (débats, médiation, aide aux jeunes, sport, école), de ton quartier ou de ton village (Yaoundé, Douala, ta commune), et explique pourquoi la paix commence par chacun de nous.
\end{exercice}

\newpage

\coursec{Corrigés des exercices}

\begin{corrige}[Exercice 1 — QCM : vocabulaire de la paix]
\begin{enumerate}
\item \textbf{b) la réconciliation}. \all{die Versöhnung} vient du verbe \all{sich versöhnen} (se réconcilier). Famille de mots : \all{versöhnen} (réconcilier), \all{versöhnlich} (conciliant).
\item \textbf{b) die Kriege}. Le pluriel de \all{der Krieg} se forme avec la terminaison \all{-e} : \all{der Krieg, die Kriege}. Attention : \all{die Krieger} signifie « les guerriers » — faux ami de forme !
\item \textbf{a) l'entente, la compréhension mutuelle}. \all{die Verständigung} vient de \all{sich verständigen} (s'entendre, se mettre d'accord). Ne pas confondre avec \all{das Verständnis} (la compréhension, la compassion).
\item \textbf{b) verzeihen}. \all{verzeihen} (verbe fort : \all{verzeiht, verzieh, hat verziehen}) signifie « pardonner ». \all{kämpfen} = combattre ; \all{zerstören} = détruire.
\end{enumerate}
\textbf{Point de vigilance :} retiens toujours chaque nom avec son article et son pluriel : \all{der Frieden} (sans pluriel), \all{der Krieg, die Kriege}, \all{die Versöhnung, -en}, \all{der Dialog, -e}, \all{die Toleranz} (sans pluriel), \all{die Verständigung, -en}.
\end{corrige}

\begin{corrige}[Exercice 2 — Vrai ou faux ?]
\begin{enumerate}
\item \textbf{richtig}. \all{bei uns} signifie « chez nous / avec nous » ; la phrase complète se traduit : « La paix commence avec nous. »
\item \textbf{falsch}. À l'impératif avec \all{Sie}, le verbe se place en \textbf{première position}, comme dans une question : \all{Lesen Sie den Text!} (« Lisez le texte ! »). C'est la même structure qu'à l'impératif \all{du} et \all{ihr} : le verbe est toujours en tête de phrase.
\item \textbf{richtig}. La phrase est construite avec l'auxiliaire \all{werden} conjugué (\all{werde}) + l'infinitif \all{organisieren} placé à la fin : c'est bien le futur I. Traduction : « J'organiserai demain un dialogue. »
\item \textbf{falsch}. \all{die Toleranz} est un nom \textbf{féminin} : presque tous les noms en \all{-anz} (comme en \all{-enz}, \all{-ion}, \all{-heit}, \all{-keit}) sont féminins. On dit donc \all{die Toleranz}, \all{die Konferenz}, \all{die Diskussion}.
\end{enumerate}
\end{corrige}

\begin{corrige}[Exercice 3 — Vocabulaire : trouve le mot]
\begin{enumerate}
\item \all{Nach dem Konflikt brauchen die beiden Familien \textbf{die Versöhnung}.} — « Après le conflit, les deux familles ont besoin de la réconciliation. » (accusatif féminin : \all{die} ne change pas.)
\item \all{Ohne \textbf{Dialog} gibt es keine Lösung : man muss miteinander sprechen !} — « Sans dialogue, il n'y a pas de solution : il faut se parler ! » (après la préposition \all{ohne} + accusatif, on peut omettre l'article ; \all{ohne einen Dialog} est aussi correct.)
\item \all{\textbf{Toleranz} bedeutet: Man akzeptiert die anderen, auch wenn sie anders denken.} — « La tolérance signifie : on accepte les autres, même s'ils pensent différemment. »
\item \all{Nach vielen Jahren \textbf{Krieg} wünschen sich alle Menschen endlich Frieden.} — « Après de nombreuses années de guerre, tous les gens souhaitent enfin la paix. » (construction sans article après \all{Jahren} ; on peut aussi écrire \all{nach vielen Jahren des Krieges}, génitif.)
\end{enumerate}
\textbf{Point de vigilance :} après \all{ohne} et \all{nach}, attention au cas : \all{ohne den Dialog} (accusatif), \all{nach dem Krieg} (datif). Ici l'absence d'article est acceptée car on parle de notions abstraites générales.
\end{corrige}

\begin{corrige}[Exercice 4 — Conjugaison à compléter]
\begin{enumerate}
\item \all{\textbf{Hör} deinem Nachbarn \textbf{zu}!} — « Écoute ton voisin ! » Verbe à particule séparable : à l'impératif \all{du}, la particule \all{zu} va à la fin. Attention au datif : \all{zuhören} se construit avec le datif (\all{deinem Nachbarn}).
\item \all{\textbf{Seid} ruhig und tolerant!} — « Soyez calmes et tolérants ! » Impératif \all{ihr} irrégulier de \all{sein} : \all{seid} (et non \emph{seit}, qui est la préposition « depuis »).
\item \all{\textbf{Lesen} Sie den Text, bitte!} — « Lisez le texte, s'il vous plaît ! » À l'impératif \all{Sie}, on reprend la forme du présent (\all{sie lesen}) suivie du pronom \all{Sie}.
\item \all{Wir \textbf{werden organisieren} nächste Woche eine Debatte.} — « Nous organiserons un débat la semaine prochaine. » Futur I : \all{werden} conjugué en 2\up{e} position + infinitif à la fin.
\item \all{Ich \textbf{werde kämpfen} für den Frieden.} — « Je combattrai pour la paix. »
\end{enumerate}
\end{corrige}

\begin{corrige}[Exercice 5 — Mets les verbes à l'impératif]
\begin{enumerate}
\item \all{\textbf{Komm} zur Friedenswoche \textbf{mit}!} — « Viens à la semaine de la paix ! » Verbe séparable : la particule \all{mit} se détache et va en fin de phrase.
\item \all{\textbf{Helft} euren Freunden!} — « Aidez vos amis ! » Impératif \all{ihr} = radical de la 2\up{e} personne du pluriel du présent sans \all{-t} final ajouté : \all{ihr helft} $\rightarrow$ \all{Helft!}. Attention : \all{helfen} + datif (\all{euren Freunden}).
\item \all{\textbf{Seien} Sie geduldig, bitte!} — « Soyez patient, s'il vous plaît ! » Forme irrégulière : \all{sein} a pour impératif \all{Sie} la forme \all{seien Sie} (et non \emph{sind Sie}).
\item \all{\textbf{Lies} den Text über die Versöhnung!} — « Lis le texte sur la réconciliation ! » Verbe fort à vocalisme changeant : \all{lesen} fait \all{du liest} au présent, donc impératif \all{Lies!} (sans \all{-e} et avec le changement e $\rightarrow$ ie).
\item \all{\textbf{Hört} dem Mediator \textbf{zu}!} — « Écoutez le médiateur ! » Particule \all{zu} en finale ; datif après \all{zuhören} : \all{dem Mediator}.
\end{enumerate}
\textbf{Point de vigilance :} les trois formes à retenir : \all{Komm!} (du), \all{Kommt!} (ihr), \all{Kommen Sie!} (Sie). Seules exceptions importantes : \all{sei / seid / seien Sie} pour \all{sein}, et les verbes forts à changement de voyelle e $\rightarrow$ i/ie qui gardent ce changement à l'impératif \all{du} : \all{Gib!}, \all{Nimm!}, \all{Sprich!}, \all{Lies!}.
\end{corrige}

\begin{corrige}[Exercice 6 — Conjugue au futur I]
\begin{enumerate}
\item \all{Wir \textbf{werden organisieren} ein Friedensfest in der Schule.} — « Nous organiserons une fête de la paix à l'école. »
\item \all{Du \textbf{wirst lernen} viel über Toleranz.} — « Tu apprendras beaucoup de choses sur la tolérance. »
\item \all{Sie \textbf{werden verzeihen} ihren Nachbarn.} — « Elles pardonneront à leurs voisins. » (\all{verzeihen} + datif : \all{ihren Nachbarn}.)
\item \all{Ich \textbf{werde kämpfen} für den Frieden in meiner Gemeinde.} — « Je combattrai pour la paix dans ma commune. »
\item \all{Ihr \textbf{werdet besuchen} nächstes Jahr Deutschland.} — « Vous visiterez l'Allemagne l'année prochaine. »
\end{enumerate}
\textbf{Rappel de la conjugaison de \all{werden} au présent :} \all{ich werde, du wirst, er/sie/es wird, wir werden, ihr werdet, sie/Sie werden}. L'infinitif se place toujours en dernière position : \all{Ich werde morgen nach Douala fahren.}
\end{corrige}

\begin{corrige}[Exercice 7 — Transforme : présent + marqueur temporel ou futur I]
\begin{enumerate}
\item \all{Morgen fahren wir nach Bonn.} $\rightarrow$ \all{\textbf{Morgen werden wir nach Bonn fahren.}} — « Demain nous irons à Bonn. » (présent + marqueur $\rightarrow$ futur I ; l'infinitif \all{fahren} passe en fin de phrase.)
\item \all{Nächste Woche beginnt der Dialog.} $\rightarrow$ \all{\textbf{Nächste Woche wird der Dialog beginnen.}} — « La semaine prochaine, le dialogue commencera. »
\item \all{Die Schüler werden ein Friedensprojekt starten.} $\rightarrow$ \all{\textbf{Nächsten Monat starten die Schüler ein Friedensprojekt.}} — « Le mois prochain, les élèves lancent un projet de paix. » (futur I $\rightarrow$ présent + marqueur temporel ; au présent, le verbe conjugué revient en 2\up{e} position.)
\item \all{Ich werde morgen mit meinem Nachbarn sprechen.} $\rightarrow$ \all{\textbf{Morgen spreche ich mit meinem Nachbarn.}} — « Demain, je parle avec mon voisin. »
\end{enumerate}
\textbf{Point de vigilance :} en allemand, le présent + un marqueur temporel (\all{morgen}, \all{nächste Woche}, \all{bald}) exprime très souvent le futur, comme en français (« demain, je pars »). Le futur I insiste davantage sur la promesse, la prédiction ou l'intention. Dans les deux cas, respecte la place du verbe : 2\up{e} position au présent, infinitif final au futur I.
\end{corrige}

\begin{corrige}[Exercice 8 — Appariements : mots et définitions]
\begin{center}
\begin{tabularx}{\linewidth}{|X|X|}
\hline
\rowcolor{popblueL} Deutsch & Français \\
\hline
1. der Frieden & c) l'état sans guerre ni violence \\
\hline
2. der Krieg & d) le conflit armé entre des groupes ou des pays \\
\hline
3. der Dialog & a) parler ensemble pour résoudre un problème \\
\hline
4. die Toleranz & b) accepter les différences des autres \\
\hline
5. die Versöhnung & e) retrouver de bonnes relations après un conflit \\
\hline
\end{tabularx}
\end{center}
\textbf{Correspondances :} 1-c, 2-d, 3-a, 4-b, 5-e.

\textbf{Pour mémoriser :} \all{der Frieden} s'oppose à \all{der Krieg} ; \all{der Dialog} et \all{die Verständigung} sont les moyens de la paix ; \all{die Toleranz} et \all{die Versöhnung} en sont les fruits. Étymologie utile : \all{Dialog} vient du grec \emph{dia} (à travers) + \emph{logos} (parole) : la parole qui passe de l'un à l'autre.
\end{corrige}

\begin{corrige}[Exercice 9 — Texte de compréhension type Bac]
\textbf{Barème indicatif (sur 20) :} compréhension globale et détaillée : 10 pts ; questions de langue : 4 pts ; qualité de l'expression allemande (phrases complètes, ordre des mots, orthographe) : 6 pts.

\textbf{Questions de compréhension — réponses modèles :}
\begin{enumerate}
\item \all{Die beiden Familien stritten über ein Stück Land.} — « Les deux familles se disputaient à propos d'un morceau de terrain. » (question globale : on reformule sans recopier tout le paragraphe.)
\item \all{Sie sprachen zehn Jahre nicht miteinander.} — « Elles ne se sont pas parlé pendant dix ans. »
\item \all{Ein neuer Lehrer organisierte den Dialog zwischen den Familien.} — « Un nouvel enseignant a organisé le dialogue entre les familles. »
\item Deux conseils du professeur (impératifs du texte) : \all{„Hört einander zu!“} (« Écoutez-vous les uns les autres ! ») et \all{„Sprecht über eure Probleme, aber bleibt ruhig und respektvoll!“} (« Parlez de vos problèmes, mais restez calmes et respectueux ! »).
\item \begin{enumerate}
\item \textbf{falsch}. Justification : \all{„Die Kinder der beiden Familien durften nicht zusammen spielen.“} — « Les enfants des deux familles n'avaient pas le droit de jouer ensemble. »
\item \textbf{richtig}. Justification : \all{„Sie werden das Land jetzt gemeinsam nutzen.“} — « Ils vont maintenant utiliser le terrain en commun. »
\end{enumerate}
\item \all{Nächstes Jahr wird das Dorf ein großes Friedensfest organisieren. Alle werden zusammen essen, singen und tanzen.} — « L'année prochaine, le village organisera une grande fête de la paix. Tous mangeront, chanteront et danseront ensemble. »
\end{enumerate}

\textbf{Questions de langue :}
\begin{enumerate}
\item Deux impératifs relevés dans le texte : \all{Hört (einander zu)!} $\rightarrow$ infinitif \all{zuhören} ; \all{Sprecht!} $\rightarrow$ infinitif \all{sprechen} ; \all{bleibt!} $\rightarrow$ infinitif \all{bleiben}. (Trois possibles ; deux suffisent.) Ce sont des impératifs à la 2\up{e} personne du pluriel (\all{ihr}) : le professeur s'adresse aux deux familles.
\item Deux phrases au futur I :
\begin{itemize}
\item \all{„Sie werden das Land jetzt gemeinsam nutzen.“} $\rightarrow$ présent : \all{Sie nutzen das Land jetzt gemeinsam.}
\item \all{„Nächstes Jahr wird das Dorf ein großes Friedensfest organisieren.“} $\rightarrow$ présent : \all{Nächstes Jahr organisiert das Dorf ein großes Friedensfest.}
\item (Autre possibilité : \all{„Diese Gemeinde wird ein Beispiel für andere Dörfer sein.“} $\rightarrow$ \all{Diese Gemeinde ist ein Beispiel für andere Dörfer.})
\end{itemize}
\end{enumerate}

\textbf{Points de vigilance :}
\begin{itemize}
\item Réponds toujours par une phrase complète avec le verbe conjugué en 2\up{e} position.
\item Pour « richtig oder falsch », la justification par une citation exacte du texte est obligatoire pour avoir tous les points.
\item Ne confonds pas \all{streiten über + Akkusativ} (se disputer à propos de) : \all{über ein Stück Land}.
\end{itemize}
\end{corrige}

\begin{corrige}[Exercice 10 — Thème et version]
\textbf{A. Thème — traductions modèles :}
\begin{enumerate}
\item \all{\textbf{Hör deinem Nachbarn zu!}} (ou \all{Höre deinem Nachbarn zu!}) — Impératif \all{du} du verbe séparable \all{zuhören} ; particule \all{zu} en finale ; \textbf{datif} obligatoire : \all{deinem Nachbarn}. Erreur fréquente : \emph{Hör deinen Nachbarn} (accusatif) est faux avec \all{zuhören}.
\item \all{\textbf{Seid tolerant!}} — Impératif \all{ihr} irrégulier de \all{sein} : \all{seid}.
\item \all{\textbf{Wir werden eine Debatte organisieren.}} — Futur I : \all{werden} en 2\up{e} position, infinitif \all{organisieren} à la fin. \all{die Debatte, -n} : nom féminin.
\item \all{\textbf{Du sollst mit ihm sprechen.}} — Conseil avec \all{sollen} : modal conjugué en 2\up{e} position (\all{sollst}) + infinitif \all{sprechen} en fin de phrase. « avec lui » = \all{mit ihm} (datif après \all{mit}).
\item \all{\textbf{Der Frieden beginnt mit dem Dialog.}} — « commencer par » se rend ici par \all{beginnen mit + Dativ} : \all{mit dem Dialog}. (On accepte aussi \all{Der Frieden beginnt beim Dialog.})
\end{enumerate}

\textbf{B. Version — traduction modèle :}

« Dans une école de Munich, il existe un nouveau projet : des élèves aident d'autres élèves en cas de conflits. „Venez nous voir quand vous avez un problème !“, disent les médiateurs. „Nous vous écouterons et nous trouverons une solution.“ L'année prochaine, vingt nouveaux élèves suivront cette formation. La directrice dit : „Ce projet va changer notre école.“ »

\textbf{Points de vigilance pour la version :}
\begin{itemize}
\item \all{Kommt zu uns, wenn ihr ein Problem habt!} : subordonnée avec \all{wenn}, verbe \all{habt} à la fin — traduire par « quand vous avez un problème ».
\item \all{diese Ausbildung machen} : « suivre cette formation » (et non « faire cette éducation » — faux ami de sens).
\item \all{wird unsere Schule verändern} : futur I, « changera / va changer notre école ».
\end{itemize}
\textbf{Barème indicatif (sur 10) :} thème 5 pts (1 pt par phrase : forme verbale 0,5 + cas et ordre des mots 0,5) ; version 5 pts (fidélité 3 pts + qualité du français 2 pts).
\end{corrige}

\begin{corrige}[Exercice 11 — Production écrite guidée type Bac]
\textbf{Barème indicatif (sur 20) :} respect du sujet et de la longueur (8 à 10 phrases) : 4 pts ; futur I correct (au moins 4 verbes différents) : 4 pts ; impératif ou conseil avec \all{sollen} : 2 pts ; vocabulaire du thème (au moins 5 mots) : 4 pts ; grammaire générale (ordre des mots, cas, orthographe, majuscules) : 4 pts ; cohérence et lien logique : 2 pts.

\textbf{Proposition de rédaction modèle :}

\begin{textebox}[Mein Beitrag zum Frieden]
Nach dem Bac werde ich mich für den Frieden in meiner Gemeinde in Yaoundé engagieren. Zuerst werde ich mit meinen Freunden einen Jugendclub gründen. Dort werden wir jeden Monat einen Dialog über Toleranz und Verständigung organisieren. Ich werde auch den jungen Schülern in meinem Quartier helfen, denn Bildung ist wichtig für den Frieden. Außerdem werde ich bei Konflikten zwischen Nachbarn vermitteln: Ich werde beiden Seiten zuhören und eine faire Lösung suchen. Mein Motto ist: „Verzeiht einander und sprecht miteinander!“ Ich glaube: Jeder soll in seiner Familie und in seiner Straße ein Beispiel sein. Die Versöhnung beginnt mit kleinen Schritten. Der Frieden beginnt bei uns — und er beginnt heute.
\end{textebox}

\textbf{Traduction du modèle :} « Après le bac, je m'engagerai pour la paix dans ma commune à Yaoundé. D'abord, je fonderai un club de jeunes avec mes amis. Là, nous organiserons chaque mois un dialogue sur la tolérance et l'entente. J'aiderai aussi les jeunes élèves de mon quartier, car l'éducation est importante pour la paix. De plus, je servirai de médiateur dans les conflits entre voisins : j'écouterai les deux parties et je chercherai une solution juste. Ma devise est : „Pardonnez-vous les uns aux autres et parlez-vous !“ Je crois que chacun doit être un exemple dans sa famille et dans sa rue. La réconciliation commence par de petits pas. La paix commence avec nous — et elle commence aujourd'hui. »

\textbf{Vérification des contraintes dans le modèle :}
\begin{itemize}
\item Futur I avec 6 verbes différents : \all{werde engagieren, werde gründen, werden organisieren, werde helfen, werde vermitteln, werde zuhören}.
\item Impératif : \all{„Verzeiht einander und sprecht miteinander!“} ; conseil avec \all{sollen} : \all{Jeder soll ein Beispiel sein}.
\item Vocabulaire du thème : \all{der Frieden, der Dialog, die Toleranz, die Verständigung, die Versöhnung, verzeihen, zuhören, helfen} — 8 mots, soit plus que les 5 exigés.
\end{itemize}

\textbf{Points de vigilance :}
\begin{itemize}
\item Au futur I, l'infinitif va toujours en fin de phrase : \all{Ich werde einen Club gründen}, jamais \emph{Ich werde gründen einen Club}.
\item Après \all{sollen}, même règle : modal en 2\up{e} position, infinitif final.
\item Majuscule à tous les noms : \all{der Frieden, der Dialog, die Toleranz, die Gemeinde, das Quartier}.
\item \all{sich engagieren für + Akkusativ} : \all{für den Frieden}.
\item Évite les phrases trop longues : 8 à 10 phrases courtes et correctes valent mieux qu'un texte ambitieux plein de fautes.
\end{itemize}
\end{corrige}

\newpage

\coursec{Fiche bilan}

\begin{popfigure}
\begin{tikzpicture}[
    mindmap,
    grow cyclic,
    every node/.style={concept, execute at begin node=\hskip0pt},
    concept color=popblue,
    level 1/.style={level distance=4.5cm, sibling angle=360/6},
    level 2/.style={level distance=3cm, sibling angle=360/6},
    texte/.style={concept, font=\small\bfseries, text width=2.8cm, align=center, inner sep=4pt},
    cat1/.style={concept color=popblue},
    cat2/.style={concept color=popgreen},
    cat3/.style={concept color=poporange},
    cat4/.style={concept color=poppurple},
    cat5/.style={concept color=poppink},
    cat6/.style={concept color=popgold}
]
\node[texte, font=\large\bfseries, text width=3.5cm] {ALA26 : Culture de la paix}
    child[cat1] { node[texte, align=center]{Texte support\\„Frieden beginnt bei uns“} }
    child[cat2] { node[texte, align=center]{Lexique\\Frieden, Krieg, Versöhnung,\\Toleranz, Dialog} }
    child[cat3] { node[texte, align=center]{Grammaire 1\\Impératif \& Conseils} }
    child[cat4] { node[texte, align=center]{Grammaire 2\\Futur (werden + Inf.)} }
    child[cat5] { node[texte, align=center]{Conjugaison\\Präsens, Futur I,\\Imperativformen} }
    child[cat6] { node[texte, align=center]{Expression\\Résumé, Questions,\\Production guidée} };
\end{tikzpicture}
\legende{Carte mentale de la leçon ALA26 — Première A}
\end{popfigure}

\begin{bilan}

\end{bilan}

\courssub{Lexique clé : Frieden und Konflikt}

\begin{center}
\begin{tabularx}{\linewidth}{|>{\bfseries}l|X|X|}\hline
\rowcolor{popblueL}
\textbf{Deutsch (Artikel, Plural)} & \textbf{Français} & \textbf{Famille / Exemple} \\ \hline
der Frieden, \textendash & la paix & friedlich (adj.), friedenstiftend \\
der Krieg, die Kriege & la guerre & kriegerisch, der Kriegsgegner \\
die Versöhnung, \textendash & la réconciliation & sich versöhnen mit (Dat.) \\
die Toleranz, \textendash & la tolérance & tolerant, tolerieren \\
der Dialog, die Dialoge & le dialogue & dialogisch, der Gesprächspartner \\
die Verständigung, \textendash & l'entente, la compréhension & sich verständigen über (Akk.) \\
der Konflikt, die Konflikte & le conflit & konfliktär, lösen \\
die Gewalt, \textendash & la violence & gewaltfrei, Gewalt anwenden \\
das Miteinander, \textendash & le vivre-ensemble & miteinander auskommen \\
die Menschenrechte (Pl.) & les droits de l'homme & die Menschenwürde \\ \hline
\end{tabularx}
\end{center}

\courssub{Règles de grammaire essentielles}

\begin{center}
\begin{tabularx}{\linewidth}{|>{\bfseries}p{3.2cm}|X|X|}\hline
\rowcolor{popblueL}
\textbf{Notion} & \textbf{Règle / Formation} & \textbf{Exemples clés} \\ \hline
\textbf{Impératif (Imperativ)} & \textbf{2e pers. sg (du) :} radical (+ e) \textit{ohne} pronom.\\
\textbf{2e pers. pl (ihr) :} verbe fini \textit{ohne} pronom.\\
\textbf{Forme polie (Sie) :} verbe à l'infinitif + \textit{Sie}. & \all{Mach das Fenster auf!}\\
\all{Seien Sie pünktlich!}\\
\all{Hört gut zu!} \\ \hline
\textbf{Conseil / Exhortation} & \textit{sollen} + infinitif (conseil moral).\\
\textit{könnten} + infinitif (Konj. II, conseil poli).\\
\textit{Lasst uns} + infinitif (proposition). & \all{Du sollst nicht lügen.}\\
\all{Wir könnten zusammen gehen.}\\
\all{Lasst uns Frieden stiften!} \\ \hline
\textbf{Futur I (Futur)} & \textbf{werden} (conjugué) + \textbf{infinitif} à la fin.\\
Emploi : action future, supposition, promesse. & \all{Ich werde lernen.}\\
\all{Es wird regnen.}\\
\all{Wir werden uns versöhnen.} \\ \hline
\textbf{Ordre des mots (Futur)} & Position 2 : \textit{werden} conjugué.\\
Position finale : \textit{infinitif principal}. & \all{Morgen \textbf{werde} ich \textbf{reisen}.} \\ \hline
\end{tabularx}
\end{center}

\courssub{Méthodes express}

\begin{itemize}
    \item \textbf{Compréhension :} 1. Lire le titre + images $\rightarrow$ hypothèses. 2. Lecture globale (idée générale). 3. Lecture détaillée (mots-clés, connecteurs \all{aber, denn, weil}). 4. Répondre en allemand \textit{mit eigenen Worten}.
    \item \textbf{Traduction (Thème/Version) :} Identifier le verbe $\rightarrow$ Cas (Sujet/Objet) $\rightarrow$ Temps (Présent/Futur) $\rightarrow$ Ordre (V2 / Verbe final) $\rightarrow$ Vocabulaire précis (faux amis : \all{konsequent} $\neq$ conséquent).
    \item \textbf{Production (Résumé/Paragraphe) :} Structure : \textit{Einleitung} (Sujet + \all{handelt von}) $\rightarrow$ \textit{Hauptteil} (3 arguments avec \all{zuerst, außerdem, schließlich}) $\rightarrow$ \textit{Schluss} (Meinung/Appel avec \all{wir müssen / wir werden}).
\end{itemize}



\begin{aretenir}[Checklist avant le Bac \textendash{} ALA26]
\begin{itemize}
    \item $\square$ Je connais le genre et le pluriel des 10 mots de vocabulaire clé (\all{der Frieden, die Versöhnung, die Menschenrechte}...).
    \item $\square$ Je sais former l'impératif aux 3 formes (\textit{du, ihr, Sie}) pour verbes réguliers, à voyelle modifiée, en \textit{-eln/-ern} et \textit{sein}.
    \item $\square$ Je distingue \textbf{conseil} (\all{sollen}) et \textbf{impératif} (ordre direct) et j'utilise \all{könnten} pour la politesse.
    \item $\square$ Je conjugue \textbf{werden} au présent sans erreur (\all{ich werde, du wirst, er wird, wir werden, ihr werdet, sie/Sie werden}).
    \item $\square$ Je construis correctement le \textbf{Futur I} : \all{werden} en position 2, infinitif à la fin.
    \item $\square$ Je place l'adverbe de temps (\all{morgen, nächstes Jahr}) en début de phrase $\rightarrow$ inversion sujet/verbe (\all{Morgen werde ich...}).
    \item $\square$ Je sais exprimer la réconciliation : \all{sich versöhnen mit (Dativ)} et \all{Verständigung erreichen}.
    \item $\square$ Je sais rédiger un résumé structuré (5-6 phrases) avec connecteurs logiques (\all{zuerst, dann, weil, deshalb}).
    \item $\square$ Je réponds aux questions de compréhension \textbf{en allemand}, phrases complètes, pas de copier-coller brut.
    \item $\square$ Je repère les verbes à particule séparable dans le texte (\all{aufgeben, anfangen, zuhören}) et je les conjugue correctement.
    \item $\square$ Je connais les faux amis du thème : \all{konsequent} (conséquent/ferme), \all{sensibel} (sensible), \all{aktuell} (actuel).
\end{itemize}
\end{aretenir}

\findecours

\end{document}
